%=====================================================================
%  CE-1115 Seguridad de la Informacion -- Proyecto, Parte 1
%  Diapositivas de la defensa oral
%  Cliente: Verdeazul Carbono S.A. -- Grupo 6
%
%  SISTEMA VISUAL (no hay plantilla corporativa):
%    Lamina   : 13,33 x 7,5 in (las medidas de PowerPoint 16:9).
%    Tipografia: Lato, dos pesos (regular y bold).
%    Escala   : cuatro tamanos -- 40 / 28 / 18 / 13 pt. Nada menor de 13.
%    Color    : UN acento (#0F6E6A) + grises. El acento marca el dato
%               importante de la lamina; nunca decora.
%    Reticula : margenes de 0,6 in en los cuatro lados; todo a la
%               izquierda.
%    Cada titulo es una afirmacion. El detalle va en las notas (\note).
%
%  Compilar:              pdflatex P1_Presentacion_grupo6.tex   (2 veces)
%  Version con notas:     pdflatex -jobname=P1_Presentacion_grupo6_notas
%                           "\def\connotas{1}%=====================================================================
%  CE-1115 Seguridad de la Informacion -- Proyecto, Parte 1
%  Diapositivas de la defensa oral
%  Cliente: Verdeazul Carbono S.A. -- Grupo 6
%
%  SISTEMA VISUAL (no hay plantilla corporativa):
%    Lamina   : 13,33 x 7,5 in (las medidas de PowerPoint 16:9).
%    Tipografia: Lato, dos pesos (regular y bold).
%    Escala   : cuatro tamanos -- 40 / 28 / 18 / 13 pt. Nada menor de 13.
%    Color    : UN acento (#0F6E6A) + grises. El acento marca el dato
%               importante de la lamina; nunca decora.
%    Reticula : margenes de 0,6 in en los cuatro lados; todo a la
%               izquierda.
%    Cada titulo es una afirmacion. El detalle va en las notas (\note).
%
%  Compilar:              pdflatex P1_Presentacion_grupo6.tex   (2 veces)
%  Version con notas:     pdflatex -jobname=P1_Presentacion_grupo6_notas
%                           "\def\connotas{1}%=====================================================================
%  CE-1115 Seguridad de la Informacion -- Proyecto, Parte 1
%  Diapositivas de la defensa oral
%  Cliente: Verdeazul Carbono S.A. -- Grupo 6
%
%  SISTEMA VISUAL (no hay plantilla corporativa):
%    Lamina   : 13,33 x 7,5 in (las medidas de PowerPoint 16:9).
%    Tipografia: Lato, dos pesos (regular y bold).
%    Escala   : cuatro tamanos -- 40 / 28 / 18 / 13 pt. Nada menor de 13.
%    Color    : UN acento (#0F6E6A) + grises. El acento marca el dato
%               importante de la lamina; nunca decora.
%    Reticula : margenes de 0,6 in en los cuatro lados; todo a la
%               izquierda.
%    Cada titulo es una afirmacion. El detalle va en las notas (\note).
%
%  Compilar:              pdflatex P1_Presentacion_grupo6.tex   (2 veces)
%  Version con notas:     pdflatex -jobname=P1_Presentacion_grupo6_notas
%                           "\def\connotas{1}%=====================================================================
%  CE-1115 Seguridad de la Informacion -- Proyecto, Parte 1
%  Diapositivas de la defensa oral
%  Cliente: Verdeazul Carbono S.A. -- Grupo 6
%
%  SISTEMA VISUAL (no hay plantilla corporativa):
%    Lamina   : 13,33 x 7,5 in (las medidas de PowerPoint 16:9).
%    Tipografia: Lato, dos pesos (regular y bold).
%    Escala   : cuatro tamanos -- 40 / 28 / 18 / 13 pt. Nada menor de 13.
%    Color    : UN acento (#0F6E6A) + grises. El acento marca el dato
%               importante de la lamina; nunca decora.
%    Reticula : margenes de 0,6 in en los cuatro lados; todo a la
%               izquierda.
%    Cada titulo es una afirmacion. El detalle va en las notas (\note).
%
%  Compilar:              pdflatex P1_Presentacion_grupo6.tex   (2 veces)
%  Version con notas:     pdflatex -jobname=P1_Presentacion_grupo6_notas
%                           "\def\connotas{1}\input{P1_Presentacion_grupo6.tex}"
%=====================================================================
\documentclass[aspectratio=169,t]{beamer}

\usepackage[utf8]{inputenc}
\usepackage[T1]{fontenc}
\usepackage[spanish,es-noshorthands,es-tabla]{babel}
\usepackage[default]{lato}
\usepackage{booktabs,array,colortbl}
\usepackage{tikz}
\usetikzlibrary{positioning,arrows.meta,calc}

\ifdefined\connotas\setbeameroption{show notes}\fi

%--------------------------- lamina y margenes -----------------------
\geometry{papersize={13.333in,7.5in}}
\setbeamersize{text margin left=0.6in,text margin right=0.6in}

%--------------------------- color -----------------------------------
\definecolor{acento}{HTML}{0F6E6A}
\definecolor{acentoclaro}{HTML}{E7F2F1}
\definecolor{tinta}{HTML}{1A1A1A}
\definecolor{gris}{HTML}{5F5E5A}
\definecolor{regla}{HTML}{B9B8B0}
\definecolor{plano}{HTML}{FAFAFA}
\definecolor{negado}{HTML}{F2F1EC}

%--------------------------- escala: cuatro tamanos ------------------
\newcommand{\tUno}{\fontsize{40}{45}\selectfont}
\newcommand{\tDos}{\fontsize{28}{35}\selectfont}
\newcommand{\tTres}{\fontsize{18}{25}\selectfont}
\newcommand{\tCuatro}{\fontsize{13}{17.5}\selectfont}
\makeatletter
\renewcommand\normalsize{\@setfontsize\normalsize{18}{25}}
\renewcommand\small{\@setfontsize\small{13}{17.5}}
\renewcommand\footnotesize{\@setfontsize\footnotesize{13}{17.5}}
\renewcommand\scriptsize{\@setfontsize\scriptsize{13}{17.5}}
\renewcommand\tiny{\@setfontsize\tiny{13}{17.5}}
\renewcommand\large{\@setfontsize\large{28}{35}}
\renewcommand\Large{\@setfontsize\Large{40}{45}}
\makeatother

%--------------------------- cromo -----------------------------------
\usetheme{default}
\setbeamertemplate{navigation symbols}{}
\setbeamercolor{structure}{fg=acento}
\setbeamercolor{normal text}{fg=tinta}
\setbeamercolor{background canvas}{bg=white}
\setbeamerfont{note page}{size=\fontsize{20}{27}\selectfont}

\setbeamertemplate{frametitle}{%
  \vskip0.0in%
  \begin{minipage}[t]{\textwidth}
    \raggedright\color{tinta}\tUno\bfseries\insertframetitle\par
  \end{minipage}\par
  \vskip0.12in}

\gdef\presentador{}
\newcommand{\quien}[1]{\gdef\presentador{#1}}
\setbeamertemplate{footline}{%
  \hbox to\paperwidth{\hskip0.6in{\tCuatro\color{gris}\insertframenumber}\hfill
    {\tCuatro\color{gris}\presentador}\hskip0.6in}%
  \vskip0.45in}

%--------------------------- atajos ----------------------------------
\newcommand{\dato}[1]{\textcolor{acento}{\textbf{#1}}}
\newcommand{\pie}[1]{\leavevmode{\tCuatro\color{gris}#1}}
\newcommand{\cabt}[1]{\leavevmode{\tCuatro\bfseries\color{gris}#1}}
\newcolumntype{L}[1]{>{\raggedright\arraybackslash}p{#1}}
\arrayrulecolor{regla}
\setlength{\arrayrulewidth}{0.8pt}
\renewcommand{\arraystretch}{1.5}
% celdas de la matriz
\newcommand{\vb}[2]{\mbox{#1\,{\color{acento}\bfseries #2}}}
\newcommand{\sinacc}{\cellcolor{negado}\leavevmode{\color{gris}---}}
\newcommand{\suj}[1]{\leavevmode{\tTres\bfseries #1}}

%=====================================================================
\begin{document}

%--------------------------- 1. PORTADA ------------------------------
\begin{frame}[plain]
\vspace*{1.25in}
{\tCuatro\color{gris}PROYECTO \ $\cdot$ \ PARTE 1 \ $\cdot$ \ DEFENSA ORAL}\\[0.3in]
{\tUno\bfseries\color{tinta}Diseño del control de acceso\\[4pt] para Verdeazul Carbono}\\[0.25in]
{\color{acento}\rule{1.6in}{4pt}}\\[0.25in]
{\tDos\color{gris}Políticas, matriz y arquitectura de seguridad}\\[0.7in]
{\tTres\bfseries Randall Bolaños \quad Javier Tenorio \quad Julio Varela \quad Axel Flores}\\[0.12in]
{\tCuatro\color{gris}Grupo 6 \ $\cdot$ \ CE-1115 Seguridad de la Información \ $\cdot$ \ Prof. MSc. Andrés
Vargas Rivera \ $\cdot$ \ ITCR \ $\cdot$ \ 15 de octubre de 2026}
\note{Reparto propuesto. Randall: cliente, activos y separación de responsabilidades. Javier:
amenazas, marco legal, fricción, riesgos residuales y respuestas del cliente. Julio: modelo de control
de acceso y matriz. Axel: controles de historial y arquitectura. Todos deben poder responder por cualquier lámina:
las preguntas del profesor son individuales.}
\end{frame}

%--------------------------- 2. EL NEGOCIO ---------------------------
\begin{frame}{Verdeazul vende confianza:\\ cada certificado respalda una tonelada real}
\quien{Randall}
\vspace{0.15in}
\begin{tikzpicture}[x=1in,y=1in,
  caja/.style={draw=regla,line width=0.8pt,fill=plano,minimum width=2.0in,minimum height=0.95in,
               align=center,font=\tTres,inner sep=2pt},
  fl/.style={-{Stealth[length=9pt,width=8pt]},draw=gris,line width=1.2pt}]
\node[caja] (a) at (1.00,0) {Finca\\forestal};
\node[caja] (b) at (3.53,0) {Técnico\\mide};
\node[caja] (c) at (6.06,0) {Verificador\\aprueba};
\node[caja,draw=acento,line width=1.6pt] (d) at (8.59,0) {\textbf{Certificado}\\de carbono};
\node[caja] (e) at (11.12,0) {Mercado\\internacional};
\foreach \p/\q in {a/b,b/c,c/d,d/e} \draw[fl] (\p) -- (\q);
\end{tikzpicture}

\vspace{0.4in}
{\tDos \dato{Un dato inflado es fraude.}}

\vspace{0.12in}
Auditores europeos pueden pedir la trazabilidad de cualquier certificado de los últimos diez años.
\note{Qué decir: Verdeazul certifica captura de carbono en fincas forestales. El técnico mide, el
verificador aprueba y se emite un certificado que el dueño vende afuera. El valor del certificado es la
confianza en el dato. Hoy todo vive en hojas de cálculo compartidas.

Si preguntan por qué empezamos por el negocio: porque el enunciado pide políticas que describan a este
cliente, y casi todas las decisiones salen de esta lámina.}
\end{frame}

%--------------------------- 3. LOS DOS SUSTOS -----------------------
\begin{frame}{Los dos sustos de la carta son fallas de integridad}
\quien{Randall}
\vspace{0.1in}
\begin{tabular}{@{}L{2.2in} L{9.5in}@{}}
\toprule
\cabt{PRIMER SUSTO} & Una medición cambió después de aprobada, y nadie sabe quién fue. \\
\midrule
\cabt{SEGUNDO SUSTO} & Un certificado se emitió dos veces sobre la misma hectárea. \\
\bottomrule
\end{tabular}

\vspace{0.45in}
{\tDos Orden de prioridad: \dato{integridad}, confidencialidad, disponibilidad.}

\vspace{0.12in}
Cuando dos pilares chocan, gana la integridad: sin registro en bitácora no hay acción.
\note{Qué decir: ninguno de los dos incidentes fue un robo de información. En los dos, el dato dejó de
ser confiable. Por eso la integridad va primero. La confidencialidad va segunda por la competencia
entre dueños y por los datos personales. La disponibilidad va tercera, con un matiz: importa más que el
dato exista en diez años que un día sin sistema.

Si preguntan por un ejemplo de «gana la integridad»: si la bitácora no se puede escribir, la operación
se rechaza aunque eso detenga el trabajo (POL-BT-06).}
\end{frame}

%--------------------------- 4. ACTIVOS ------------------------------
\begin{frame}{Lo que más vale es lo que sostiene un certificado}
\quien{Randall}
\vspace{0.05in}
\begin{tabular}{@{}L{4.6in} L{2.6in} L{4.4in}@{}}
\toprule
\cabt{ACTIVO} & \cabt{CLASIFICACIÓN} & \cabt{LO QUE NO PUEDE FALLAR} \\
\midrule
Mediciones y certificados & Confidencial & \dato{Integridad} \\
Catálogo de fincas: la cartera & Confidencial & Integridad y confidencialidad \\
Datos personales de los dueños & Confidencial & Confidencialidad \\
Bitácora y llave de sello & Restringido & \dato{Integridad} \\
\bottomrule
\end{tabular}

\vspace{0.3in}
\pie{El inventario completo tiene once activos, cada uno valorado en los tres pilares (informe de políticas, sección 2).}
\note{Qué decir: clasificamos en cuatro niveles (público, interno, confidencial, restringido) y valoramos
cada activo con la escala de FIPS 199. La Gerencia de Certificación es la responsable del dato; el
administrador y el operador son custodios.

Si preguntan por un activo que no es obvio: los parámetros de cálculo (A-10). Un factor alterado infla
todos los certificados a la vez. Y la llave de sello (A-09): quien la tenga fabrica certificados
creíbles.}
\end{frame}

%--------------------------- 5. AMENAZAS -----------------------------
\begin{frame}{\textcolor{acento}{7 de 11} amenazas son internas:\\ el atacante más probable ya tiene una cuenta}
\quien{Javier}
\vspace{0.05in}
\begin{tabular}{@{}L{1.4in} L{6.6in} L{3.6in}@{}}
\toprule
\cabt{ID} & \cabt{YA OCURRIÓ EN VERDEAZUL} & \cabt{PILAR} \\
\midrule
AM-03 & Una medición editada después de aprobarse & Integridad \\
AM-06 & Un certificado emitido dos veces & Integridad \\
AM-04 & Cientos de fincas consultadas en un fin de semana & Confidencialidad \\
\bottomrule
\end{tabular}

\vspace{0.3in}
El diseño no mira el perímetro: mira qué puede hacer cada persona autorizada.
\note{Qué decir: el modelo tiene once amenazas, cada una con agente, motivo, activo, la debilidad que
hoy lo permite y el riesgo. Siete son internas: técnico que infla, misma persona que mide y aprueba,
edición silenciosa, extracción de la cartera, administrador u operador que abusa, doble emisión por
error y, desde la consulta al cliente, el empleado con un vínculo con un dueño (AM-11): no existe un
registro de esos vínculos. Tres ya ocurrieron.

Externas: dueño que quiere ver a un competidor, robo de credenciales, ransomware y certificado
falsificado o cuenta de auditor tomada.

Si preguntan cómo calculamos el riesgo: probabilidad por impacto. Probabilidad alta es que ya ocurrió
aquí o que es la vía habitual de ataque en el país.}
\end{frame}

%--------------------------- 6. MARCO LEGAL --------------------------
\begin{frame}{Cada obligación legal termina\\ en una regla que se puede comprobar}
\quien{Javier}
\vspace{0.05in}
\begin{tabular}{@{}L{2.9in} L{4.2in} L{4.5in}@{}}
\toprule
\cabt{NORMA} & \cabt{EXIGE} & \cabt{REGLA EN VERDEAZUL} \\
\midrule
Ley 8968, art. 14 & Autorización para transferir datos & El auditor consulta en la plataforma \\
Reglamento, art. 38 & Avisar una fuga en 5 días hábiles & Procedimiento con ese plazo \\
Ley 7975, art. 2 & Medidas razonables de secreto & Acceso por necesidad y registro \\
ISO/IEC 17029 & Quien revisa no participó & \dato{Medir y aprobar, separados} \\
\bottomrule
\end{tabular}

\vspace{0.2in}
\pie{El cliente aclaró que no audita bajo una norma determinada: la ISO se cita como referencia. Su exigencia firme es la trazabilidad de diez años.}
\note{Qué decir: la Ley 8968 aplica porque los datos de los dueños se comparten con auditores externos;
no es una base de uso exclusivamente interno. Derechos del titular en cinco días hábiles; consentimiento
informado al dar de alta a un dueño; conservación hasta diez años y después se desasocia.

La Ley 7975 es la que permite reclamar si la cartera llega a la competencia, pero solo si la empresa
puede demostrar que tomó medidas razonables: por eso el acceso por asignación y la bitácora tienen
además valor legal.

También: Código Penal (arts. 196 bis, 217 bis, 229 bis y 231), Código de Trabajo art. 71 g) y Ley 8454
de firma digital.

Auditores, según la minuta (P-18): una vez al año, varios días, en sitio o a distancia. Pueden exportar
evidencia, pero como paquetes que genera el sistema, sellados, registrados y sin datos personales. El
acceso vence con la ventana autorizada; prorrogarlo es una autorización nueva.}
\end{frame}

%--------------------------- 7. MODELO: TRES CAPAS -------------------
\begin{frame}{Roles con restricciones,\\ y condiciones que solo pueden quitar}
\quien{Julio}
\vspace{0.2in}
\begin{tikzpicture}[x=1in,y=1in,
  capa/.style={draw=acento,line width=1.4pt,fill=acentoclaro,minimum width=2.4in,minimum height=1.1in,
               align=center,inner sep=3pt},
  ext/.style={draw=regla,line width=0.8pt,fill=plano,minimum height=1.1in,align=center,inner sep=3pt},
  fl/.style={-{Stealth[length=9pt,width=8pt]},draw=gris,line width=1.2pt},
  no/.style={-{Stealth[length=9pt,width=8pt]},draw=gris,line width=1.0pt,densely dashed}]
\node[ext,minimum width=1.6in] (p) at (0.8,0) {Petición\\autenticada};
\node[capa] (r) at (3.25,0) {\textbf{1 $\cdot$ Roles}\\{\tCuatro ¿Este rol puede?}};
\node[capa] (c) at (6.10,0) {\textbf{2 $\cdot$ Condiciones}\\{\tCuatro ¿Aquí y ahora?}};
\node[capa] (h) at (8.95,0) {\textbf{3 $\cdot$ Historial}\\{\tCuatro ¿Encaja con lo previo?}};
\node[ext,minimum width=1.5in] (d) at (11.35,0) {Conceder};
\draw[fl] (p) -- (r);
\draw[fl] (r) -- node[above,font=\tCuatro,text=gris]{sí} (c);
\draw[fl] (c) -- node[above,font=\tCuatro,text=gris]{sí} (h);
\draw[fl] (h) -- node[above,font=\tCuatro,text=gris]{sí} (d);
\node[draw=regla,line width=0.8pt,fill=negado,minimum width=8.1in,minimum height=0.6in] (n) at (6.10,-1.55)
  {Denegar y registrar en bitácora};
\foreach \k in {r,c,h} \draw[no] (\k.south) -- node[right,font=\tCuatro,text=gris]{no} (\k.south |- n.north);
\end{tikzpicture}

\vspace{0.3in}
\dato{Solo la primera capa concede.} La matriz muestra lo máximo que cada rol puede llegar a hacer.
\note{Qué decir: elegimos RBAC con restricciones, nivel RBAC2 de Sandhu. Sobre esa base, un conjunto
cerrado de atributos decide si el permiso aplica a esta finca, en este estado y en este momento.
Es el enfoque centrado en roles de Kuhn, Coyne y Weil.

La frontera del híbrido: la capa 1 dice si el rol tiene el verbo; la capa 2 si aplica a este objeto ahora;
la capa 3 si encaja con el historial. Se combinan sin contradecirse porque solo la tabla de roles concede
y las otras dos son filtros: una negativa en cualquier capa gana.

Si preguntan por qué RBAC: la carta describe funciones, no personas; la gente rota; y hay que poder
explicarlo a un auditor leyendo una fila.}
\end{frame}

%--------------------------- 8. DESCARTES ----------------------------
\begin{frame}{Los otros tres modelos fallan\\ por razones de este caso}
\quien{Julio}
\vspace{0.05in}
\begin{tabular}{@{}L{2.2in} L{9.5in}@{}}
\toprule
\cabt{MODELO} & \cabt{POR QUÉ NO EN VERDEAZUL} \\
\midrule
DAC & Es lo que hay hoy: las hojas compartidas donde ocurrieron los dos sustos. \\
MAC & No hay niveles de secreto. Haría falta un compartimento por finca. \\
ABAC puro & Desaparece la fila que responde «¿qué puede hacer un verificador?». \\
\bottomrule
\end{tabular}

\vspace{0.3in}
\pie{De cada uno se conserva algo: la titularidad como atributo (DAC), la política que el usuario no puede
alterar (MAC) y las condiciones acotadas (ABAC).}
\note{DAC: el técnico que crea la medición sería su dueño y podría seguir modificándola o compartirla.
Es el primer susto.

MAC: Bell-LaPadula ordena por niveles de secreto y protege la confidencialidad. Aquí el riesgo principal
es de integridad. La variante de integridad, Biba, pone una etiqueta fija, pero la confianza en una
medición cambia con su estado y depende de quién hizo el paso anterior. Es el argumento de Clark y
Wilson. MAC sí lo usamos en la plataforma: SELinux confina los contenedores.

ABAC puro: podría expresarlo todo, pero no se puede revisar leyendo una fila, y la empresa no tiene
quién gobierne un catálogo de atributos.}
\end{frame}

%--------------------------- 9. NIVEL Y EXPLOSION --------------------
\begin{frame}{\textcolor{acento}{Siete} roles hoy, y siete con el doble de fincas}
\quien{Julio}
\vspace{0.25in}
{\tDos El alcance nunca es un rol: es un atributo.}

\vspace{0.35in}
\begin{tabular}{@{}L{3.3in} L{8.4in}@{}}
\toprule
RBAC$_2$, sin jerarquía & Si el verificador heredara del técnico, podría registrar: se rompe la separación. \\
Explosión de roles & Aparece al usar un rol para decir dónde o cuándo. Aquí eso vive en las condiciones. \\
La consulta lo probó & Cinco necesidades nuevas: \dato{un rol y cuatro atributos}. \\
\bottomrule
\end{tabular}
\note{Los siete roles: técnico, verificador, supervisor de campo, gerencia de certificación,
administrador de catálogo y cuentas, auditor externo y dueño de finca. Más una identidad de servicio
(SIS), el operador de plataforma, que no tiene rol en la aplicación, y el público sin autenticar.

El séptimo rol salió de la consulta: el cliente dijo que las fincas las asignan los supervisores del
personal de campo (minuta, P-05). Las otras cuatro necesidades nuevas (el verificador que mide por
excepción, el técnico que apoya a un compañero, la sociedad con un representante y el empleado con un
vínculo con un dueño) no son funciones de trabajo: se resolvieron con atributos.

Por qué RBAC2 y no RBAC3: no hay cadena de mando entre las siete funciones; la herencia vuelve
implícitos los permisos y la matriz no debe tener nada implícito; y con siete roles no ahorra nada.

Restricciones: un rol por persona, una identidad por persona, al menos dos verificadores, cuenta de
auditor con ventana. La pareja técnico--verificador es la única excepción: es dinámica. Un verificador
puede recibir una habilitación temporal de técnico, con fecha de fin, y nunca aprueba lo que midió
(minuta, P-01).

Si preguntan dónde está el riesgo de explosión: un rol por finca para aislar a los dueños, o roles como
«Técnico Zona Sur 2027».}
\end{frame}

%--------------------------- 10. HBAC --------------------------------
\begin{frame}{Cuatro controles de historial, uno de cada tipo}
\quien{Axel}
\vspace{0.0in}
\begin{tabular}{@{}L{0.7in} L{1.7in} L{4.6in} L{4.4in}@{}}
\toprule
 & \cabt{TIPO} & \cabt{QUÉ LO DISPARA} & \cabt{QUÉ DECIDE} \\
\midrule
H1 & Frecuencia & Muchas fincas distintas en un día & \dato{Tres niveles}: avisa, confirma, restringe \\
H2 & Tiempo & Una decisión fuera de jornada & Pide segundo factor y motivo \\
H3 & Combinada & Fuera de jornada, muchas fincas, casi todas nuevas & Confirma y avisa; si insiste, restringe \\
H4 & Secuencia & Quien participó antes quiere decidir & \dato{Deniega} \\
\bottomrule
\end{tabular}

\vspace{0.2in}
\pie{Ninguna señal aislada restringe: el cliente pidió avisar primero y no bloquear de inmediato (minuta, P-10).}
\note{Son reglas explícitas, no aprendizaje automático: se le pueden leer a un auditor y se pueden probar.

Cómo se distingue lo legítimo: se cuentan señales: volumen, novedad, momento e insistencia. El trabajo
extraordinario legítimo muestra una; la extracción de la cartera las muestra juntas.

Tres niveles: aviso (el usuario no nota nada), verificación adicional (segundo factor y motivo) y
restricción temporal (no abre fincas nuevas hasta el siguiente día hábil; la cuenta sigue activa).

H1 mitiga AM-04. Umbrales por rol, calibrados con datos reales; iniciales 40, 60 y 80 fincas por día
(supuesto P-13). Los intentos denegados cuentan doble.

H2 mitiga AM-08, la cuenta robada usada de madrugada. H4 es la separación de responsabilidades como
control de historial: reglas C3, C8 y C10.

H3 es la firma del incidente: fin de semana, cientos de fincas, ninguna suya. Pide confirmación en la
finca once y avisa; restringe en la veintiuna.

Falsos positivos: verificador con la cola acumulada al cierre, o nuevo y sin historial. El primer
período corre en observación para calibrar.

Descartada: Muralla China. Todos los dueños compiten: cada empleado quedaría con una sola finca.}
\end{frame}

%--------------------------- 11. EL TECNICO --------------------------
\begin{frame}{El técnico del fin de semana\\ hoy no vería ni una finca ajena}
\quien{Axel}
\vspace{0.05in}
\begin{tabular}{@{}L{4.6in} L{7.1in}@{}}
\toprule
\cabt{LO QUE HACE} & \cabt{LO QUE HACE EL SISTEMA} \\
\midrule
Abre una finca que no tiene asignada & \dato{Deniega}: su rol no ve fincas ajenas \\
Insiste tres, cinco, diez veces & Avisa a su supervisor; pide segundo factor; lo detiene media hora \\
La necesitaba de verdad & El supervisor le da una asignación temporal, que vence sola \\
Presenta su renuncia & Reporte de sus consultas de 90 días \\
\bottomrule
\end{tabular}

\vspace{0.25in}
\pie{El cliente lo confirmó: fincas asignadas con flexibilidad, y respuesta según la gravedad (minuta, P-10 y P-11).}
\note{Qué decir: la carta dice «no sabemos qué hacer con eso, pero nos gustaría que el sistema al menos se
diera cuenta». Nuestra primera respuesta no es vigilar el permiso: es quitarlo. Un técnico solo ve las
fincas que tiene asignadas (regla C1).

El cliente lo confirmó, con una condición: flexibilidad para apoyar a un compañero, sustituir a un
ausente o atender una visita extraordinaria, sin acceso permanente a todas las fincas. Se resuelve con
asignaciones temporales que crea el supervisor de campo, con motivo y fecha de fin (regla C19).

El diseño original suspendía la sesión a la quinta denegación. El cliente pidió no bloquear de
inmediato, y se cambió: aviso al supervisor a la tercera, segundo factor a la quinta, media hora de
restricción a la décima.

Los controles de historial quedan para quienes sí necesitan leer muchas fincas: verificador, supervisor,
gerencia, administrador y auditor. Si el técnico hubiera sido verificador, H3 le pide confirmación en la
finca once y lo restringe en la veintiuna.

Además hay una vía legal: confidencialidad firmada, que persiste después de la salida (Ley 8968 art. 11
y Ley 7975).}
\end{frame}

%--------------------------- 12. NADIE EMITE -------------------------
\begin{frame}{Nadie emite un certificado: lo hace el sistema}
\quien{Randall}
\vspace{0.25in}
\begin{tikzpicture}[x=1in,y=1in,
  paso/.style={draw=acento,line width=1.4pt,fill=acentoclaro,minimum width=3.3in,minimum height=1.1in,
               align=center,inner sep=3pt},
  sis/.style={draw=gris,line width=1.2pt,fill=negado,minimum width=3.3in,minimum height=1.1in,
              align=center,inner sep=3pt},
  fl/.style={-{Stealth[length=9pt,width=8pt]},draw=gris,line width=1.2pt},
  nota/.style={font=\tCuatro,text=gris,align=center}]
\node[paso] (a) at (1.65,0) {\textbf{Técnico}\\registra y envía};
\node[paso] (b) at (6.05,0) {\textbf{Verificador}\\aprueba o devuelve};
\node[sis]  (c) at (10.45,0) {\textbf{Sistema}\\emite};
\draw[fl] (a) -- (b); \draw[fl] (b) -- (c);
\node[nota,below=0.08in of a] {solo en fincas asignadas\\y con el período abierto};
\node[nota,below=0.08in of b] {no registró ninguna versión\\de esa medición};
\node[nota,below=0.08in of c] {no hay otro certificado vigente\\para la parcela y el período};
\end{tikzpicture}

\vspace{0.35in}
Ninguna persona tiene el permiso de emitir. \dato{Hacen falta dos, y distintas.}
\note{La acción crítica: emitir un certificado. Si una sola persona pudiera medir y aprobar, podría
inflar la biomasa, aprobar su propio dato y obtener un certificado por toneladas que no existen.

La segunda acción crítica: corregir una medición ya congelada. Es el primer susto, y se trata igual,
con el orden que pidió el cliente (minuta, P-07): solicita el técnico responsable, un verificador o la
Gerencia; corrige el técnico; revisa un verificador que no es autor ni solicitante; y la Gerencia da la
autorización final. Son al menos tres personas.

Si preguntan qué pasa con el certificado tras una enmienda (minuta, P-08): si las toneladas no cambian,
sigue vigente con la enmienda anotada. Si cambian, el sistema lo revoca y emite el reemplazo en un solo
acto, enlazados: nunca hay dos vigentes y ninguno se borra. Si ya se vendió, la Gerencia lo registra y
coordina con las partes.}
\end{frame}

%--------------------------- 13. POR MEDICION ------------------------
\begin{frame}{La separación se comprueba\\ por medición, no por puesto}
\quien{Randall}
\vspace{0.05in}
\begin{tabular}{@{}L{3.4in} L{8.3in}@{}}
\toprule
\cabt{SI SE COMPRUEBA\ldots} & \cabt{UN TÉCNICO ASCENDIDO A VERIFICADOR\ldots} \\
\midrule
por el puesto de hoy & aprueba lo que él mismo midió el año anterior. \\
\dato{por la medición} & \dato{no puede: figura como autor.} \\
\bottomrule
\end{tabular}

\vspace{0.35in}
Técnica: \textbf{separación por etapas del flujo}, reforzada con roles mutuamente excluyentes.
\note{Qué decir: quien inicia el proceso no puede cerrarlo. El sistema no pregunta qué rol tiene la
persona, pregunta si participó en esta medición. En la literatura es la separación basada en el objeto
o en el historial (Nash y Poland 1990; Simon y Zurko 1997).

Por qué es la adecuada: apunta al riesgo tal como lo dijo el cliente; sobrevive a los cambios de puesto
(la lección de Kerviel); no depende de que los roles estén bien administrados; y deja dos nombres
distintos en cada certificado.

El cliente confirmó que algunos verificadores fueron técnicos y que por excepción pueden medir, y pidió
«garantizar una revisión independiente, sin impedir que el personal pueda asumir otras
responsabilidades» (minuta, P-01). Una separación por puesto obligaría a escoger; esta da las dos
cosas: el verificador mide con una habilitación temporal y no podrá aprobar nunca lo que midió.

Coincide con el principio de imparcialidad de ISO/IEC 17029, que citamos solo como referencia.}
\end{frame}

%--------------------------- 14. ALTERNATIVAS ------------------------
\begin{frame}{Las alternativas se descartaron\\ por lo que cuestan aquí}
\quien{Randall}
\vspace{0.0in}
\begin{tabular}{@{}L{4.3in} L{7.4in}@{}}
\toprule
\cabt{TÉCNICA} & \cabt{POR QUÉ NO ES LA PRINCIPAL} \\
\midrule
Roles excluyentes, solos & No cubren a quien cambia de puesto \\
Doble control (cuatro ojos) & \dato{Duplica la verificación de cada certificado} \\
Segregación temporal & Limita cuándo, no quién \\
Rotación & No impide un acto; se recomienda contra la colusión \\
\bottomrule
\end{tabular}

\vspace{0.2in}
\pie{La segregación temporal sí se usa donde el problema es la duración: la asignación temporal, la habilitación del verificador, el plazo de la enmienda y la cuenta del auditor.}
\note{Roles mutuamente excluyentes estáticos (RBAC2): se conservan como refuerzo. Li, Tripunitara y Bizri
muestran que excluir roles es un mecanismo para aproximar la separación, no la separación misma.

Exclusión dinámica: basta cerrar sesión y abrir otra. Sola, demora el fraude y no lo impide. Se usa
como refuerzo en la pareja técnico--verificador, detrás de la regla C3.

Doble control: era la opción más fuerte contra el fraude. Se descartó por costo operativo: duplica la
verificación de cada certificado, con cualquier tamaño de equipo.

Rotación: cuesta conocimiento de la finca y depende de cuántas personas haya, un dato que el cliente
pidió no fijar (minuta, P-02).}
\end{frame}

%--------------------------- 15. FRICCION ----------------------------
\begin{frame}{La fricción existe y está acotada}
\quien{Javier}
\vspace{0.0in}
\begin{tabular}{@{}L{5.2in} L{6.5in}@{}}
\toprule
\cabt{LO QUE CUESTA} & \cabt{CÓMO SE MANEJA} \\
\midrule
Siempre hacen falta dos personas & Mínimo dos verificadores; \dato{sin atajo de emergencia} \\
Las devoluciones quitan tiempo de campo & Validación de rangos al registrar \\
Una enmienda pide tres personas & Plazo de diez días; el sistema la clasifica \\
La Gerencia concentra las segundas firmas & Nunca inicia lo que aprueba \\
\bottomrule
\end{tabular}
\note{Qué decir: no existe un atajo para aprobar lo propio. Un atajo así sería el hueco que el auditor
buscaría. Si el único verificador presente participó, la medición espera; la cola avisa a la Gerencia.

Lo que la técnica no resuelve: la colusión. Si un técnico y un verificador se ponen de acuerdo, el flujo
se cumple. La separación sube el costo del fraude de una persona a dos. Quedan la rotación, un indicador
de parejas repetidas y la auditoría externa.

La Gerencia es un cuello de botella y un punto de confianza: firma en varios flujos, pero nunca inicia
lo que aprueba, no registra ni aprueba mediciones, ya no asigna fincas y no resuelve sus propias alertas.

Enmienda: revisa un verificador distinto del original solo cuando es material (cambia toneladas o ya hay
certificado); lo determina el sistema. Con solo dos verificadores, si uno solicita sobre lo que aprobó
el otro no queda revisor: la presenta el técnico o la Gerencia.

La flexibilidad también cuesta un trámite: la asignación temporal la crea el supervisor y la
habilitación de un verificador para medir pasa por el administrador y la Gerencia.}
\end{frame}

%--------------------------- 16. MATRIZ (1 de 2) ---------------------
\begin{frame}{Lo que no está escrito está denegado}
\quien{Julio}
\vspace{-0.02in}
{\tCuatro
\renewcommand{\arraystretch}{1.3}
\setlength{\tabcolsep}{5pt}
\begin{tabular}{@{}L{0.75in}|L{1.95in}|L{2.25in}|L{2.0in}|L{1.85in}|L{2.3in}@{}}
\rowcolor{acentoclaro}
 & \textbf{Fincas y parcelas} & \textbf{Datos personales del dueño} & \textbf{Períodos de medición} & \textbf{Asignaciones} & \textbf{Mediciones} \\
\hline
\suj{TEC} & \vb{leer}{C1} & \vb{leer}{C1} (solo contacto) & leer & \vb{leer}{C1} (las propias) & \vb{leer}{C1}, \vb{crear}{C2}, \vb{modificar}{C2}, \vb{enviar}{C2} \\
\hline
\suj{VER} & \vb{leer}{H} & \sinacc & leer & \sinacc & \vb{leer}{H}, \vb{aprobar}{C3}, \vb{devolver}{C3} \\
\hline
\suj{SUP} & \vb{leer}{H} & \sinacc & leer & leer, \vb{crear}{C19}, \vb{modificar}{C19} & \vb{leer}{H} (solo estado y fechas) \\
\hline
\suj{GER} & \vb{leer}{H}, \vb{aprobar}{C11} & \vb{leer}{H} & leer, \vb{abrir}{C14}, \vb{cerrar}{C14} & leer & \vb{leer}{H} \\
\hline
\suj{ADM} & \vb{leer}{H}, \vb{proponer}{C11} & \vb{leer}{H}, crear, modificar, \vb{eliminar}{C16} & leer, crear, \vb{proponer}{C14} & \sinacc & \sinacc \\
\hline
\suj{AUD} & \vb{leer}{C6}, \vb{exportar}{C18} & \vb{leer}{C6} & \vb{leer}{C6}, \vb{exportar}{C18} & \vb{leer}{C6}, \vb{exportar}{C18} & \vb{leer}{C6}, \vb{exportar}{C18} \\
\hline
\suj{DUE} & \vb{leer}{C5} & \vb{leer}{C5} & leer & \vb{leer}{C5} & \vb{leer}{C5} (solo aprobadas) \\
\hline
\suj{SIS} & leer & leer (solo el nombre del titular) & leer & leer & leer \\
\hline
\suj{OPE} & \sinacc & \sinacc & \sinacc & \sinacc & \sinacc \\
\hline
\suj{PUB} & \sinacc & \sinacc & \sinacc & \sinacc & \sinacc \\
\end{tabular}

\vspace{0.06in}
{\color{gris}Matriz real, objetos 1 a 5 de 10. \quad --- sin acceso \quad {\color{acento}\textbf{C}$n$} permiso sujeto a la
regla $n$ \quad {\color{acento}\textbf{H}} lectura vigilada por historial}
}
\note{Cómo leerla: sujetos en las filas, objetos en las columnas. Una celda gris es sin acceso. Una marca C
significa que el permiso solo vale si se cumple esa regla; las condiciones solo pueden quitar.

Recorrido sugerido: fila TEC, columna Mediciones: crea, modifica y envía, pero no aprueba. Fila VER:
aprueba y devuelve con C3, no crea ni modifica, y no ve los datos personales del dueño. Fila SUP, que
es nueva: asigna fincas (C19) y de las mediciones solo ve estado y fechas. Fila GER: abre y cierra
períodos que registró el administrador (C14); ya no crea asignaciones. Fila AUD: leer y exportar
evidencia, con ventana (C6 y C18), y ningún verbo de escritura. Fila DUE: solo lo suyo (C5). Filas OPE
y PUB: nada en esta mitad.

Diez sujetos: siete roles, el sistema, el operador y el público sin autenticar.

Mínimo privilegio: cada permiso amplio tiene su justificación en la sección 6 del documento de la matriz.}
\end{frame}

%--------------------------- 17. MATRIZ (2 de 2) ---------------------
\begin{frame}{Nadie tiene «eliminar», y nadie emite}
\quien{Julio}
\vspace{-0.02in}
{\tCuatro
\renewcommand{\arraystretch}{1.3}
\setlength{\tabcolsep}{5pt}
\begin{tabular}{@{}L{0.75in}|L{2.3in}|L{1.9in}|L{2.75in}|L{1.4in}|L{2.0in}@{}}
\rowcolor{acentoclaro}
 & \textbf{Enmiendas} & \textbf{Certificados} & \textbf{Cuentas y roles} & \textbf{Bitácora} & \textbf{Alertas de historial} \\
\hline
\suj{TEC} & \vb{leer}{C1}, \vb{solicitar}{C7}, \vb{corregir}{C9} & \sinacc & \vb{modificar}{C13} & \sinacc & \sinacc \\
\hline
\suj{VER} & leer, \vb{solicitar}{C7}, \vb{revisar}{C10} & \vb{leer}{H} & \vb{modificar}{C13} & \sinacc & \sinacc \\
\hline
\suj{SUP} & leer (solo estado) & \sinacc & \vb{modificar}{C13} & \sinacc & \vb{leer}{C15}, \vb{resolver}{C15} \\
\hline
\suj{GER} & leer, \vb{solicitar}{C7}, \vb{autorizar}{C8} & \vb{leer}{H} & leer, \vb{aprobar}{C12}, \vb{modificar}{C13} & leer & leer, \vb{resolver}{C15} \\
\hline
\suj{ADM} & \sinacc & \sinacc & leer, \vb{crear}{C12}, \vb{desactivar}{C12}, \vb{proponer}{C12}, \vb{modificar}{C13} & \sinacc & \sinacc \\
\hline
\suj{AUD} & \vb{leer}{C6}, \vb{exportar}{C18} & \vb{leer}{C6}, \vb{exportar}{C18} & \vb{leer}{C6}, \vb{modificar}{C13} & \vb{leer}{C6}, \vb{exportar}{C18} & \vb{leer}{C6}, \vb{exportar}{C18} \\
\hline
\suj{DUE} & \vb{leer}{C5} & \vb{leer}{C5} & \vb{modificar}{C13} & \sinacc & \sinacc \\
\hline
\suj{SIS} & leer & \dato{emitir C4}, \vb{revocar}{C4} & \sinacc & anexar & crear \\
\hline
\suj{OPE} & \sinacc & \sinacc & \sinacc & \sinacc & \sinacc \\
\hline
\suj{PUB} & \sinacc & \vb{comprobar}{C20} & \sinacc & \sinacc & \sinacc \\
\end{tabular}

\vspace{0.06in}
{\color{gris}Matriz real, objetos 6 a 10 de 10. \quad Toda decisión exige segundo factor reciente. \quad {\color{acento}\textbf{C17}}: nadie decide sobre una finca con cuyo dueño declaró un vínculo.}
}
\note{Dos cosas que se ven en esta mitad. Una: el verbo «eliminar» no aparece en enmiendas, certificados,
cuentas ni bitácora. Nada se borra; todo cambio es una versión nueva o un cambio de estado. Dos: «emitir»
solo está en la fila SIS. Ninguna fila de persona lo tiene.

Columna Enmiendas: solicitar (técnico, verificador o Gerencia), corregir (técnico), revisar
(verificador) y autorizar (Gerencia), con las reglas C7 a C10. Cuatro pasos, al menos tres personas.

Columna Cuentas y roles: C13 es «credenciales y declaración propias»; C12 es «nadie se administra a sí
mismo»: el administrador propone el rol y la Gerencia lo aprueba.

Columna Bitácora: solo se agrega. La leen la Gerencia y el auditor.

Fila PUB: quien tenga el código de un certificado comprueba si es auténtico y si sigue vigente, y nada
más (C20). Sale de la respuesta del cliente sobre autenticidad y estado del certificado (minuta, P-04).

C17, abstención por vínculo: quien declaró un vínculo familiar o comercial con un dueño no mide, no
aprueba, no asigna ni autoriza nada sobre esa finca (minuta, P-15).

Al lado de la matriz hay una tabla de exclusiones: cada persona tiene un solo rol, salvo la habilitación
temporal de un verificador como técnico.}
\end{frame}

%--------------------------- 18. TRES CELDAS -------------------------
\begin{frame}{Tres celdas bastan para leer la separación}
\quien{Julio}
\vspace{0.0in}
\begin{tabular}{@{}L{1.0in} L{5.9in} L{4.8in}@{}}
\toprule
 & \cabt{MEDICIONES} & \cabt{CERTIFICADOS} \\
\midrule
\textbf{TEC} & Crea, modifica y envía. Solo sus borradores, con el período abierto. & Sin acceso \\
\textbf{VER} & \dato{Aprueba solo si no registró ninguna versión.} & Lee \\
\textbf{SIS} & Lee & \dato{Emite solo si no hay otro vigente.} \\
\bottomrule
\end{tabular}

\vspace{0.25in}
\pie{Caso de prueba T-06: la verificadora que antes fue técnica intenta aprobar una medición que registró. Resultado esperado: deniega.}
\note{Son las reglas C2, C3 y C4. Cada regla de la matriz tiene un caso de prueba con resultado esperado:
treinta y siete casos, T-01 a T-37. El equipo que construya debe convertirlos en pruebas automáticas, más una por
cada celda sin acceso.

Otros casos para citar: T-07, segunda emisión sobre la misma parcela y período: no emite. T-08, el dueño
cambia el identificador en la dirección: deniega y responde «no encontrado». T-10, el auditor envía una
petición de modificación: deniega.

Caso nuevo para citar: T-17, después del cierre el verificador aprueba una medición que se envió a
tiempo: concede. El cliente pidió distinguir registrar, modificar, revisar y aprobar (minuta, P-06).

El mismo patrón, quien inicia no cierra, se repite en enmiendas (C7 a C10), períodos (C14),
asignaciones (C19), catálogo (C11) y cuentas (C12).}
\end{frame}

%--------------------------- 19. ARQUITECTURA ------------------------
\begin{frame}{Un solo punto decide, en el servidor,\\ antes de tocar el dato}
\quien{Axel}
\vspace{0.1in}
\begin{tikzpicture}[x=1in,y=1in,
  comp/.style={draw=regla,line width=0.8pt,fill=white,align=center,inner sep=3pt,minimum height=0.8in},
  pdp/.style={draw=acento,line width=1.6pt,fill=acentoclaro,align=center,inner sep=3pt,minimum height=0.8in},
  zona/.style={draw=gris,densely dashed,line width=0.7pt},
  tz/.style={font=\tCuatro\bfseries,text=gris,anchor=north west},
  dos/.style={{Stealth[length=8pt,width=7pt]}-{Stealth[length=8pt,width=7pt]},draw=gris,line width=1.1pt}]
\draw[zona] (0,0) rectangle (2.0,3.3);
\draw[zona] (2.35,0) rectangle (3.95,3.3);
\draw[zona] (4.3,0) rectangle (8.9,3.3);
\draw[zona] (9.25,0) rectangle (12.1,3.3);
\node[tz] at (0.08,3.25) {NO CONFIABLE};
\node[tz] at (2.43,3.25) {BORDE};
\node[tz] at (4.38,3.25) {APLICACIÓN};
\node[tz] at (9.33,3.25) {DATOS Y SECRETOS};
\node[comp,minimum width=1.6in] (nav) at (1.0,1.25) {Navegador};
\node[comp,minimum width=1.2in] (px) at (3.15,1.25) {Proxy};
\node[comp,minimum width=4.0in] (app) at (6.6,1.25) {\textbf{Aplicación}\\{\tCuatro identifica, valida y hace cumplir}};
\node[pdp,minimum width=4.0in] (opa) at (6.6,2.45) {\textbf{Motor de decisión}\\{\tCuatro la matriz, como política}};
\node[comp,minimum width=2.4in] (pg) at (10.68,2.1) {\textbf{Base de datos}\\{\tCuatro segundo candado}};
\node[comp,minimum width=2.4in] (bao) at (10.68,0.75) {\textbf{Bóveda}\\{\tCuatro secretos y sello}};
\draw[dos] (nav) -- (px);
\draw[dos] (px) -- (app);
\draw[dos,draw=acento] (app) -- (opa);
\draw[dos] ($(app.east)+(0,0.18)$) -- ++(0.3,0) |- (pg.west);
\draw[dos] ($(app.east)+(0,-0.18)$) -- ++(0.3,0) |- (bao.west);
\end{tikzpicture}

\vspace{0.2in}
\dato{Esconder un botón no es autorización.} Si el servidor no decide, cualquiera llama directo.
\note{Flujo de una petición: el proxy termina la conexión cifrada; la aplicación identifica la sesión,
valida la forma de los datos y lee de la base los hechos; entrega al motor el sujeto, la acción, el
objeto y esos hechos; el motor devuelve la decisión con la regla que la determinó. Solo entonces se toca
el dato, en una transacción que también escribe la bitácora. La base aplica sus propias restricciones.

Fronteras de confianza: nada de lo que llega del navegador se da por bueno; los hechos se toman de la
base, nunca de la petición; si el motor no responde, se deniega; la base no confía en la aplicación.

Si preguntan qué pasa si el control está solo en la interfaz: el técnico envía a mano la petición de
aprobar su propia medición y el servidor la ejecuta. Es el error de referencia directa a objetos de la
clase 13.

Hospedaje, que el cliente dejó a nuestro criterio (minuta, P-23): un servidor virtual contratado, operado
por personal que la empresa designe. Hay además una ruta pública de solo lectura para comprobar
certificados.}
\end{frame}

%--------------------------- 20. TECNOLOGIAS -------------------------
\begin{frame}{Seis tecnologías libres,\\ elegidas por lo que protegen}
\quien{Axel}
\vspace{-0.05in}
{\renewcommand{\arraystretch}{1.36}
\begin{tabular}{@{}L{2.9in} L{3.7in} L{5.1in}@{}}
\toprule
\cabt{DECISIÓN} & \cabt{ELECCIÓN} & \cabt{LO QUE APORTA} \\
\midrule
Lenguaje y framework & Python 3.14, Django 5.2 LTS & Consultas parametrizadas; escape por omisión \\
Base de datos & PostgreSQL 18 & Seguridad por fila; emisión única \\
Sesión & De servidor, cookie opaca & \dato{Se revoca en el acto} \\
Autorización & Open Policy Agent 1.21 & Un solo lugar; la matriz como política \\
Contenedores & Podman sobre AlmaLinux 10 & Sin privilegios; confinados por SELinux \\
\bottomrule
\end{tabular}}
\note{Para cada una el documento indica versión, alternativas descartadas, aporte, limitaciones y
herramientas de análisis.

Descartes por licencia: HashiCorp Vault (BUSL desde 2023) y Docker Desktop. Se usa OpenBao, que es
software libre.

Descartes por seguridad: tokens JWT, porque no se pueden revocar en el acto al dar de baja o cambiar
de rol; MariaDB, porque no tiene seguridad por fila; Docker Engine, por su servicio con privilegios; una
cadena de bloques, porque resuelve la desconfianza entre organizaciones, que aquí no existe.

Limitaciones que hay que saber decir: los permisos de Django son por tipo de objeto, no por objeto; su
sitio de administración no se instala; el superusuario de PostgreSQL se salta todo; el motor no conoce el
historial y no filtra listados.

Contraseñas: Argon2id con sal por usuario. Segundo factor: llave de acceso o códigos temporales; nunca SMS.}
\end{frame}

%--------------------------- 21. BITACORA ----------------------------
\begin{frame}{La bitácora no la reescribe\\ ni quien administra el servidor}
\quien{Axel}
\vspace{0.2in}
\begin{tikzpicture}[x=1in,y=1in,
  reg/.style={draw=regla,line width=0.8pt,fill=plano,minimum width=2.25in,minimum height=1.0in,
              align=center,inner sep=3pt},
  ext/.style={draw=acento,line width=1.6pt,fill=acentoclaro,minimum width=2.6in,minimum height=1.0in,
              align=center,inner sep=3pt},
  fl/.style={-{Stealth[length=9pt,width=8pt]},draw=gris,line width=1.2pt}]
\draw[draw=gris,densely dashed,line width=0.7pt] (-0.1,-0.8) rectangle (8.55,1.15);
\node[font=\tCuatro\bfseries,text=gris,anchor=north west] at (0,1.1) {SERVIDOR: SOLO SE PUEDE INSERTAR};
\node[reg] (a) at (1.25,-0.05) {Registro 1041\\{\tCuatro huella del 1040}};
\node[reg] (b) at (4.2,-0.05) {Registro 1042\\{\tCuatro huella del 1041}};
\node[reg] (c) at (7.15,-0.05) {Registro 1043\\{\tCuatro huella del 1042}};
\draw[fl] (b) -- (a); \draw[fl] (c) -- (b);
\node[ext,minimum width=2.25in] (t) at (10.97,-0.05) {\textbf{Testigo externo}\\{\tCuatro sello de tiempo y auditor}};
\draw[fl,draw=acento] (c) -- node[above,pos=0.66,font=\tCuatro,text=gris]{cada día} (t);
\end{tikzpicture}

\vspace{0.3in}
Cambiar un registro rompe la cadena. Reescribirla entera no coincide con lo que ya salió.
\note{Cuatro mecanismos que se apoyan: (1) la aplicación solo puede insertar, no modificar ni borrar;
(2) cada registro incluye la huella SHA-256 del anterior; (3) una vez al día la última huella sale del
servidor, recibe un sello de tiempo RFC 3161 y va al buzón de la Gerencia y del auditor; (4) la acción y
su registro se confirman en la misma transacción.

Qué se registra: autenticación, cambios, consultas de datos sensibles, administración y toda decisión
denegada. Con qué campos: instante, cuenta y rol, acción, objeto y versión, resultado y regla, origen,
huellas.

Conservación: diez años para cambios y decisiones, que es la trazabilidad que el cliente debe poder
reconstruir (minuta, P-18); dos para consultas (supuesto P-19).

Secretos: viven en una bóveda; la llave de sello no sale, se le pide a la bóveda que firme. La bóveda se
abre con dos de tres custodios.}
\end{frame}

%--------------------------- 22. RESIDUALES --------------------------
\begin{frame}{Lo que no se cubre está declarado}
\quien{Javier}
\vspace{0.0in}
\begin{tabular}{@{}L{5.0in} L{6.7in}@{}}
\toprule
\cabt{RIESGO RESIDUAL} & \cabt{POR QUÉ SE ACEPTA} \\
\midrule
Colusión entre técnico y verificador & Evitarla exige una tercera persona por certificado \\
Una medición falsa desde el origen & El sistema no ve el árbol \\
Extracción lenta, bajo los umbrales & Bajarlos castiga a quien trabaja bien \\
Un vínculo con un dueño que nadie declaró & No existe un registro que lo descarte \\
El histórico de las hojas de cálculo & No se le puede dar la trazabilidad que no tuvo \\
\bottomrule
\end{tabular}
\note{Son trece riesgos residuales en el documento de arquitectura. Otros: equipo de usuario infectado; el
operador puede detener o destruir (se detecta, no se impide); un solo servidor; sin hardware de
seguridad dedicado; la fecha de medición es declarada, porque no hay aplicación sin conexión; el
proveedor de hospedaje tiene acceso físico; la aprobación no lleva firma digital certificada; y la
misma regla vive en dos lugares, el motor y la base.

Si preguntan qué requerimiento quedó sin cubrir: la trazabilidad de diez años hacia atrás. Las hojas de
cálculo se congelan, se les calcula la huella y lo migrado queda marcado como histórico.

Antes de entrar en operación hay diez criterios de aceptación, entre ellos una prueba de penetración
autorizada por escrito.}
\end{frame}

%--------------------------- 23. RESPUESTAS DEL CLIENTE -------------
\begin{frame}{El cliente respondió,\\ y el diseño cambió en cinco puntos}
\quien{Javier}
\vspace{0.0in}
\begin{tabular}{@{}L{0.9in} L{5.0in} L{5.6in}@{}}
\toprule
 & \cabt{LO QUE RESPONDIÓ} & \cabt{QUÉ CAMBIÓ} \\
\midrule
P-01 & Un verificador puede medir, por excepción & Habilitación temporal; no aprueba lo que midió \\
P-05 & Las fincas las asignan los supervisores & \dato{Rol nuevo: supervisor de campo} \\
P-06 & La revisión sigue después del cierre & El cierre detiene registrar, no aprobar \\
P-07 & Revisión técnica antes de autorizar & Solicitar, corregir, revisar, autorizar \\
P-10 & Avisar primero; no bloquear de inmediato & Aviso, verificación, restricción temporal \\
\bottomrule
\end{tabular}

\vspace{0.2in}
\pie{Catorce de veintiséis consultas, respondidas por escrito. Las otras doce son supuestos declarados. Minuta: informe de políticas, Anexo A.}
\note{Qué decir: el diseño se hizo completo con supuestos declarados, y la consulta sirvió como validación.
De catorce respuestas: dos confirmaron el supuesto (P-03, el verificador devuelve y no edita; P-11, el
técnico solo ve sus fincas), cuatro lo ajustaron, seis lo cambiaron y dos quedaron a nuestro criterio.

Lo que resistió: la técnica de separación. La respuesta a P-01 no la cambió; solo volvió dinámica la
exclusión entre técnico y verificador.

Lo que hubo que rehacer: el orden de la enmienda (P-07), lo que el cierre detiene (P-06) y quién asigna
las fincas (P-05). Y el diseño era más estricto de lo que el cliente quería: bloqueaba de inmediato
donde bastaba avisar (P-10).

Otros cambios para citar si preguntan: sociedades y representantes como dueños, y declaración de
vínculos con abstención (P-15); ninguna norma de auditoría se da por supuesta, y los auditores exportan
evidencia sellada (P-18); registro después de la visita, con fecha de medición y fecha de registro
(P-22); diseño sin cantidades fijas (P-02).

Cada cambio llegó hasta la matriz: diecisiete casos de prueba nuevos, T-21 a T-37, y cinco modificados.}
\end{frame}

%--------------------------- 24. DECISIONES DELEGADAS ----------------
\begin{frame}{Cuatro decisiones que el cliente\\ dejó en nuestras manos}
\quien{Javier}
\vspace{0.0in}
\begin{tabular}{@{}L{4.7in} L{7.0in}@{}}
\toprule
\cabt{LA PREGUNTA} & \cabt{LO QUE DECIDIMOS} \\
\midrule
¿Revisa un verificador distinto del original? & \dato{Sí, cuando la enmienda es material} \\
¿Hace falta firma digital certificada? & No: sello de la empresa y comprobación en línea \\
¿Llaves físicas para todo el personal? & No: llave de acceso del propio equipo \\
¿Dónde se hospeda y quién lo opera? & Servidor virtual contratado; personal designado \\
\bottomrule
\end{tabular}

\vspace{0.2in}
\pie{El cliente pidió que cada decisión tenga una justificación clara y responda a un riesgo real, no el sistema más complejo.}
\note{Verificador distinto (minuta, P-07): tres opciones. Siempre distinto: toda corrección, hasta una de
forma, espera a otra persona. El mismo: quien aprobó el dato erróneo dictamina sobre su propio error.
Elegimos según lo que cambia: es material si cambia un dato del cálculo de toneladas o si ya hay un
certificado; entonces revisa otro. Lo determina el sistema, no una persona.

Firma digital (P-04): la firma certificada de cada verificador da autenticidad y autoría con presunción
legal, pero no dice si el certificado fue revocado después, y cuesta por persona. Elegimos sello de la
empresa, comprobación en línea por código y aprobación personal con segundo factor, que queda en la
bitácora.

Llaves físicas (P-23): no hay presupuesto para todos. La llave de acceso del propio dispositivo
(passkey) resiste igual el correo falso y no cuesta nada. Se recomiendan llaves físicas solo para
Gerencia, administrador, operador y custodios de la bóveda.

Hospedaje (P-23): servidor propio, servidor virtual contratado o servicios administrados, con los
cuatro criterios del cliente. El propio exige inversión y cuidar energía y enlace en una oficina
rural; los administrados no son software libre. Lo cotidiano está automatizado.}
\end{frame}

%--------------------------- 25. REFERENCIAS -------------------------
\begin{frame}{Referencias principales}
\quien{}
\vspace{0.0in}
{\tCuatro
\begin{columns}[T,onlytextwidth]
\begin{column}{0.485\textwidth}
\raggedright
Abadi, M., \& Fournet, C. (2003). Access control based on execution history. \emph{NDSS}.\\[5pt]
Clark, D. D., \& Wilson, D. R. (1987). A comparison of commercial and military computer security
policies. \emph{IEEE S\&P}.\\[5pt]
Kuhn, D. R., Coyne, E. J., \& Weil, T. R. (2010). Adding attributes to role-based access control.
\emph{Computer, 43}(6).\\[5pt]
Li, N., Tripunitara, M. V., \& Bizri, Z. (2007). On mutually exclusive roles and separation-of-duty.
\emph{ACM TISSEC, 10}(2).\\[5pt]
Saltzer, J. H., \& Schroeder, M. D. (1975). The protection of information in computer systems.
\emph{Proc. IEEE, 63}(9).\\[5pt]
Sandhu, R. S., Coyne, E. J., Feinstein, H. L., \& Youman, C. E. (1996). Role-based access control
models. \emph{Computer, 29}(2).
\end{column}
\begin{column}{0.485\textwidth}
\raggedright
Schneier, B., \& Kelsey, J. (1999). Secure audit logs to support computer forensics. \emph{ACM
TISSEC, 2}(2).\\[5pt]
Simon, R. T., \& Zurko, M. E. (1997). Separation of duty in role-based environments. \emph{CSFW}.\\[5pt]
Temoshok, D., et al. (2025). \emph{Digital identity guidelines: Authentication and authenticator
management} (NIST SP 800-63B-4).\\[5pt]
Ley N.º 8968 (2011) y su reglamento, Decreto N.º 37554-JP. Ley N.º 7975 (2000). Ley N.º 8454
(2005). Ley N.º 9048 (2012).\\[5pt]
ISO/IEC 17029:2019. ISO 14065:2020. ISO/IEC 27001:2022.\\[5pt]
Vargas Rivera, A. (2026). \emph{Material de estudio del curso CE-1115 y enunciado del proyecto,
parte 1}. ITCR.
\end{column}
\end{columns}
}
\note{La lista completa, en formato APA, está en cada uno de los cinco documentos.}
\end{frame}

\end{document}
"
%=====================================================================
\documentclass[aspectratio=169,t]{beamer}

\usepackage[utf8]{inputenc}
\usepackage[T1]{fontenc}
\usepackage[spanish,es-noshorthands,es-tabla]{babel}
\usepackage[default]{lato}
\usepackage{booktabs,array,colortbl}
\usepackage{tikz}
\usetikzlibrary{positioning,arrows.meta,calc}

\ifdefined\connotas\setbeameroption{show notes}\fi

%--------------------------- lamina y margenes -----------------------
\geometry{papersize={13.333in,7.5in}}
\setbeamersize{text margin left=0.6in,text margin right=0.6in}

%--------------------------- color -----------------------------------
\definecolor{acento}{HTML}{0F6E6A}
\definecolor{acentoclaro}{HTML}{E7F2F1}
\definecolor{tinta}{HTML}{1A1A1A}
\definecolor{gris}{HTML}{5F5E5A}
\definecolor{regla}{HTML}{B9B8B0}
\definecolor{plano}{HTML}{FAFAFA}
\definecolor{negado}{HTML}{F2F1EC}

%--------------------------- escala: cuatro tamanos ------------------
\newcommand{\tUno}{\fontsize{40}{45}\selectfont}
\newcommand{\tDos}{\fontsize{28}{35}\selectfont}
\newcommand{\tTres}{\fontsize{18}{25}\selectfont}
\newcommand{\tCuatro}{\fontsize{13}{17.5}\selectfont}
\makeatletter
\renewcommand\normalsize{\@setfontsize\normalsize{18}{25}}
\renewcommand\small{\@setfontsize\small{13}{17.5}}
\renewcommand\footnotesize{\@setfontsize\footnotesize{13}{17.5}}
\renewcommand\scriptsize{\@setfontsize\scriptsize{13}{17.5}}
\renewcommand\tiny{\@setfontsize\tiny{13}{17.5}}
\renewcommand\large{\@setfontsize\large{28}{35}}
\renewcommand\Large{\@setfontsize\Large{40}{45}}
\makeatother

%--------------------------- cromo -----------------------------------
\usetheme{default}
\setbeamertemplate{navigation symbols}{}
\setbeamercolor{structure}{fg=acento}
\setbeamercolor{normal text}{fg=tinta}
\setbeamercolor{background canvas}{bg=white}
\setbeamerfont{note page}{size=\fontsize{20}{27}\selectfont}

\setbeamertemplate{frametitle}{%
  \vskip0.0in%
  \begin{minipage}[t]{\textwidth}
    \raggedright\color{tinta}\tUno\bfseries\insertframetitle\par
  \end{minipage}\par
  \vskip0.12in}

\gdef\presentador{}
\newcommand{\quien}[1]{\gdef\presentador{#1}}
\setbeamertemplate{footline}{%
  \hbox to\paperwidth{\hskip0.6in{\tCuatro\color{gris}\insertframenumber}\hfill
    {\tCuatro\color{gris}\presentador}\hskip0.6in}%
  \vskip0.45in}

%--------------------------- atajos ----------------------------------
\newcommand{\dato}[1]{\textcolor{acento}{\textbf{#1}}}
\newcommand{\pie}[1]{\leavevmode{\tCuatro\color{gris}#1}}
\newcommand{\cabt}[1]{\leavevmode{\tCuatro\bfseries\color{gris}#1}}
\newcolumntype{L}[1]{>{\raggedright\arraybackslash}p{#1}}
\arrayrulecolor{regla}
\setlength{\arrayrulewidth}{0.8pt}
\renewcommand{\arraystretch}{1.5}
% celdas de la matriz
\newcommand{\vb}[2]{\mbox{#1\,{\color{acento}\bfseries #2}}}
\newcommand{\sinacc}{\cellcolor{negado}\leavevmode{\color{gris}---}}
\newcommand{\suj}[1]{\leavevmode{\tTres\bfseries #1}}

%=====================================================================
\begin{document}

%--------------------------- 1. PORTADA ------------------------------
\begin{frame}[plain]
\vspace*{1.25in}
{\tCuatro\color{gris}PROYECTO \ $\cdot$ \ PARTE 1 \ $\cdot$ \ DEFENSA ORAL}\\[0.3in]
{\tUno\bfseries\color{tinta}Diseño del control de acceso\\[4pt] para Verdeazul Carbono}\\[0.25in]
{\color{acento}\rule{1.6in}{4pt}}\\[0.25in]
{\tDos\color{gris}Políticas, matriz y arquitectura de seguridad}\\[0.7in]
{\tTres\bfseries Randall Bolaños \quad Javier Tenorio \quad Julio Varela \quad Axel Flores}\\[0.12in]
{\tCuatro\color{gris}Grupo 6 \ $\cdot$ \ CE-1115 Seguridad de la Información \ $\cdot$ \ Prof. MSc. Andrés
Vargas Rivera \ $\cdot$ \ ITCR \ $\cdot$ \ 15 de octubre de 2026}
\note{Reparto propuesto. Randall: cliente, activos y separación de responsabilidades. Javier:
amenazas, marco legal, fricción, riesgos residuales y respuestas del cliente. Julio: modelo de control
de acceso y matriz. Axel: controles de historial y arquitectura. Todos deben poder responder por cualquier lámina:
las preguntas del profesor son individuales.}
\end{frame}

%--------------------------- 2. EL NEGOCIO ---------------------------
\begin{frame}{Verdeazul vende confianza:\\ cada certificado respalda una tonelada real}
\quien{Randall}
\vspace{0.15in}
\begin{tikzpicture}[x=1in,y=1in,
  caja/.style={draw=regla,line width=0.8pt,fill=plano,minimum width=2.0in,minimum height=0.95in,
               align=center,font=\tTres,inner sep=2pt},
  fl/.style={-{Stealth[length=9pt,width=8pt]},draw=gris,line width=1.2pt}]
\node[caja] (a) at (1.00,0) {Finca\\forestal};
\node[caja] (b) at (3.53,0) {Técnico\\mide};
\node[caja] (c) at (6.06,0) {Verificador\\aprueba};
\node[caja,draw=acento,line width=1.6pt] (d) at (8.59,0) {\textbf{Certificado}\\de carbono};
\node[caja] (e) at (11.12,0) {Mercado\\internacional};
\foreach \p/\q in {a/b,b/c,c/d,d/e} \draw[fl] (\p) -- (\q);
\end{tikzpicture}

\vspace{0.4in}
{\tDos \dato{Un dato inflado es fraude.}}

\vspace{0.12in}
Auditores europeos pueden pedir la trazabilidad de cualquier certificado de los últimos diez años.
\note{Qué decir: Verdeazul certifica captura de carbono en fincas forestales. El técnico mide, el
verificador aprueba y se emite un certificado que el dueño vende afuera. El valor del certificado es la
confianza en el dato. Hoy todo vive en hojas de cálculo compartidas.

Si preguntan por qué empezamos por el negocio: porque el enunciado pide políticas que describan a este
cliente, y casi todas las decisiones salen de esta lámina.}
\end{frame}

%--------------------------- 3. LOS DOS SUSTOS -----------------------
\begin{frame}{Los dos sustos de la carta son fallas de integridad}
\quien{Randall}
\vspace{0.1in}
\begin{tabular}{@{}L{2.2in} L{9.5in}@{}}
\toprule
\cabt{PRIMER SUSTO} & Una medición cambió después de aprobada, y nadie sabe quién fue. \\
\midrule
\cabt{SEGUNDO SUSTO} & Un certificado se emitió dos veces sobre la misma hectárea. \\
\bottomrule
\end{tabular}

\vspace{0.45in}
{\tDos Orden de prioridad: \dato{integridad}, confidencialidad, disponibilidad.}

\vspace{0.12in}
Cuando dos pilares chocan, gana la integridad: sin registro en bitácora no hay acción.
\note{Qué decir: ninguno de los dos incidentes fue un robo de información. En los dos, el dato dejó de
ser confiable. Por eso la integridad va primero. La confidencialidad va segunda por la competencia
entre dueños y por los datos personales. La disponibilidad va tercera, con un matiz: importa más que el
dato exista en diez años que un día sin sistema.

Si preguntan por un ejemplo de «gana la integridad»: si la bitácora no se puede escribir, la operación
se rechaza aunque eso detenga el trabajo (POL-BT-06).}
\end{frame}

%--------------------------- 4. ACTIVOS ------------------------------
\begin{frame}{Lo que más vale es lo que sostiene un certificado}
\quien{Randall}
\vspace{0.05in}
\begin{tabular}{@{}L{4.6in} L{2.6in} L{4.4in}@{}}
\toprule
\cabt{ACTIVO} & \cabt{CLASIFICACIÓN} & \cabt{LO QUE NO PUEDE FALLAR} \\
\midrule
Mediciones y certificados & Confidencial & \dato{Integridad} \\
Catálogo de fincas: la cartera & Confidencial & Integridad y confidencialidad \\
Datos personales de los dueños & Confidencial & Confidencialidad \\
Bitácora y llave de sello & Restringido & \dato{Integridad} \\
\bottomrule
\end{tabular}

\vspace{0.3in}
\pie{El inventario completo tiene once activos, cada uno valorado en los tres pilares (informe de políticas, sección 2).}
\note{Qué decir: clasificamos en cuatro niveles (público, interno, confidencial, restringido) y valoramos
cada activo con la escala de FIPS 199. La Gerencia de Certificación es la responsable del dato; el
administrador y el operador son custodios.

Si preguntan por un activo que no es obvio: los parámetros de cálculo (A-10). Un factor alterado infla
todos los certificados a la vez. Y la llave de sello (A-09): quien la tenga fabrica certificados
creíbles.}
\end{frame}

%--------------------------- 5. AMENAZAS -----------------------------
\begin{frame}{\textcolor{acento}{7 de 11} amenazas son internas:\\ el atacante más probable ya tiene una cuenta}
\quien{Javier}
\vspace{0.05in}
\begin{tabular}{@{}L{1.4in} L{6.6in} L{3.6in}@{}}
\toprule
\cabt{ID} & \cabt{YA OCURRIÓ EN VERDEAZUL} & \cabt{PILAR} \\
\midrule
AM-03 & Una medición editada después de aprobarse & Integridad \\
AM-06 & Un certificado emitido dos veces & Integridad \\
AM-04 & Cientos de fincas consultadas en un fin de semana & Confidencialidad \\
\bottomrule
\end{tabular}

\vspace{0.3in}
El diseño no mira el perímetro: mira qué puede hacer cada persona autorizada.
\note{Qué decir: el modelo tiene once amenazas, cada una con agente, motivo, activo, la debilidad que
hoy lo permite y el riesgo. Siete son internas: técnico que infla, misma persona que mide y aprueba,
edición silenciosa, extracción de la cartera, administrador u operador que abusa, doble emisión por
error y, desde la consulta al cliente, el empleado con un vínculo con un dueño (AM-11): no existe un
registro de esos vínculos. Tres ya ocurrieron.

Externas: dueño que quiere ver a un competidor, robo de credenciales, ransomware y certificado
falsificado o cuenta de auditor tomada.

Si preguntan cómo calculamos el riesgo: probabilidad por impacto. Probabilidad alta es que ya ocurrió
aquí o que es la vía habitual de ataque en el país.}
\end{frame}

%--------------------------- 6. MARCO LEGAL --------------------------
\begin{frame}{Cada obligación legal termina\\ en una regla que se puede comprobar}
\quien{Javier}
\vspace{0.05in}
\begin{tabular}{@{}L{2.9in} L{4.2in} L{4.5in}@{}}
\toprule
\cabt{NORMA} & \cabt{EXIGE} & \cabt{REGLA EN VERDEAZUL} \\
\midrule
Ley 8968, art. 14 & Autorización para transferir datos & El auditor consulta en la plataforma \\
Reglamento, art. 38 & Avisar una fuga en 5 días hábiles & Procedimiento con ese plazo \\
Ley 7975, art. 2 & Medidas razonables de secreto & Acceso por necesidad y registro \\
ISO/IEC 17029 & Quien revisa no participó & \dato{Medir y aprobar, separados} \\
\bottomrule
\end{tabular}

\vspace{0.2in}
\pie{El cliente aclaró que no audita bajo una norma determinada: la ISO se cita como referencia. Su exigencia firme es la trazabilidad de diez años.}
\note{Qué decir: la Ley 8968 aplica porque los datos de los dueños se comparten con auditores externos;
no es una base de uso exclusivamente interno. Derechos del titular en cinco días hábiles; consentimiento
informado al dar de alta a un dueño; conservación hasta diez años y después se desasocia.

La Ley 7975 es la que permite reclamar si la cartera llega a la competencia, pero solo si la empresa
puede demostrar que tomó medidas razonables: por eso el acceso por asignación y la bitácora tienen
además valor legal.

También: Código Penal (arts. 196 bis, 217 bis, 229 bis y 231), Código de Trabajo art. 71 g) y Ley 8454
de firma digital.

Auditores, según la minuta (P-18): una vez al año, varios días, en sitio o a distancia. Pueden exportar
evidencia, pero como paquetes que genera el sistema, sellados, registrados y sin datos personales. El
acceso vence con la ventana autorizada; prorrogarlo es una autorización nueva.}
\end{frame}

%--------------------------- 7. MODELO: TRES CAPAS -------------------
\begin{frame}{Roles con restricciones,\\ y condiciones que solo pueden quitar}
\quien{Julio}
\vspace{0.2in}
\begin{tikzpicture}[x=1in,y=1in,
  capa/.style={draw=acento,line width=1.4pt,fill=acentoclaro,minimum width=2.4in,minimum height=1.1in,
               align=center,inner sep=3pt},
  ext/.style={draw=regla,line width=0.8pt,fill=plano,minimum height=1.1in,align=center,inner sep=3pt},
  fl/.style={-{Stealth[length=9pt,width=8pt]},draw=gris,line width=1.2pt},
  no/.style={-{Stealth[length=9pt,width=8pt]},draw=gris,line width=1.0pt,densely dashed}]
\node[ext,minimum width=1.6in] (p) at (0.8,0) {Petición\\autenticada};
\node[capa] (r) at (3.25,0) {\textbf{1 $\cdot$ Roles}\\{\tCuatro ¿Este rol puede?}};
\node[capa] (c) at (6.10,0) {\textbf{2 $\cdot$ Condiciones}\\{\tCuatro ¿Aquí y ahora?}};
\node[capa] (h) at (8.95,0) {\textbf{3 $\cdot$ Historial}\\{\tCuatro ¿Encaja con lo previo?}};
\node[ext,minimum width=1.5in] (d) at (11.35,0) {Conceder};
\draw[fl] (p) -- (r);
\draw[fl] (r) -- node[above,font=\tCuatro,text=gris]{sí} (c);
\draw[fl] (c) -- node[above,font=\tCuatro,text=gris]{sí} (h);
\draw[fl] (h) -- node[above,font=\tCuatro,text=gris]{sí} (d);
\node[draw=regla,line width=0.8pt,fill=negado,minimum width=8.1in,minimum height=0.6in] (n) at (6.10,-1.55)
  {Denegar y registrar en bitácora};
\foreach \k in {r,c,h} \draw[no] (\k.south) -- node[right,font=\tCuatro,text=gris]{no} (\k.south |- n.north);
\end{tikzpicture}

\vspace{0.3in}
\dato{Solo la primera capa concede.} La matriz muestra lo máximo que cada rol puede llegar a hacer.
\note{Qué decir: elegimos RBAC con restricciones, nivel RBAC2 de Sandhu. Sobre esa base, un conjunto
cerrado de atributos decide si el permiso aplica a esta finca, en este estado y en este momento.
Es el enfoque centrado en roles de Kuhn, Coyne y Weil.

La frontera del híbrido: la capa 1 dice si el rol tiene el verbo; la capa 2 si aplica a este objeto ahora;
la capa 3 si encaja con el historial. Se combinan sin contradecirse porque solo la tabla de roles concede
y las otras dos son filtros: una negativa en cualquier capa gana.

Si preguntan por qué RBAC: la carta describe funciones, no personas; la gente rota; y hay que poder
explicarlo a un auditor leyendo una fila.}
\end{frame}

%--------------------------- 8. DESCARTES ----------------------------
\begin{frame}{Los otros tres modelos fallan\\ por razones de este caso}
\quien{Julio}
\vspace{0.05in}
\begin{tabular}{@{}L{2.2in} L{9.5in}@{}}
\toprule
\cabt{MODELO} & \cabt{POR QUÉ NO EN VERDEAZUL} \\
\midrule
DAC & Es lo que hay hoy: las hojas compartidas donde ocurrieron los dos sustos. \\
MAC & No hay niveles de secreto. Haría falta un compartimento por finca. \\
ABAC puro & Desaparece la fila que responde «¿qué puede hacer un verificador?». \\
\bottomrule
\end{tabular}

\vspace{0.3in}
\pie{De cada uno se conserva algo: la titularidad como atributo (DAC), la política que el usuario no puede
alterar (MAC) y las condiciones acotadas (ABAC).}
\note{DAC: el técnico que crea la medición sería su dueño y podría seguir modificándola o compartirla.
Es el primer susto.

MAC: Bell-LaPadula ordena por niveles de secreto y protege la confidencialidad. Aquí el riesgo principal
es de integridad. La variante de integridad, Biba, pone una etiqueta fija, pero la confianza en una
medición cambia con su estado y depende de quién hizo el paso anterior. Es el argumento de Clark y
Wilson. MAC sí lo usamos en la plataforma: SELinux confina los contenedores.

ABAC puro: podría expresarlo todo, pero no se puede revisar leyendo una fila, y la empresa no tiene
quién gobierne un catálogo de atributos.}
\end{frame}

%--------------------------- 9. NIVEL Y EXPLOSION --------------------
\begin{frame}{\textcolor{acento}{Siete} roles hoy, y siete con el doble de fincas}
\quien{Julio}
\vspace{0.25in}
{\tDos El alcance nunca es un rol: es un atributo.}

\vspace{0.35in}
\begin{tabular}{@{}L{3.3in} L{8.4in}@{}}
\toprule
RBAC$_2$, sin jerarquía & Si el verificador heredara del técnico, podría registrar: se rompe la separación. \\
Explosión de roles & Aparece al usar un rol para decir dónde o cuándo. Aquí eso vive en las condiciones. \\
La consulta lo probó & Cinco necesidades nuevas: \dato{un rol y cuatro atributos}. \\
\bottomrule
\end{tabular}
\note{Los siete roles: técnico, verificador, supervisor de campo, gerencia de certificación,
administrador de catálogo y cuentas, auditor externo y dueño de finca. Más una identidad de servicio
(SIS), el operador de plataforma, que no tiene rol en la aplicación, y el público sin autenticar.

El séptimo rol salió de la consulta: el cliente dijo que las fincas las asignan los supervisores del
personal de campo (minuta, P-05). Las otras cuatro necesidades nuevas (el verificador que mide por
excepción, el técnico que apoya a un compañero, la sociedad con un representante y el empleado con un
vínculo con un dueño) no son funciones de trabajo: se resolvieron con atributos.

Por qué RBAC2 y no RBAC3: no hay cadena de mando entre las siete funciones; la herencia vuelve
implícitos los permisos y la matriz no debe tener nada implícito; y con siete roles no ahorra nada.

Restricciones: un rol por persona, una identidad por persona, al menos dos verificadores, cuenta de
auditor con ventana. La pareja técnico--verificador es la única excepción: es dinámica. Un verificador
puede recibir una habilitación temporal de técnico, con fecha de fin, y nunca aprueba lo que midió
(minuta, P-01).

Si preguntan dónde está el riesgo de explosión: un rol por finca para aislar a los dueños, o roles como
«Técnico Zona Sur 2027».}
\end{frame}

%--------------------------- 10. HBAC --------------------------------
\begin{frame}{Cuatro controles de historial, uno de cada tipo}
\quien{Axel}
\vspace{0.0in}
\begin{tabular}{@{}L{0.7in} L{1.7in} L{4.6in} L{4.4in}@{}}
\toprule
 & \cabt{TIPO} & \cabt{QUÉ LO DISPARA} & \cabt{QUÉ DECIDE} \\
\midrule
H1 & Frecuencia & Muchas fincas distintas en un día & \dato{Tres niveles}: avisa, confirma, restringe \\
H2 & Tiempo & Una decisión fuera de jornada & Pide segundo factor y motivo \\
H3 & Combinada & Fuera de jornada, muchas fincas, casi todas nuevas & Confirma y avisa; si insiste, restringe \\
H4 & Secuencia & Quien participó antes quiere decidir & \dato{Deniega} \\
\bottomrule
\end{tabular}

\vspace{0.2in}
\pie{Ninguna señal aislada restringe: el cliente pidió avisar primero y no bloquear de inmediato (minuta, P-10).}
\note{Son reglas explícitas, no aprendizaje automático: se le pueden leer a un auditor y se pueden probar.

Cómo se distingue lo legítimo: se cuentan señales: volumen, novedad, momento e insistencia. El trabajo
extraordinario legítimo muestra una; la extracción de la cartera las muestra juntas.

Tres niveles: aviso (el usuario no nota nada), verificación adicional (segundo factor y motivo) y
restricción temporal (no abre fincas nuevas hasta el siguiente día hábil; la cuenta sigue activa).

H1 mitiga AM-04. Umbrales por rol, calibrados con datos reales; iniciales 40, 60 y 80 fincas por día
(supuesto P-13). Los intentos denegados cuentan doble.

H2 mitiga AM-08, la cuenta robada usada de madrugada. H4 es la separación de responsabilidades como
control de historial: reglas C3, C8 y C10.

H3 es la firma del incidente: fin de semana, cientos de fincas, ninguna suya. Pide confirmación en la
finca once y avisa; restringe en la veintiuna.

Falsos positivos: verificador con la cola acumulada al cierre, o nuevo y sin historial. El primer
período corre en observación para calibrar.

Descartada: Muralla China. Todos los dueños compiten: cada empleado quedaría con una sola finca.}
\end{frame}

%--------------------------- 11. EL TECNICO --------------------------
\begin{frame}{El técnico del fin de semana\\ hoy no vería ni una finca ajena}
\quien{Axel}
\vspace{0.05in}
\begin{tabular}{@{}L{4.6in} L{7.1in}@{}}
\toprule
\cabt{LO QUE HACE} & \cabt{LO QUE HACE EL SISTEMA} \\
\midrule
Abre una finca que no tiene asignada & \dato{Deniega}: su rol no ve fincas ajenas \\
Insiste tres, cinco, diez veces & Avisa a su supervisor; pide segundo factor; lo detiene media hora \\
La necesitaba de verdad & El supervisor le da una asignación temporal, que vence sola \\
Presenta su renuncia & Reporte de sus consultas de 90 días \\
\bottomrule
\end{tabular}

\vspace{0.25in}
\pie{El cliente lo confirmó: fincas asignadas con flexibilidad, y respuesta según la gravedad (minuta, P-10 y P-11).}
\note{Qué decir: la carta dice «no sabemos qué hacer con eso, pero nos gustaría que el sistema al menos se
diera cuenta». Nuestra primera respuesta no es vigilar el permiso: es quitarlo. Un técnico solo ve las
fincas que tiene asignadas (regla C1).

El cliente lo confirmó, con una condición: flexibilidad para apoyar a un compañero, sustituir a un
ausente o atender una visita extraordinaria, sin acceso permanente a todas las fincas. Se resuelve con
asignaciones temporales que crea el supervisor de campo, con motivo y fecha de fin (regla C19).

El diseño original suspendía la sesión a la quinta denegación. El cliente pidió no bloquear de
inmediato, y se cambió: aviso al supervisor a la tercera, segundo factor a la quinta, media hora de
restricción a la décima.

Los controles de historial quedan para quienes sí necesitan leer muchas fincas: verificador, supervisor,
gerencia, administrador y auditor. Si el técnico hubiera sido verificador, H3 le pide confirmación en la
finca once y lo restringe en la veintiuna.

Además hay una vía legal: confidencialidad firmada, que persiste después de la salida (Ley 8968 art. 11
y Ley 7975).}
\end{frame}

%--------------------------- 12. NADIE EMITE -------------------------
\begin{frame}{Nadie emite un certificado: lo hace el sistema}
\quien{Randall}
\vspace{0.25in}
\begin{tikzpicture}[x=1in,y=1in,
  paso/.style={draw=acento,line width=1.4pt,fill=acentoclaro,minimum width=3.3in,minimum height=1.1in,
               align=center,inner sep=3pt},
  sis/.style={draw=gris,line width=1.2pt,fill=negado,minimum width=3.3in,minimum height=1.1in,
              align=center,inner sep=3pt},
  fl/.style={-{Stealth[length=9pt,width=8pt]},draw=gris,line width=1.2pt},
  nota/.style={font=\tCuatro,text=gris,align=center}]
\node[paso] (a) at (1.65,0) {\textbf{Técnico}\\registra y envía};
\node[paso] (b) at (6.05,0) {\textbf{Verificador}\\aprueba o devuelve};
\node[sis]  (c) at (10.45,0) {\textbf{Sistema}\\emite};
\draw[fl] (a) -- (b); \draw[fl] (b) -- (c);
\node[nota,below=0.08in of a] {solo en fincas asignadas\\y con el período abierto};
\node[nota,below=0.08in of b] {no registró ninguna versión\\de esa medición};
\node[nota,below=0.08in of c] {no hay otro certificado vigente\\para la parcela y el período};
\end{tikzpicture}

\vspace{0.35in}
Ninguna persona tiene el permiso de emitir. \dato{Hacen falta dos, y distintas.}
\note{La acción crítica: emitir un certificado. Si una sola persona pudiera medir y aprobar, podría
inflar la biomasa, aprobar su propio dato y obtener un certificado por toneladas que no existen.

La segunda acción crítica: corregir una medición ya congelada. Es el primer susto, y se trata igual,
con el orden que pidió el cliente (minuta, P-07): solicita el técnico responsable, un verificador o la
Gerencia; corrige el técnico; revisa un verificador que no es autor ni solicitante; y la Gerencia da la
autorización final. Son al menos tres personas.

Si preguntan qué pasa con el certificado tras una enmienda (minuta, P-08): si las toneladas no cambian,
sigue vigente con la enmienda anotada. Si cambian, el sistema lo revoca y emite el reemplazo en un solo
acto, enlazados: nunca hay dos vigentes y ninguno se borra. Si ya se vendió, la Gerencia lo registra y
coordina con las partes.}
\end{frame}

%--------------------------- 13. POR MEDICION ------------------------
\begin{frame}{La separación se comprueba\\ por medición, no por puesto}
\quien{Randall}
\vspace{0.05in}
\begin{tabular}{@{}L{3.4in} L{8.3in}@{}}
\toprule
\cabt{SI SE COMPRUEBA\ldots} & \cabt{UN TÉCNICO ASCENDIDO A VERIFICADOR\ldots} \\
\midrule
por el puesto de hoy & aprueba lo que él mismo midió el año anterior. \\
\dato{por la medición} & \dato{no puede: figura como autor.} \\
\bottomrule
\end{tabular}

\vspace{0.35in}
Técnica: \textbf{separación por etapas del flujo}, reforzada con roles mutuamente excluyentes.
\note{Qué decir: quien inicia el proceso no puede cerrarlo. El sistema no pregunta qué rol tiene la
persona, pregunta si participó en esta medición. En la literatura es la separación basada en el objeto
o en el historial (Nash y Poland 1990; Simon y Zurko 1997).

Por qué es la adecuada: apunta al riesgo tal como lo dijo el cliente; sobrevive a los cambios de puesto
(la lección de Kerviel); no depende de que los roles estén bien administrados; y deja dos nombres
distintos en cada certificado.

El cliente confirmó que algunos verificadores fueron técnicos y que por excepción pueden medir, y pidió
«garantizar una revisión independiente, sin impedir que el personal pueda asumir otras
responsabilidades» (minuta, P-01). Una separación por puesto obligaría a escoger; esta da las dos
cosas: el verificador mide con una habilitación temporal y no podrá aprobar nunca lo que midió.

Coincide con el principio de imparcialidad de ISO/IEC 17029, que citamos solo como referencia.}
\end{frame}

%--------------------------- 14. ALTERNATIVAS ------------------------
\begin{frame}{Las alternativas se descartaron\\ por lo que cuestan aquí}
\quien{Randall}
\vspace{0.0in}
\begin{tabular}{@{}L{4.3in} L{7.4in}@{}}
\toprule
\cabt{TÉCNICA} & \cabt{POR QUÉ NO ES LA PRINCIPAL} \\
\midrule
Roles excluyentes, solos & No cubren a quien cambia de puesto \\
Doble control (cuatro ojos) & \dato{Duplica la verificación de cada certificado} \\
Segregación temporal & Limita cuándo, no quién \\
Rotación & No impide un acto; se recomienda contra la colusión \\
\bottomrule
\end{tabular}

\vspace{0.2in}
\pie{La segregación temporal sí se usa donde el problema es la duración: la asignación temporal, la habilitación del verificador, el plazo de la enmienda y la cuenta del auditor.}
\note{Roles mutuamente excluyentes estáticos (RBAC2): se conservan como refuerzo. Li, Tripunitara y Bizri
muestran que excluir roles es un mecanismo para aproximar la separación, no la separación misma.

Exclusión dinámica: basta cerrar sesión y abrir otra. Sola, demora el fraude y no lo impide. Se usa
como refuerzo en la pareja técnico--verificador, detrás de la regla C3.

Doble control: era la opción más fuerte contra el fraude. Se descartó por costo operativo: duplica la
verificación de cada certificado, con cualquier tamaño de equipo.

Rotación: cuesta conocimiento de la finca y depende de cuántas personas haya, un dato que el cliente
pidió no fijar (minuta, P-02).}
\end{frame}

%--------------------------- 15. FRICCION ----------------------------
\begin{frame}{La fricción existe y está acotada}
\quien{Javier}
\vspace{0.0in}
\begin{tabular}{@{}L{5.2in} L{6.5in}@{}}
\toprule
\cabt{LO QUE CUESTA} & \cabt{CÓMO SE MANEJA} \\
\midrule
Siempre hacen falta dos personas & Mínimo dos verificadores; \dato{sin atajo de emergencia} \\
Las devoluciones quitan tiempo de campo & Validación de rangos al registrar \\
Una enmienda pide tres personas & Plazo de diez días; el sistema la clasifica \\
La Gerencia concentra las segundas firmas & Nunca inicia lo que aprueba \\
\bottomrule
\end{tabular}
\note{Qué decir: no existe un atajo para aprobar lo propio. Un atajo así sería el hueco que el auditor
buscaría. Si el único verificador presente participó, la medición espera; la cola avisa a la Gerencia.

Lo que la técnica no resuelve: la colusión. Si un técnico y un verificador se ponen de acuerdo, el flujo
se cumple. La separación sube el costo del fraude de una persona a dos. Quedan la rotación, un indicador
de parejas repetidas y la auditoría externa.

La Gerencia es un cuello de botella y un punto de confianza: firma en varios flujos, pero nunca inicia
lo que aprueba, no registra ni aprueba mediciones, ya no asigna fincas y no resuelve sus propias alertas.

Enmienda: revisa un verificador distinto del original solo cuando es material (cambia toneladas o ya hay
certificado); lo determina el sistema. Con solo dos verificadores, si uno solicita sobre lo que aprobó
el otro no queda revisor: la presenta el técnico o la Gerencia.

La flexibilidad también cuesta un trámite: la asignación temporal la crea el supervisor y la
habilitación de un verificador para medir pasa por el administrador y la Gerencia.}
\end{frame}

%--------------------------- 16. MATRIZ (1 de 2) ---------------------
\begin{frame}{Lo que no está escrito está denegado}
\quien{Julio}
\vspace{-0.02in}
{\tCuatro
\renewcommand{\arraystretch}{1.3}
\setlength{\tabcolsep}{5pt}
\begin{tabular}{@{}L{0.75in}|L{1.95in}|L{2.25in}|L{2.0in}|L{1.85in}|L{2.3in}@{}}
\rowcolor{acentoclaro}
 & \textbf{Fincas y parcelas} & \textbf{Datos personales del dueño} & \textbf{Períodos de medición} & \textbf{Asignaciones} & \textbf{Mediciones} \\
\hline
\suj{TEC} & \vb{leer}{C1} & \vb{leer}{C1} (solo contacto) & leer & \vb{leer}{C1} (las propias) & \vb{leer}{C1}, \vb{crear}{C2}, \vb{modificar}{C2}, \vb{enviar}{C2} \\
\hline
\suj{VER} & \vb{leer}{H} & \sinacc & leer & \sinacc & \vb{leer}{H}, \vb{aprobar}{C3}, \vb{devolver}{C3} \\
\hline
\suj{SUP} & \vb{leer}{H} & \sinacc & leer & leer, \vb{crear}{C19}, \vb{modificar}{C19} & \vb{leer}{H} (solo estado y fechas) \\
\hline
\suj{GER} & \vb{leer}{H}, \vb{aprobar}{C11} & \vb{leer}{H} & leer, \vb{abrir}{C14}, \vb{cerrar}{C14} & leer & \vb{leer}{H} \\
\hline
\suj{ADM} & \vb{leer}{H}, \vb{proponer}{C11} & \vb{leer}{H}, crear, modificar, \vb{eliminar}{C16} & leer, crear, \vb{proponer}{C14} & \sinacc & \sinacc \\
\hline
\suj{AUD} & \vb{leer}{C6}, \vb{exportar}{C18} & \vb{leer}{C6} & \vb{leer}{C6}, \vb{exportar}{C18} & \vb{leer}{C6}, \vb{exportar}{C18} & \vb{leer}{C6}, \vb{exportar}{C18} \\
\hline
\suj{DUE} & \vb{leer}{C5} & \vb{leer}{C5} & leer & \vb{leer}{C5} & \vb{leer}{C5} (solo aprobadas) \\
\hline
\suj{SIS} & leer & leer (solo el nombre del titular) & leer & leer & leer \\
\hline
\suj{OPE} & \sinacc & \sinacc & \sinacc & \sinacc & \sinacc \\
\hline
\suj{PUB} & \sinacc & \sinacc & \sinacc & \sinacc & \sinacc \\
\end{tabular}

\vspace{0.06in}
{\color{gris}Matriz real, objetos 1 a 5 de 10. \quad --- sin acceso \quad {\color{acento}\textbf{C}$n$} permiso sujeto a la
regla $n$ \quad {\color{acento}\textbf{H}} lectura vigilada por historial}
}
\note{Cómo leerla: sujetos en las filas, objetos en las columnas. Una celda gris es sin acceso. Una marca C
significa que el permiso solo vale si se cumple esa regla; las condiciones solo pueden quitar.

Recorrido sugerido: fila TEC, columna Mediciones: crea, modifica y envía, pero no aprueba. Fila VER:
aprueba y devuelve con C3, no crea ni modifica, y no ve los datos personales del dueño. Fila SUP, que
es nueva: asigna fincas (C19) y de las mediciones solo ve estado y fechas. Fila GER: abre y cierra
períodos que registró el administrador (C14); ya no crea asignaciones. Fila AUD: leer y exportar
evidencia, con ventana (C6 y C18), y ningún verbo de escritura. Fila DUE: solo lo suyo (C5). Filas OPE
y PUB: nada en esta mitad.

Diez sujetos: siete roles, el sistema, el operador y el público sin autenticar.

Mínimo privilegio: cada permiso amplio tiene su justificación en la sección 6 del documento de la matriz.}
\end{frame}

%--------------------------- 17. MATRIZ (2 de 2) ---------------------
\begin{frame}{Nadie tiene «eliminar», y nadie emite}
\quien{Julio}
\vspace{-0.02in}
{\tCuatro
\renewcommand{\arraystretch}{1.3}
\setlength{\tabcolsep}{5pt}
\begin{tabular}{@{}L{0.75in}|L{2.3in}|L{1.9in}|L{2.75in}|L{1.4in}|L{2.0in}@{}}
\rowcolor{acentoclaro}
 & \textbf{Enmiendas} & \textbf{Certificados} & \textbf{Cuentas y roles} & \textbf{Bitácora} & \textbf{Alertas de historial} \\
\hline
\suj{TEC} & \vb{leer}{C1}, \vb{solicitar}{C7}, \vb{corregir}{C9} & \sinacc & \vb{modificar}{C13} & \sinacc & \sinacc \\
\hline
\suj{VER} & leer, \vb{solicitar}{C7}, \vb{revisar}{C10} & \vb{leer}{H} & \vb{modificar}{C13} & \sinacc & \sinacc \\
\hline
\suj{SUP} & leer (solo estado) & \sinacc & \vb{modificar}{C13} & \sinacc & \vb{leer}{C15}, \vb{resolver}{C15} \\
\hline
\suj{GER} & leer, \vb{solicitar}{C7}, \vb{autorizar}{C8} & \vb{leer}{H} & leer, \vb{aprobar}{C12}, \vb{modificar}{C13} & leer & leer, \vb{resolver}{C15} \\
\hline
\suj{ADM} & \sinacc & \sinacc & leer, \vb{crear}{C12}, \vb{desactivar}{C12}, \vb{proponer}{C12}, \vb{modificar}{C13} & \sinacc & \sinacc \\
\hline
\suj{AUD} & \vb{leer}{C6}, \vb{exportar}{C18} & \vb{leer}{C6}, \vb{exportar}{C18} & \vb{leer}{C6}, \vb{modificar}{C13} & \vb{leer}{C6}, \vb{exportar}{C18} & \vb{leer}{C6}, \vb{exportar}{C18} \\
\hline
\suj{DUE} & \vb{leer}{C5} & \vb{leer}{C5} & \vb{modificar}{C13} & \sinacc & \sinacc \\
\hline
\suj{SIS} & leer & \dato{emitir C4}, \vb{revocar}{C4} & \sinacc & anexar & crear \\
\hline
\suj{OPE} & \sinacc & \sinacc & \sinacc & \sinacc & \sinacc \\
\hline
\suj{PUB} & \sinacc & \vb{comprobar}{C20} & \sinacc & \sinacc & \sinacc \\
\end{tabular}

\vspace{0.06in}
{\color{gris}Matriz real, objetos 6 a 10 de 10. \quad Toda decisión exige segundo factor reciente. \quad {\color{acento}\textbf{C17}}: nadie decide sobre una finca con cuyo dueño declaró un vínculo.}
}
\note{Dos cosas que se ven en esta mitad. Una: el verbo «eliminar» no aparece en enmiendas, certificados,
cuentas ni bitácora. Nada se borra; todo cambio es una versión nueva o un cambio de estado. Dos: «emitir»
solo está en la fila SIS. Ninguna fila de persona lo tiene.

Columna Enmiendas: solicitar (técnico, verificador o Gerencia), corregir (técnico), revisar
(verificador) y autorizar (Gerencia), con las reglas C7 a C10. Cuatro pasos, al menos tres personas.

Columna Cuentas y roles: C13 es «credenciales y declaración propias»; C12 es «nadie se administra a sí
mismo»: el administrador propone el rol y la Gerencia lo aprueba.

Columna Bitácora: solo se agrega. La leen la Gerencia y el auditor.

Fila PUB: quien tenga el código de un certificado comprueba si es auténtico y si sigue vigente, y nada
más (C20). Sale de la respuesta del cliente sobre autenticidad y estado del certificado (minuta, P-04).

C17, abstención por vínculo: quien declaró un vínculo familiar o comercial con un dueño no mide, no
aprueba, no asigna ni autoriza nada sobre esa finca (minuta, P-15).

Al lado de la matriz hay una tabla de exclusiones: cada persona tiene un solo rol, salvo la habilitación
temporal de un verificador como técnico.}
\end{frame}

%--------------------------- 18. TRES CELDAS -------------------------
\begin{frame}{Tres celdas bastan para leer la separación}
\quien{Julio}
\vspace{0.0in}
\begin{tabular}{@{}L{1.0in} L{5.9in} L{4.8in}@{}}
\toprule
 & \cabt{MEDICIONES} & \cabt{CERTIFICADOS} \\
\midrule
\textbf{TEC} & Crea, modifica y envía. Solo sus borradores, con el período abierto. & Sin acceso \\
\textbf{VER} & \dato{Aprueba solo si no registró ninguna versión.} & Lee \\
\textbf{SIS} & Lee & \dato{Emite solo si no hay otro vigente.} \\
\bottomrule
\end{tabular}

\vspace{0.25in}
\pie{Caso de prueba T-06: la verificadora que antes fue técnica intenta aprobar una medición que registró. Resultado esperado: deniega.}
\note{Son las reglas C2, C3 y C4. Cada regla de la matriz tiene un caso de prueba con resultado esperado:
treinta y siete casos, T-01 a T-37. El equipo que construya debe convertirlos en pruebas automáticas, más una por
cada celda sin acceso.

Otros casos para citar: T-07, segunda emisión sobre la misma parcela y período: no emite. T-08, el dueño
cambia el identificador en la dirección: deniega y responde «no encontrado». T-10, el auditor envía una
petición de modificación: deniega.

Caso nuevo para citar: T-17, después del cierre el verificador aprueba una medición que se envió a
tiempo: concede. El cliente pidió distinguir registrar, modificar, revisar y aprobar (minuta, P-06).

El mismo patrón, quien inicia no cierra, se repite en enmiendas (C7 a C10), períodos (C14),
asignaciones (C19), catálogo (C11) y cuentas (C12).}
\end{frame}

%--------------------------- 19. ARQUITECTURA ------------------------
\begin{frame}{Un solo punto decide, en el servidor,\\ antes de tocar el dato}
\quien{Axel}
\vspace{0.1in}
\begin{tikzpicture}[x=1in,y=1in,
  comp/.style={draw=regla,line width=0.8pt,fill=white,align=center,inner sep=3pt,minimum height=0.8in},
  pdp/.style={draw=acento,line width=1.6pt,fill=acentoclaro,align=center,inner sep=3pt,minimum height=0.8in},
  zona/.style={draw=gris,densely dashed,line width=0.7pt},
  tz/.style={font=\tCuatro\bfseries,text=gris,anchor=north west},
  dos/.style={{Stealth[length=8pt,width=7pt]}-{Stealth[length=8pt,width=7pt]},draw=gris,line width=1.1pt}]
\draw[zona] (0,0) rectangle (2.0,3.3);
\draw[zona] (2.35,0) rectangle (3.95,3.3);
\draw[zona] (4.3,0) rectangle (8.9,3.3);
\draw[zona] (9.25,0) rectangle (12.1,3.3);
\node[tz] at (0.08,3.25) {NO CONFIABLE};
\node[tz] at (2.43,3.25) {BORDE};
\node[tz] at (4.38,3.25) {APLICACIÓN};
\node[tz] at (9.33,3.25) {DATOS Y SECRETOS};
\node[comp,minimum width=1.6in] (nav) at (1.0,1.25) {Navegador};
\node[comp,minimum width=1.2in] (px) at (3.15,1.25) {Proxy};
\node[comp,minimum width=4.0in] (app) at (6.6,1.25) {\textbf{Aplicación}\\{\tCuatro identifica, valida y hace cumplir}};
\node[pdp,minimum width=4.0in] (opa) at (6.6,2.45) {\textbf{Motor de decisión}\\{\tCuatro la matriz, como política}};
\node[comp,minimum width=2.4in] (pg) at (10.68,2.1) {\textbf{Base de datos}\\{\tCuatro segundo candado}};
\node[comp,minimum width=2.4in] (bao) at (10.68,0.75) {\textbf{Bóveda}\\{\tCuatro secretos y sello}};
\draw[dos] (nav) -- (px);
\draw[dos] (px) -- (app);
\draw[dos,draw=acento] (app) -- (opa);
\draw[dos] ($(app.east)+(0,0.18)$) -- ++(0.3,0) |- (pg.west);
\draw[dos] ($(app.east)+(0,-0.18)$) -- ++(0.3,0) |- (bao.west);
\end{tikzpicture}

\vspace{0.2in}
\dato{Esconder un botón no es autorización.} Si el servidor no decide, cualquiera llama directo.
\note{Flujo de una petición: el proxy termina la conexión cifrada; la aplicación identifica la sesión,
valida la forma de los datos y lee de la base los hechos; entrega al motor el sujeto, la acción, el
objeto y esos hechos; el motor devuelve la decisión con la regla que la determinó. Solo entonces se toca
el dato, en una transacción que también escribe la bitácora. La base aplica sus propias restricciones.

Fronteras de confianza: nada de lo que llega del navegador se da por bueno; los hechos se toman de la
base, nunca de la petición; si el motor no responde, se deniega; la base no confía en la aplicación.

Si preguntan qué pasa si el control está solo en la interfaz: el técnico envía a mano la petición de
aprobar su propia medición y el servidor la ejecuta. Es el error de referencia directa a objetos de la
clase 13.

Hospedaje, que el cliente dejó a nuestro criterio (minuta, P-23): un servidor virtual contratado, operado
por personal que la empresa designe. Hay además una ruta pública de solo lectura para comprobar
certificados.}
\end{frame}

%--------------------------- 20. TECNOLOGIAS -------------------------
\begin{frame}{Seis tecnologías libres,\\ elegidas por lo que protegen}
\quien{Axel}
\vspace{-0.05in}
{\renewcommand{\arraystretch}{1.36}
\begin{tabular}{@{}L{2.9in} L{3.7in} L{5.1in}@{}}
\toprule
\cabt{DECISIÓN} & \cabt{ELECCIÓN} & \cabt{LO QUE APORTA} \\
\midrule
Lenguaje y framework & Python 3.14, Django 5.2 LTS & Consultas parametrizadas; escape por omisión \\
Base de datos & PostgreSQL 18 & Seguridad por fila; emisión única \\
Sesión & De servidor, cookie opaca & \dato{Se revoca en el acto} \\
Autorización & Open Policy Agent 1.21 & Un solo lugar; la matriz como política \\
Contenedores & Podman sobre AlmaLinux 10 & Sin privilegios; confinados por SELinux \\
\bottomrule
\end{tabular}}
\note{Para cada una el documento indica versión, alternativas descartadas, aporte, limitaciones y
herramientas de análisis.

Descartes por licencia: HashiCorp Vault (BUSL desde 2023) y Docker Desktop. Se usa OpenBao, que es
software libre.

Descartes por seguridad: tokens JWT, porque no se pueden revocar en el acto al dar de baja o cambiar
de rol; MariaDB, porque no tiene seguridad por fila; Docker Engine, por su servicio con privilegios; una
cadena de bloques, porque resuelve la desconfianza entre organizaciones, que aquí no existe.

Limitaciones que hay que saber decir: los permisos de Django son por tipo de objeto, no por objeto; su
sitio de administración no se instala; el superusuario de PostgreSQL se salta todo; el motor no conoce el
historial y no filtra listados.

Contraseñas: Argon2id con sal por usuario. Segundo factor: llave de acceso o códigos temporales; nunca SMS.}
\end{frame}

%--------------------------- 21. BITACORA ----------------------------
\begin{frame}{La bitácora no la reescribe\\ ni quien administra el servidor}
\quien{Axel}
\vspace{0.2in}
\begin{tikzpicture}[x=1in,y=1in,
  reg/.style={draw=regla,line width=0.8pt,fill=plano,minimum width=2.25in,minimum height=1.0in,
              align=center,inner sep=3pt},
  ext/.style={draw=acento,line width=1.6pt,fill=acentoclaro,minimum width=2.6in,minimum height=1.0in,
              align=center,inner sep=3pt},
  fl/.style={-{Stealth[length=9pt,width=8pt]},draw=gris,line width=1.2pt}]
\draw[draw=gris,densely dashed,line width=0.7pt] (-0.1,-0.8) rectangle (8.55,1.15);
\node[font=\tCuatro\bfseries,text=gris,anchor=north west] at (0,1.1) {SERVIDOR: SOLO SE PUEDE INSERTAR};
\node[reg] (a) at (1.25,-0.05) {Registro 1041\\{\tCuatro huella del 1040}};
\node[reg] (b) at (4.2,-0.05) {Registro 1042\\{\tCuatro huella del 1041}};
\node[reg] (c) at (7.15,-0.05) {Registro 1043\\{\tCuatro huella del 1042}};
\draw[fl] (b) -- (a); \draw[fl] (c) -- (b);
\node[ext,minimum width=2.25in] (t) at (10.97,-0.05) {\textbf{Testigo externo}\\{\tCuatro sello de tiempo y auditor}};
\draw[fl,draw=acento] (c) -- node[above,pos=0.66,font=\tCuatro,text=gris]{cada día} (t);
\end{tikzpicture}

\vspace{0.3in}
Cambiar un registro rompe la cadena. Reescribirla entera no coincide con lo que ya salió.
\note{Cuatro mecanismos que se apoyan: (1) la aplicación solo puede insertar, no modificar ni borrar;
(2) cada registro incluye la huella SHA-256 del anterior; (3) una vez al día la última huella sale del
servidor, recibe un sello de tiempo RFC 3161 y va al buzón de la Gerencia y del auditor; (4) la acción y
su registro se confirman en la misma transacción.

Qué se registra: autenticación, cambios, consultas de datos sensibles, administración y toda decisión
denegada. Con qué campos: instante, cuenta y rol, acción, objeto y versión, resultado y regla, origen,
huellas.

Conservación: diez años para cambios y decisiones, que es la trazabilidad que el cliente debe poder
reconstruir (minuta, P-18); dos para consultas (supuesto P-19).

Secretos: viven en una bóveda; la llave de sello no sale, se le pide a la bóveda que firme. La bóveda se
abre con dos de tres custodios.}
\end{frame}

%--------------------------- 22. RESIDUALES --------------------------
\begin{frame}{Lo que no se cubre está declarado}
\quien{Javier}
\vspace{0.0in}
\begin{tabular}{@{}L{5.0in} L{6.7in}@{}}
\toprule
\cabt{RIESGO RESIDUAL} & \cabt{POR QUÉ SE ACEPTA} \\
\midrule
Colusión entre técnico y verificador & Evitarla exige una tercera persona por certificado \\
Una medición falsa desde el origen & El sistema no ve el árbol \\
Extracción lenta, bajo los umbrales & Bajarlos castiga a quien trabaja bien \\
Un vínculo con un dueño que nadie declaró & No existe un registro que lo descarte \\
El histórico de las hojas de cálculo & No se le puede dar la trazabilidad que no tuvo \\
\bottomrule
\end{tabular}
\note{Son trece riesgos residuales en el documento de arquitectura. Otros: equipo de usuario infectado; el
operador puede detener o destruir (se detecta, no se impide); un solo servidor; sin hardware de
seguridad dedicado; la fecha de medición es declarada, porque no hay aplicación sin conexión; el
proveedor de hospedaje tiene acceso físico; la aprobación no lleva firma digital certificada; y la
misma regla vive en dos lugares, el motor y la base.

Si preguntan qué requerimiento quedó sin cubrir: la trazabilidad de diez años hacia atrás. Las hojas de
cálculo se congelan, se les calcula la huella y lo migrado queda marcado como histórico.

Antes de entrar en operación hay diez criterios de aceptación, entre ellos una prueba de penetración
autorizada por escrito.}
\end{frame}

%--------------------------- 23. RESPUESTAS DEL CLIENTE -------------
\begin{frame}{El cliente respondió,\\ y el diseño cambió en cinco puntos}
\quien{Javier}
\vspace{0.0in}
\begin{tabular}{@{}L{0.9in} L{5.0in} L{5.6in}@{}}
\toprule
 & \cabt{LO QUE RESPONDIÓ} & \cabt{QUÉ CAMBIÓ} \\
\midrule
P-01 & Un verificador puede medir, por excepción & Habilitación temporal; no aprueba lo que midió \\
P-05 & Las fincas las asignan los supervisores & \dato{Rol nuevo: supervisor de campo} \\
P-06 & La revisión sigue después del cierre & El cierre detiene registrar, no aprobar \\
P-07 & Revisión técnica antes de autorizar & Solicitar, corregir, revisar, autorizar \\
P-10 & Avisar primero; no bloquear de inmediato & Aviso, verificación, restricción temporal \\
\bottomrule
\end{tabular}

\vspace{0.2in}
\pie{Catorce de veintiséis consultas, respondidas por escrito. Las otras doce son supuestos declarados. Minuta: informe de políticas, Anexo A.}
\note{Qué decir: el diseño se hizo completo con supuestos declarados, y la consulta sirvió como validación.
De catorce respuestas: dos confirmaron el supuesto (P-03, el verificador devuelve y no edita; P-11, el
técnico solo ve sus fincas), cuatro lo ajustaron, seis lo cambiaron y dos quedaron a nuestro criterio.

Lo que resistió: la técnica de separación. La respuesta a P-01 no la cambió; solo volvió dinámica la
exclusión entre técnico y verificador.

Lo que hubo que rehacer: el orden de la enmienda (P-07), lo que el cierre detiene (P-06) y quién asigna
las fincas (P-05). Y el diseño era más estricto de lo que el cliente quería: bloqueaba de inmediato
donde bastaba avisar (P-10).

Otros cambios para citar si preguntan: sociedades y representantes como dueños, y declaración de
vínculos con abstención (P-15); ninguna norma de auditoría se da por supuesta, y los auditores exportan
evidencia sellada (P-18); registro después de la visita, con fecha de medición y fecha de registro
(P-22); diseño sin cantidades fijas (P-02).

Cada cambio llegó hasta la matriz: diecisiete casos de prueba nuevos, T-21 a T-37, y cinco modificados.}
\end{frame}

%--------------------------- 24. DECISIONES DELEGADAS ----------------
\begin{frame}{Cuatro decisiones que el cliente\\ dejó en nuestras manos}
\quien{Javier}
\vspace{0.0in}
\begin{tabular}{@{}L{4.7in} L{7.0in}@{}}
\toprule
\cabt{LA PREGUNTA} & \cabt{LO QUE DECIDIMOS} \\
\midrule
¿Revisa un verificador distinto del original? & \dato{Sí, cuando la enmienda es material} \\
¿Hace falta firma digital certificada? & No: sello de la empresa y comprobación en línea \\
¿Llaves físicas para todo el personal? & No: llave de acceso del propio equipo \\
¿Dónde se hospeda y quién lo opera? & Servidor virtual contratado; personal designado \\
\bottomrule
\end{tabular}

\vspace{0.2in}
\pie{El cliente pidió que cada decisión tenga una justificación clara y responda a un riesgo real, no el sistema más complejo.}
\note{Verificador distinto (minuta, P-07): tres opciones. Siempre distinto: toda corrección, hasta una de
forma, espera a otra persona. El mismo: quien aprobó el dato erróneo dictamina sobre su propio error.
Elegimos según lo que cambia: es material si cambia un dato del cálculo de toneladas o si ya hay un
certificado; entonces revisa otro. Lo determina el sistema, no una persona.

Firma digital (P-04): la firma certificada de cada verificador da autenticidad y autoría con presunción
legal, pero no dice si el certificado fue revocado después, y cuesta por persona. Elegimos sello de la
empresa, comprobación en línea por código y aprobación personal con segundo factor, que queda en la
bitácora.

Llaves físicas (P-23): no hay presupuesto para todos. La llave de acceso del propio dispositivo
(passkey) resiste igual el correo falso y no cuesta nada. Se recomiendan llaves físicas solo para
Gerencia, administrador, operador y custodios de la bóveda.

Hospedaje (P-23): servidor propio, servidor virtual contratado o servicios administrados, con los
cuatro criterios del cliente. El propio exige inversión y cuidar energía y enlace en una oficina
rural; los administrados no son software libre. Lo cotidiano está automatizado.}
\end{frame}

%--------------------------- 25. REFERENCIAS -------------------------
\begin{frame}{Referencias principales}
\quien{}
\vspace{0.0in}
{\tCuatro
\begin{columns}[T,onlytextwidth]
\begin{column}{0.485\textwidth}
\raggedright
Abadi, M., \& Fournet, C. (2003). Access control based on execution history. \emph{NDSS}.\\[5pt]
Clark, D. D., \& Wilson, D. R. (1987). A comparison of commercial and military computer security
policies. \emph{IEEE S\&P}.\\[5pt]
Kuhn, D. R., Coyne, E. J., \& Weil, T. R. (2010). Adding attributes to role-based access control.
\emph{Computer, 43}(6).\\[5pt]
Li, N., Tripunitara, M. V., \& Bizri, Z. (2007). On mutually exclusive roles and separation-of-duty.
\emph{ACM TISSEC, 10}(2).\\[5pt]
Saltzer, J. H., \& Schroeder, M. D. (1975). The protection of information in computer systems.
\emph{Proc. IEEE, 63}(9).\\[5pt]
Sandhu, R. S., Coyne, E. J., Feinstein, H. L., \& Youman, C. E. (1996). Role-based access control
models. \emph{Computer, 29}(2).
\end{column}
\begin{column}{0.485\textwidth}
\raggedright
Schneier, B., \& Kelsey, J. (1999). Secure audit logs to support computer forensics. \emph{ACM
TISSEC, 2}(2).\\[5pt]
Simon, R. T., \& Zurko, M. E. (1997). Separation of duty in role-based environments. \emph{CSFW}.\\[5pt]
Temoshok, D., et al. (2025). \emph{Digital identity guidelines: Authentication and authenticator
management} (NIST SP 800-63B-4).\\[5pt]
Ley N.º 8968 (2011) y su reglamento, Decreto N.º 37554-JP. Ley N.º 7975 (2000). Ley N.º 8454
(2005). Ley N.º 9048 (2012).\\[5pt]
ISO/IEC 17029:2019. ISO 14065:2020. ISO/IEC 27001:2022.\\[5pt]
Vargas Rivera, A. (2026). \emph{Material de estudio del curso CE-1115 y enunciado del proyecto,
parte 1}. ITCR.
\end{column}
\end{columns}
}
\note{La lista completa, en formato APA, está en cada uno de los cinco documentos.}
\end{frame}

\end{document}
"
%=====================================================================
\documentclass[aspectratio=169,t]{beamer}

\usepackage[utf8]{inputenc}
\usepackage[T1]{fontenc}
\usepackage[spanish,es-noshorthands,es-tabla]{babel}
\usepackage[default]{lato}
\usepackage{booktabs,array,colortbl}
\usepackage{tikz}
\usetikzlibrary{positioning,arrows.meta,calc}

\ifdefined\connotas\setbeameroption{show notes}\fi

%--------------------------- lamina y margenes -----------------------
\geometry{papersize={13.333in,7.5in}}
\setbeamersize{text margin left=0.6in,text margin right=0.6in}

%--------------------------- color -----------------------------------
\definecolor{acento}{HTML}{0F6E6A}
\definecolor{acentoclaro}{HTML}{E7F2F1}
\definecolor{tinta}{HTML}{1A1A1A}
\definecolor{gris}{HTML}{5F5E5A}
\definecolor{regla}{HTML}{B9B8B0}
\definecolor{plano}{HTML}{FAFAFA}
\definecolor{negado}{HTML}{F2F1EC}

%--------------------------- escala: cuatro tamanos ------------------
\newcommand{\tUno}{\fontsize{40}{45}\selectfont}
\newcommand{\tDos}{\fontsize{28}{35}\selectfont}
\newcommand{\tTres}{\fontsize{18}{25}\selectfont}
\newcommand{\tCuatro}{\fontsize{13}{17.5}\selectfont}
\makeatletter
\renewcommand\normalsize{\@setfontsize\normalsize{18}{25}}
\renewcommand\small{\@setfontsize\small{13}{17.5}}
\renewcommand\footnotesize{\@setfontsize\footnotesize{13}{17.5}}
\renewcommand\scriptsize{\@setfontsize\scriptsize{13}{17.5}}
\renewcommand\tiny{\@setfontsize\tiny{13}{17.5}}
\renewcommand\large{\@setfontsize\large{28}{35}}
\renewcommand\Large{\@setfontsize\Large{40}{45}}
\makeatother

%--------------------------- cromo -----------------------------------
\usetheme{default}
\setbeamertemplate{navigation symbols}{}
\setbeamercolor{structure}{fg=acento}
\setbeamercolor{normal text}{fg=tinta}
\setbeamercolor{background canvas}{bg=white}
\setbeamerfont{note page}{size=\fontsize{20}{27}\selectfont}

\setbeamertemplate{frametitle}{%
  \vskip0.0in%
  \begin{minipage}[t]{\textwidth}
    \raggedright\color{tinta}\tUno\bfseries\insertframetitle\par
  \end{minipage}\par
  \vskip0.12in}

\gdef\presentador{}
\newcommand{\quien}[1]{\gdef\presentador{#1}}
\setbeamertemplate{footline}{%
  \hbox to\paperwidth{\hskip0.6in{\tCuatro\color{gris}\insertframenumber}\hfill
    {\tCuatro\color{gris}\presentador}\hskip0.6in}%
  \vskip0.45in}

%--------------------------- atajos ----------------------------------
\newcommand{\dato}[1]{\textcolor{acento}{\textbf{#1}}}
\newcommand{\pie}[1]{\leavevmode{\tCuatro\color{gris}#1}}
\newcommand{\cabt}[1]{\leavevmode{\tCuatro\bfseries\color{gris}#1}}
\newcolumntype{L}[1]{>{\raggedright\arraybackslash}p{#1}}
\arrayrulecolor{regla}
\setlength{\arrayrulewidth}{0.8pt}
\renewcommand{\arraystretch}{1.5}
% celdas de la matriz
\newcommand{\vb}[2]{\mbox{#1\,{\color{acento}\bfseries #2}}}
\newcommand{\sinacc}{\cellcolor{negado}\leavevmode{\color{gris}---}}
\newcommand{\suj}[1]{\leavevmode{\tTres\bfseries #1}}

%=====================================================================
\begin{document}

%--------------------------- 1. PORTADA ------------------------------
\begin{frame}[plain]
\vspace*{1.25in}
{\tCuatro\color{gris}PROYECTO \ $\cdot$ \ PARTE 1 \ $\cdot$ \ DEFENSA ORAL}\\[0.3in]
{\tUno\bfseries\color{tinta}Diseño del control de acceso\\[4pt] para Verdeazul Carbono}\\[0.25in]
{\color{acento}\rule{1.6in}{4pt}}\\[0.25in]
{\tDos\color{gris}Políticas, matriz y arquitectura de seguridad}\\[0.7in]
{\tTres\bfseries Randall Bolaños \quad Javier Tenorio \quad Julio Varela \quad Axel Flores}\\[0.12in]
{\tCuatro\color{gris}Grupo 6 \ $\cdot$ \ CE-1115 Seguridad de la Información \ $\cdot$ \ Prof. MSc. Andrés
Vargas Rivera \ $\cdot$ \ ITCR \ $\cdot$ \ 15 de octubre de 2026}
\note{Reparto propuesto. Randall: cliente, activos y separación de responsabilidades. Javier:
amenazas, marco legal, fricción, riesgos residuales y respuestas del cliente. Julio: modelo de control
de acceso y matriz. Axel: controles de historial y arquitectura. Todos deben poder responder por cualquier lámina:
las preguntas del profesor son individuales.}
\end{frame}

%--------------------------- 2. EL NEGOCIO ---------------------------
\begin{frame}{Verdeazul vende confianza:\\ cada certificado respalda una tonelada real}
\quien{Randall}
\vspace{0.15in}
\begin{tikzpicture}[x=1in,y=1in,
  caja/.style={draw=regla,line width=0.8pt,fill=plano,minimum width=2.0in,minimum height=0.95in,
               align=center,font=\tTres,inner sep=2pt},
  fl/.style={-{Stealth[length=9pt,width=8pt]},draw=gris,line width=1.2pt}]
\node[caja] (a) at (1.00,0) {Finca\\forestal};
\node[caja] (b) at (3.53,0) {Técnico\\mide};
\node[caja] (c) at (6.06,0) {Verificador\\aprueba};
\node[caja,draw=acento,line width=1.6pt] (d) at (8.59,0) {\textbf{Certificado}\\de carbono};
\node[caja] (e) at (11.12,0) {Mercado\\internacional};
\foreach \p/\q in {a/b,b/c,c/d,d/e} \draw[fl] (\p) -- (\q);
\end{tikzpicture}

\vspace{0.4in}
{\tDos \dato{Un dato inflado es fraude.}}

\vspace{0.12in}
Auditores europeos pueden pedir la trazabilidad de cualquier certificado de los últimos diez años.
\note{Qué decir: Verdeazul certifica captura de carbono en fincas forestales. El técnico mide, el
verificador aprueba y se emite un certificado que el dueño vende afuera. El valor del certificado es la
confianza en el dato. Hoy todo vive en hojas de cálculo compartidas.

Si preguntan por qué empezamos por el negocio: porque el enunciado pide políticas que describan a este
cliente, y casi todas las decisiones salen de esta lámina.}
\end{frame}

%--------------------------- 3. LOS DOS SUSTOS -----------------------
\begin{frame}{Los dos sustos de la carta son fallas de integridad}
\quien{Randall}
\vspace{0.1in}
\begin{tabular}{@{}L{2.2in} L{9.5in}@{}}
\toprule
\cabt{PRIMER SUSTO} & Una medición cambió después de aprobada, y nadie sabe quién fue. \\
\midrule
\cabt{SEGUNDO SUSTO} & Un certificado se emitió dos veces sobre la misma hectárea. \\
\bottomrule
\end{tabular}

\vspace{0.45in}
{\tDos Orden de prioridad: \dato{integridad}, confidencialidad, disponibilidad.}

\vspace{0.12in}
Cuando dos pilares chocan, gana la integridad: sin registro en bitácora no hay acción.
\note{Qué decir: ninguno de los dos incidentes fue un robo de información. En los dos, el dato dejó de
ser confiable. Por eso la integridad va primero. La confidencialidad va segunda por la competencia
entre dueños y por los datos personales. La disponibilidad va tercera, con un matiz: importa más que el
dato exista en diez años que un día sin sistema.

Si preguntan por un ejemplo de «gana la integridad»: si la bitácora no se puede escribir, la operación
se rechaza aunque eso detenga el trabajo (POL-BT-06).}
\end{frame}

%--------------------------- 4. ACTIVOS ------------------------------
\begin{frame}{Lo que más vale es lo que sostiene un certificado}
\quien{Randall}
\vspace{0.05in}
\begin{tabular}{@{}L{4.6in} L{2.6in} L{4.4in}@{}}
\toprule
\cabt{ACTIVO} & \cabt{CLASIFICACIÓN} & \cabt{LO QUE NO PUEDE FALLAR} \\
\midrule
Mediciones y certificados & Confidencial & \dato{Integridad} \\
Catálogo de fincas: la cartera & Confidencial & Integridad y confidencialidad \\
Datos personales de los dueños & Confidencial & Confidencialidad \\
Bitácora y llave de sello & Restringido & \dato{Integridad} \\
\bottomrule
\end{tabular}

\vspace{0.3in}
\pie{El inventario completo tiene once activos, cada uno valorado en los tres pilares (informe de políticas, sección 2).}
\note{Qué decir: clasificamos en cuatro niveles (público, interno, confidencial, restringido) y valoramos
cada activo con la escala de FIPS 199. La Gerencia de Certificación es la responsable del dato; el
administrador y el operador son custodios.

Si preguntan por un activo que no es obvio: los parámetros de cálculo (A-10). Un factor alterado infla
todos los certificados a la vez. Y la llave de sello (A-09): quien la tenga fabrica certificados
creíbles.}
\end{frame}

%--------------------------- 5. AMENAZAS -----------------------------
\begin{frame}{\textcolor{acento}{7 de 11} amenazas son internas:\\ el atacante más probable ya tiene una cuenta}
\quien{Javier}
\vspace{0.05in}
\begin{tabular}{@{}L{1.4in} L{6.6in} L{3.6in}@{}}
\toprule
\cabt{ID} & \cabt{YA OCURRIÓ EN VERDEAZUL} & \cabt{PILAR} \\
\midrule
AM-03 & Una medición editada después de aprobarse & Integridad \\
AM-06 & Un certificado emitido dos veces & Integridad \\
AM-04 & Cientos de fincas consultadas en un fin de semana & Confidencialidad \\
\bottomrule
\end{tabular}

\vspace{0.3in}
El diseño no mira el perímetro: mira qué puede hacer cada persona autorizada.
\note{Qué decir: el modelo tiene once amenazas, cada una con agente, motivo, activo, la debilidad que
hoy lo permite y el riesgo. Siete son internas: técnico que infla, misma persona que mide y aprueba,
edición silenciosa, extracción de la cartera, administrador u operador que abusa, doble emisión por
error y, desde la consulta al cliente, el empleado con un vínculo con un dueño (AM-11): no existe un
registro de esos vínculos. Tres ya ocurrieron.

Externas: dueño que quiere ver a un competidor, robo de credenciales, ransomware y certificado
falsificado o cuenta de auditor tomada.

Si preguntan cómo calculamos el riesgo: probabilidad por impacto. Probabilidad alta es que ya ocurrió
aquí o que es la vía habitual de ataque en el país.}
\end{frame}

%--------------------------- 6. MARCO LEGAL --------------------------
\begin{frame}{Cada obligación legal termina\\ en una regla que se puede comprobar}
\quien{Javier}
\vspace{0.05in}
\begin{tabular}{@{}L{2.9in} L{4.2in} L{4.5in}@{}}
\toprule
\cabt{NORMA} & \cabt{EXIGE} & \cabt{REGLA EN VERDEAZUL} \\
\midrule
Ley 8968, art. 14 & Autorización para transferir datos & El auditor consulta en la plataforma \\
Reglamento, art. 38 & Avisar una fuga en 5 días hábiles & Procedimiento con ese plazo \\
Ley 7975, art. 2 & Medidas razonables de secreto & Acceso por necesidad y registro \\
ISO/IEC 17029 & Quien revisa no participó & \dato{Medir y aprobar, separados} \\
\bottomrule
\end{tabular}

\vspace{0.2in}
\pie{El cliente aclaró que no audita bajo una norma determinada: la ISO se cita como referencia. Su exigencia firme es la trazabilidad de diez años.}
\note{Qué decir: la Ley 8968 aplica porque los datos de los dueños se comparten con auditores externos;
no es una base de uso exclusivamente interno. Derechos del titular en cinco días hábiles; consentimiento
informado al dar de alta a un dueño; conservación hasta diez años y después se desasocia.

La Ley 7975 es la que permite reclamar si la cartera llega a la competencia, pero solo si la empresa
puede demostrar que tomó medidas razonables: por eso el acceso por asignación y la bitácora tienen
además valor legal.

También: Código Penal (arts. 196 bis, 217 bis, 229 bis y 231), Código de Trabajo art. 71 g) y Ley 8454
de firma digital.

Auditores, según la minuta (P-18): una vez al año, varios días, en sitio o a distancia. Pueden exportar
evidencia, pero como paquetes que genera el sistema, sellados, registrados y sin datos personales. El
acceso vence con la ventana autorizada; prorrogarlo es una autorización nueva.}
\end{frame}

%--------------------------- 7. MODELO: TRES CAPAS -------------------
\begin{frame}{Roles con restricciones,\\ y condiciones que solo pueden quitar}
\quien{Julio}
\vspace{0.2in}
\begin{tikzpicture}[x=1in,y=1in,
  capa/.style={draw=acento,line width=1.4pt,fill=acentoclaro,minimum width=2.4in,minimum height=1.1in,
               align=center,inner sep=3pt},
  ext/.style={draw=regla,line width=0.8pt,fill=plano,minimum height=1.1in,align=center,inner sep=3pt},
  fl/.style={-{Stealth[length=9pt,width=8pt]},draw=gris,line width=1.2pt},
  no/.style={-{Stealth[length=9pt,width=8pt]},draw=gris,line width=1.0pt,densely dashed}]
\node[ext,minimum width=1.6in] (p) at (0.8,0) {Petición\\autenticada};
\node[capa] (r) at (3.25,0) {\textbf{1 $\cdot$ Roles}\\{\tCuatro ¿Este rol puede?}};
\node[capa] (c) at (6.10,0) {\textbf{2 $\cdot$ Condiciones}\\{\tCuatro ¿Aquí y ahora?}};
\node[capa] (h) at (8.95,0) {\textbf{3 $\cdot$ Historial}\\{\tCuatro ¿Encaja con lo previo?}};
\node[ext,minimum width=1.5in] (d) at (11.35,0) {Conceder};
\draw[fl] (p) -- (r);
\draw[fl] (r) -- node[above,font=\tCuatro,text=gris]{sí} (c);
\draw[fl] (c) -- node[above,font=\tCuatro,text=gris]{sí} (h);
\draw[fl] (h) -- node[above,font=\tCuatro,text=gris]{sí} (d);
\node[draw=regla,line width=0.8pt,fill=negado,minimum width=8.1in,minimum height=0.6in] (n) at (6.10,-1.55)
  {Denegar y registrar en bitácora};
\foreach \k in {r,c,h} \draw[no] (\k.south) -- node[right,font=\tCuatro,text=gris]{no} (\k.south |- n.north);
\end{tikzpicture}

\vspace{0.3in}
\dato{Solo la primera capa concede.} La matriz muestra lo máximo que cada rol puede llegar a hacer.
\note{Qué decir: elegimos RBAC con restricciones, nivel RBAC2 de Sandhu. Sobre esa base, un conjunto
cerrado de atributos decide si el permiso aplica a esta finca, en este estado y en este momento.
Es el enfoque centrado en roles de Kuhn, Coyne y Weil.

La frontera del híbrido: la capa 1 dice si el rol tiene el verbo; la capa 2 si aplica a este objeto ahora;
la capa 3 si encaja con el historial. Se combinan sin contradecirse porque solo la tabla de roles concede
y las otras dos son filtros: una negativa en cualquier capa gana.

Si preguntan por qué RBAC: la carta describe funciones, no personas; la gente rota; y hay que poder
explicarlo a un auditor leyendo una fila.}
\end{frame}

%--------------------------- 8. DESCARTES ----------------------------
\begin{frame}{Los otros tres modelos fallan\\ por razones de este caso}
\quien{Julio}
\vspace{0.05in}
\begin{tabular}{@{}L{2.2in} L{9.5in}@{}}
\toprule
\cabt{MODELO} & \cabt{POR QUÉ NO EN VERDEAZUL} \\
\midrule
DAC & Es lo que hay hoy: las hojas compartidas donde ocurrieron los dos sustos. \\
MAC & No hay niveles de secreto. Haría falta un compartimento por finca. \\
ABAC puro & Desaparece la fila que responde «¿qué puede hacer un verificador?». \\
\bottomrule
\end{tabular}

\vspace{0.3in}
\pie{De cada uno se conserva algo: la titularidad como atributo (DAC), la política que el usuario no puede
alterar (MAC) y las condiciones acotadas (ABAC).}
\note{DAC: el técnico que crea la medición sería su dueño y podría seguir modificándola o compartirla.
Es el primer susto.

MAC: Bell-LaPadula ordena por niveles de secreto y protege la confidencialidad. Aquí el riesgo principal
es de integridad. La variante de integridad, Biba, pone una etiqueta fija, pero la confianza en una
medición cambia con su estado y depende de quién hizo el paso anterior. Es el argumento de Clark y
Wilson. MAC sí lo usamos en la plataforma: SELinux confina los contenedores.

ABAC puro: podría expresarlo todo, pero no se puede revisar leyendo una fila, y la empresa no tiene
quién gobierne un catálogo de atributos.}
\end{frame}

%--------------------------- 9. NIVEL Y EXPLOSION --------------------
\begin{frame}{\textcolor{acento}{Siete} roles hoy, y siete con el doble de fincas}
\quien{Julio}
\vspace{0.25in}
{\tDos El alcance nunca es un rol: es un atributo.}

\vspace{0.35in}
\begin{tabular}{@{}L{3.3in} L{8.4in}@{}}
\toprule
RBAC$_2$, sin jerarquía & Si el verificador heredara del técnico, podría registrar: se rompe la separación. \\
Explosión de roles & Aparece al usar un rol para decir dónde o cuándo. Aquí eso vive en las condiciones. \\
La consulta lo probó & Cinco necesidades nuevas: \dato{un rol y cuatro atributos}. \\
\bottomrule
\end{tabular}
\note{Los siete roles: técnico, verificador, supervisor de campo, gerencia de certificación,
administrador de catálogo y cuentas, auditor externo y dueño de finca. Más una identidad de servicio
(SIS), el operador de plataforma, que no tiene rol en la aplicación, y el público sin autenticar.

El séptimo rol salió de la consulta: el cliente dijo que las fincas las asignan los supervisores del
personal de campo (minuta, P-05). Las otras cuatro necesidades nuevas (el verificador que mide por
excepción, el técnico que apoya a un compañero, la sociedad con un representante y el empleado con un
vínculo con un dueño) no son funciones de trabajo: se resolvieron con atributos.

Por qué RBAC2 y no RBAC3: no hay cadena de mando entre las siete funciones; la herencia vuelve
implícitos los permisos y la matriz no debe tener nada implícito; y con siete roles no ahorra nada.

Restricciones: un rol por persona, una identidad por persona, al menos dos verificadores, cuenta de
auditor con ventana. La pareja técnico--verificador es la única excepción: es dinámica. Un verificador
puede recibir una habilitación temporal de técnico, con fecha de fin, y nunca aprueba lo que midió
(minuta, P-01).

Si preguntan dónde está el riesgo de explosión: un rol por finca para aislar a los dueños, o roles como
«Técnico Zona Sur 2027».}
\end{frame}

%--------------------------- 10. HBAC --------------------------------
\begin{frame}{Cuatro controles de historial, uno de cada tipo}
\quien{Axel}
\vspace{0.0in}
\begin{tabular}{@{}L{0.7in} L{1.7in} L{4.6in} L{4.4in}@{}}
\toprule
 & \cabt{TIPO} & \cabt{QUÉ LO DISPARA} & \cabt{QUÉ DECIDE} \\
\midrule
H1 & Frecuencia & Muchas fincas distintas en un día & \dato{Tres niveles}: avisa, confirma, restringe \\
H2 & Tiempo & Una decisión fuera de jornada & Pide segundo factor y motivo \\
H3 & Combinada & Fuera de jornada, muchas fincas, casi todas nuevas & Confirma y avisa; si insiste, restringe \\
H4 & Secuencia & Quien participó antes quiere decidir & \dato{Deniega} \\
\bottomrule
\end{tabular}

\vspace{0.2in}
\pie{Ninguna señal aislada restringe: el cliente pidió avisar primero y no bloquear de inmediato (minuta, P-10).}
\note{Son reglas explícitas, no aprendizaje automático: se le pueden leer a un auditor y se pueden probar.

Cómo se distingue lo legítimo: se cuentan señales: volumen, novedad, momento e insistencia. El trabajo
extraordinario legítimo muestra una; la extracción de la cartera las muestra juntas.

Tres niveles: aviso (el usuario no nota nada), verificación adicional (segundo factor y motivo) y
restricción temporal (no abre fincas nuevas hasta el siguiente día hábil; la cuenta sigue activa).

H1 mitiga AM-04. Umbrales por rol, calibrados con datos reales; iniciales 40, 60 y 80 fincas por día
(supuesto P-13). Los intentos denegados cuentan doble.

H2 mitiga AM-08, la cuenta robada usada de madrugada. H4 es la separación de responsabilidades como
control de historial: reglas C3, C8 y C10.

H3 es la firma del incidente: fin de semana, cientos de fincas, ninguna suya. Pide confirmación en la
finca once y avisa; restringe en la veintiuna.

Falsos positivos: verificador con la cola acumulada al cierre, o nuevo y sin historial. El primer
período corre en observación para calibrar.

Descartada: Muralla China. Todos los dueños compiten: cada empleado quedaría con una sola finca.}
\end{frame}

%--------------------------- 11. EL TECNICO --------------------------
\begin{frame}{El técnico del fin de semana\\ hoy no vería ni una finca ajena}
\quien{Axel}
\vspace{0.05in}
\begin{tabular}{@{}L{4.6in} L{7.1in}@{}}
\toprule
\cabt{LO QUE HACE} & \cabt{LO QUE HACE EL SISTEMA} \\
\midrule
Abre una finca que no tiene asignada & \dato{Deniega}: su rol no ve fincas ajenas \\
Insiste tres, cinco, diez veces & Avisa a su supervisor; pide segundo factor; lo detiene media hora \\
La necesitaba de verdad & El supervisor le da una asignación temporal, que vence sola \\
Presenta su renuncia & Reporte de sus consultas de 90 días \\
\bottomrule
\end{tabular}

\vspace{0.25in}
\pie{El cliente lo confirmó: fincas asignadas con flexibilidad, y respuesta según la gravedad (minuta, P-10 y P-11).}
\note{Qué decir: la carta dice «no sabemos qué hacer con eso, pero nos gustaría que el sistema al menos se
diera cuenta». Nuestra primera respuesta no es vigilar el permiso: es quitarlo. Un técnico solo ve las
fincas que tiene asignadas (regla C1).

El cliente lo confirmó, con una condición: flexibilidad para apoyar a un compañero, sustituir a un
ausente o atender una visita extraordinaria, sin acceso permanente a todas las fincas. Se resuelve con
asignaciones temporales que crea el supervisor de campo, con motivo y fecha de fin (regla C19).

El diseño original suspendía la sesión a la quinta denegación. El cliente pidió no bloquear de
inmediato, y se cambió: aviso al supervisor a la tercera, segundo factor a la quinta, media hora de
restricción a la décima.

Los controles de historial quedan para quienes sí necesitan leer muchas fincas: verificador, supervisor,
gerencia, administrador y auditor. Si el técnico hubiera sido verificador, H3 le pide confirmación en la
finca once y lo restringe en la veintiuna.

Además hay una vía legal: confidencialidad firmada, que persiste después de la salida (Ley 8968 art. 11
y Ley 7975).}
\end{frame}

%--------------------------- 12. NADIE EMITE -------------------------
\begin{frame}{Nadie emite un certificado: lo hace el sistema}
\quien{Randall}
\vspace{0.25in}
\begin{tikzpicture}[x=1in,y=1in,
  paso/.style={draw=acento,line width=1.4pt,fill=acentoclaro,minimum width=3.3in,minimum height=1.1in,
               align=center,inner sep=3pt},
  sis/.style={draw=gris,line width=1.2pt,fill=negado,minimum width=3.3in,minimum height=1.1in,
              align=center,inner sep=3pt},
  fl/.style={-{Stealth[length=9pt,width=8pt]},draw=gris,line width=1.2pt},
  nota/.style={font=\tCuatro,text=gris,align=center}]
\node[paso] (a) at (1.65,0) {\textbf{Técnico}\\registra y envía};
\node[paso] (b) at (6.05,0) {\textbf{Verificador}\\aprueba o devuelve};
\node[sis]  (c) at (10.45,0) {\textbf{Sistema}\\emite};
\draw[fl] (a) -- (b); \draw[fl] (b) -- (c);
\node[nota,below=0.08in of a] {solo en fincas asignadas\\y con el período abierto};
\node[nota,below=0.08in of b] {no registró ninguna versión\\de esa medición};
\node[nota,below=0.08in of c] {no hay otro certificado vigente\\para la parcela y el período};
\end{tikzpicture}

\vspace{0.35in}
Ninguna persona tiene el permiso de emitir. \dato{Hacen falta dos, y distintas.}
\note{La acción crítica: emitir un certificado. Si una sola persona pudiera medir y aprobar, podría
inflar la biomasa, aprobar su propio dato y obtener un certificado por toneladas que no existen.

La segunda acción crítica: corregir una medición ya congelada. Es el primer susto, y se trata igual,
con el orden que pidió el cliente (minuta, P-07): solicita el técnico responsable, un verificador o la
Gerencia; corrige el técnico; revisa un verificador que no es autor ni solicitante; y la Gerencia da la
autorización final. Son al menos tres personas.

Si preguntan qué pasa con el certificado tras una enmienda (minuta, P-08): si las toneladas no cambian,
sigue vigente con la enmienda anotada. Si cambian, el sistema lo revoca y emite el reemplazo en un solo
acto, enlazados: nunca hay dos vigentes y ninguno se borra. Si ya se vendió, la Gerencia lo registra y
coordina con las partes.}
\end{frame}

%--------------------------- 13. POR MEDICION ------------------------
\begin{frame}{La separación se comprueba\\ por medición, no por puesto}
\quien{Randall}
\vspace{0.05in}
\begin{tabular}{@{}L{3.4in} L{8.3in}@{}}
\toprule
\cabt{SI SE COMPRUEBA\ldots} & \cabt{UN TÉCNICO ASCENDIDO A VERIFICADOR\ldots} \\
\midrule
por el puesto de hoy & aprueba lo que él mismo midió el año anterior. \\
\dato{por la medición} & \dato{no puede: figura como autor.} \\
\bottomrule
\end{tabular}

\vspace{0.35in}
Técnica: \textbf{separación por etapas del flujo}, reforzada con roles mutuamente excluyentes.
\note{Qué decir: quien inicia el proceso no puede cerrarlo. El sistema no pregunta qué rol tiene la
persona, pregunta si participó en esta medición. En la literatura es la separación basada en el objeto
o en el historial (Nash y Poland 1990; Simon y Zurko 1997).

Por qué es la adecuada: apunta al riesgo tal como lo dijo el cliente; sobrevive a los cambios de puesto
(la lección de Kerviel); no depende de que los roles estén bien administrados; y deja dos nombres
distintos en cada certificado.

El cliente confirmó que algunos verificadores fueron técnicos y que por excepción pueden medir, y pidió
«garantizar una revisión independiente, sin impedir que el personal pueda asumir otras
responsabilidades» (minuta, P-01). Una separación por puesto obligaría a escoger; esta da las dos
cosas: el verificador mide con una habilitación temporal y no podrá aprobar nunca lo que midió.

Coincide con el principio de imparcialidad de ISO/IEC 17029, que citamos solo como referencia.}
\end{frame}

%--------------------------- 14. ALTERNATIVAS ------------------------
\begin{frame}{Las alternativas se descartaron\\ por lo que cuestan aquí}
\quien{Randall}
\vspace{0.0in}
\begin{tabular}{@{}L{4.3in} L{7.4in}@{}}
\toprule
\cabt{TÉCNICA} & \cabt{POR QUÉ NO ES LA PRINCIPAL} \\
\midrule
Roles excluyentes, solos & No cubren a quien cambia de puesto \\
Doble control (cuatro ojos) & \dato{Duplica la verificación de cada certificado} \\
Segregación temporal & Limita cuándo, no quién \\
Rotación & No impide un acto; se recomienda contra la colusión \\
\bottomrule
\end{tabular}

\vspace{0.2in}
\pie{La segregación temporal sí se usa donde el problema es la duración: la asignación temporal, la habilitación del verificador, el plazo de la enmienda y la cuenta del auditor.}
\note{Roles mutuamente excluyentes estáticos (RBAC2): se conservan como refuerzo. Li, Tripunitara y Bizri
muestran que excluir roles es un mecanismo para aproximar la separación, no la separación misma.

Exclusión dinámica: basta cerrar sesión y abrir otra. Sola, demora el fraude y no lo impide. Se usa
como refuerzo en la pareja técnico--verificador, detrás de la regla C3.

Doble control: era la opción más fuerte contra el fraude. Se descartó por costo operativo: duplica la
verificación de cada certificado, con cualquier tamaño de equipo.

Rotación: cuesta conocimiento de la finca y depende de cuántas personas haya, un dato que el cliente
pidió no fijar (minuta, P-02).}
\end{frame}

%--------------------------- 15. FRICCION ----------------------------
\begin{frame}{La fricción existe y está acotada}
\quien{Javier}
\vspace{0.0in}
\begin{tabular}{@{}L{5.2in} L{6.5in}@{}}
\toprule
\cabt{LO QUE CUESTA} & \cabt{CÓMO SE MANEJA} \\
\midrule
Siempre hacen falta dos personas & Mínimo dos verificadores; \dato{sin atajo de emergencia} \\
Las devoluciones quitan tiempo de campo & Validación de rangos al registrar \\
Una enmienda pide tres personas & Plazo de diez días; el sistema la clasifica \\
La Gerencia concentra las segundas firmas & Nunca inicia lo que aprueba \\
\bottomrule
\end{tabular}
\note{Qué decir: no existe un atajo para aprobar lo propio. Un atajo así sería el hueco que el auditor
buscaría. Si el único verificador presente participó, la medición espera; la cola avisa a la Gerencia.

Lo que la técnica no resuelve: la colusión. Si un técnico y un verificador se ponen de acuerdo, el flujo
se cumple. La separación sube el costo del fraude de una persona a dos. Quedan la rotación, un indicador
de parejas repetidas y la auditoría externa.

La Gerencia es un cuello de botella y un punto de confianza: firma en varios flujos, pero nunca inicia
lo que aprueba, no registra ni aprueba mediciones, ya no asigna fincas y no resuelve sus propias alertas.

Enmienda: revisa un verificador distinto del original solo cuando es material (cambia toneladas o ya hay
certificado); lo determina el sistema. Con solo dos verificadores, si uno solicita sobre lo que aprobó
el otro no queda revisor: la presenta el técnico o la Gerencia.

La flexibilidad también cuesta un trámite: la asignación temporal la crea el supervisor y la
habilitación de un verificador para medir pasa por el administrador y la Gerencia.}
\end{frame}

%--------------------------- 16. MATRIZ (1 de 2) ---------------------
\begin{frame}{Lo que no está escrito está denegado}
\quien{Julio}
\vspace{-0.02in}
{\tCuatro
\renewcommand{\arraystretch}{1.3}
\setlength{\tabcolsep}{5pt}
\begin{tabular}{@{}L{0.75in}|L{1.95in}|L{2.25in}|L{2.0in}|L{1.85in}|L{2.3in}@{}}
\rowcolor{acentoclaro}
 & \textbf{Fincas y parcelas} & \textbf{Datos personales del dueño} & \textbf{Períodos de medición} & \textbf{Asignaciones} & \textbf{Mediciones} \\
\hline
\suj{TEC} & \vb{leer}{C1} & \vb{leer}{C1} (solo contacto) & leer & \vb{leer}{C1} (las propias) & \vb{leer}{C1}, \vb{crear}{C2}, \vb{modificar}{C2}, \vb{enviar}{C2} \\
\hline
\suj{VER} & \vb{leer}{H} & \sinacc & leer & \sinacc & \vb{leer}{H}, \vb{aprobar}{C3}, \vb{devolver}{C3} \\
\hline
\suj{SUP} & \vb{leer}{H} & \sinacc & leer & leer, \vb{crear}{C19}, \vb{modificar}{C19} & \vb{leer}{H} (solo estado y fechas) \\
\hline
\suj{GER} & \vb{leer}{H}, \vb{aprobar}{C11} & \vb{leer}{H} & leer, \vb{abrir}{C14}, \vb{cerrar}{C14} & leer & \vb{leer}{H} \\
\hline
\suj{ADM} & \vb{leer}{H}, \vb{proponer}{C11} & \vb{leer}{H}, crear, modificar, \vb{eliminar}{C16} & leer, crear, \vb{proponer}{C14} & \sinacc & \sinacc \\
\hline
\suj{AUD} & \vb{leer}{C6}, \vb{exportar}{C18} & \vb{leer}{C6} & \vb{leer}{C6}, \vb{exportar}{C18} & \vb{leer}{C6}, \vb{exportar}{C18} & \vb{leer}{C6}, \vb{exportar}{C18} \\
\hline
\suj{DUE} & \vb{leer}{C5} & \vb{leer}{C5} & leer & \vb{leer}{C5} & \vb{leer}{C5} (solo aprobadas) \\
\hline
\suj{SIS} & leer & leer (solo el nombre del titular) & leer & leer & leer \\
\hline
\suj{OPE} & \sinacc & \sinacc & \sinacc & \sinacc & \sinacc \\
\hline
\suj{PUB} & \sinacc & \sinacc & \sinacc & \sinacc & \sinacc \\
\end{tabular}

\vspace{0.06in}
{\color{gris}Matriz real, objetos 1 a 5 de 10. \quad --- sin acceso \quad {\color{acento}\textbf{C}$n$} permiso sujeto a la
regla $n$ \quad {\color{acento}\textbf{H}} lectura vigilada por historial}
}
\note{Cómo leerla: sujetos en las filas, objetos en las columnas. Una celda gris es sin acceso. Una marca C
significa que el permiso solo vale si se cumple esa regla; las condiciones solo pueden quitar.

Recorrido sugerido: fila TEC, columna Mediciones: crea, modifica y envía, pero no aprueba. Fila VER:
aprueba y devuelve con C3, no crea ni modifica, y no ve los datos personales del dueño. Fila SUP, que
es nueva: asigna fincas (C19) y de las mediciones solo ve estado y fechas. Fila GER: abre y cierra
períodos que registró el administrador (C14); ya no crea asignaciones. Fila AUD: leer y exportar
evidencia, con ventana (C6 y C18), y ningún verbo de escritura. Fila DUE: solo lo suyo (C5). Filas OPE
y PUB: nada en esta mitad.

Diez sujetos: siete roles, el sistema, el operador y el público sin autenticar.

Mínimo privilegio: cada permiso amplio tiene su justificación en la sección 6 del documento de la matriz.}
\end{frame}

%--------------------------- 17. MATRIZ (2 de 2) ---------------------
\begin{frame}{Nadie tiene «eliminar», y nadie emite}
\quien{Julio}
\vspace{-0.02in}
{\tCuatro
\renewcommand{\arraystretch}{1.3}
\setlength{\tabcolsep}{5pt}
\begin{tabular}{@{}L{0.75in}|L{2.3in}|L{1.9in}|L{2.75in}|L{1.4in}|L{2.0in}@{}}
\rowcolor{acentoclaro}
 & \textbf{Enmiendas} & \textbf{Certificados} & \textbf{Cuentas y roles} & \textbf{Bitácora} & \textbf{Alertas de historial} \\
\hline
\suj{TEC} & \vb{leer}{C1}, \vb{solicitar}{C7}, \vb{corregir}{C9} & \sinacc & \vb{modificar}{C13} & \sinacc & \sinacc \\
\hline
\suj{VER} & leer, \vb{solicitar}{C7}, \vb{revisar}{C10} & \vb{leer}{H} & \vb{modificar}{C13} & \sinacc & \sinacc \\
\hline
\suj{SUP} & leer (solo estado) & \sinacc & \vb{modificar}{C13} & \sinacc & \vb{leer}{C15}, \vb{resolver}{C15} \\
\hline
\suj{GER} & leer, \vb{solicitar}{C7}, \vb{autorizar}{C8} & \vb{leer}{H} & leer, \vb{aprobar}{C12}, \vb{modificar}{C13} & leer & leer, \vb{resolver}{C15} \\
\hline
\suj{ADM} & \sinacc & \sinacc & leer, \vb{crear}{C12}, \vb{desactivar}{C12}, \vb{proponer}{C12}, \vb{modificar}{C13} & \sinacc & \sinacc \\
\hline
\suj{AUD} & \vb{leer}{C6}, \vb{exportar}{C18} & \vb{leer}{C6}, \vb{exportar}{C18} & \vb{leer}{C6}, \vb{modificar}{C13} & \vb{leer}{C6}, \vb{exportar}{C18} & \vb{leer}{C6}, \vb{exportar}{C18} \\
\hline
\suj{DUE} & \vb{leer}{C5} & \vb{leer}{C5} & \vb{modificar}{C13} & \sinacc & \sinacc \\
\hline
\suj{SIS} & leer & \dato{emitir C4}, \vb{revocar}{C4} & \sinacc & anexar & crear \\
\hline
\suj{OPE} & \sinacc & \sinacc & \sinacc & \sinacc & \sinacc \\
\hline
\suj{PUB} & \sinacc & \vb{comprobar}{C20} & \sinacc & \sinacc & \sinacc \\
\end{tabular}

\vspace{0.06in}
{\color{gris}Matriz real, objetos 6 a 10 de 10. \quad Toda decisión exige segundo factor reciente. \quad {\color{acento}\textbf{C17}}: nadie decide sobre una finca con cuyo dueño declaró un vínculo.}
}
\note{Dos cosas que se ven en esta mitad. Una: el verbo «eliminar» no aparece en enmiendas, certificados,
cuentas ni bitácora. Nada se borra; todo cambio es una versión nueva o un cambio de estado. Dos: «emitir»
solo está en la fila SIS. Ninguna fila de persona lo tiene.

Columna Enmiendas: solicitar (técnico, verificador o Gerencia), corregir (técnico), revisar
(verificador) y autorizar (Gerencia), con las reglas C7 a C10. Cuatro pasos, al menos tres personas.

Columna Cuentas y roles: C13 es «credenciales y declaración propias»; C12 es «nadie se administra a sí
mismo»: el administrador propone el rol y la Gerencia lo aprueba.

Columna Bitácora: solo se agrega. La leen la Gerencia y el auditor.

Fila PUB: quien tenga el código de un certificado comprueba si es auténtico y si sigue vigente, y nada
más (C20). Sale de la respuesta del cliente sobre autenticidad y estado del certificado (minuta, P-04).

C17, abstención por vínculo: quien declaró un vínculo familiar o comercial con un dueño no mide, no
aprueba, no asigna ni autoriza nada sobre esa finca (minuta, P-15).

Al lado de la matriz hay una tabla de exclusiones: cada persona tiene un solo rol, salvo la habilitación
temporal de un verificador como técnico.}
\end{frame}

%--------------------------- 18. TRES CELDAS -------------------------
\begin{frame}{Tres celdas bastan para leer la separación}
\quien{Julio}
\vspace{0.0in}
\begin{tabular}{@{}L{1.0in} L{5.9in} L{4.8in}@{}}
\toprule
 & \cabt{MEDICIONES} & \cabt{CERTIFICADOS} \\
\midrule
\textbf{TEC} & Crea, modifica y envía. Solo sus borradores, con el período abierto. & Sin acceso \\
\textbf{VER} & \dato{Aprueba solo si no registró ninguna versión.} & Lee \\
\textbf{SIS} & Lee & \dato{Emite solo si no hay otro vigente.} \\
\bottomrule
\end{tabular}

\vspace{0.25in}
\pie{Caso de prueba T-06: la verificadora que antes fue técnica intenta aprobar una medición que registró. Resultado esperado: deniega.}
\note{Son las reglas C2, C3 y C4. Cada regla de la matriz tiene un caso de prueba con resultado esperado:
treinta y siete casos, T-01 a T-37. El equipo que construya debe convertirlos en pruebas automáticas, más una por
cada celda sin acceso.

Otros casos para citar: T-07, segunda emisión sobre la misma parcela y período: no emite. T-08, el dueño
cambia el identificador en la dirección: deniega y responde «no encontrado». T-10, el auditor envía una
petición de modificación: deniega.

Caso nuevo para citar: T-17, después del cierre el verificador aprueba una medición que se envió a
tiempo: concede. El cliente pidió distinguir registrar, modificar, revisar y aprobar (minuta, P-06).

El mismo patrón, quien inicia no cierra, se repite en enmiendas (C7 a C10), períodos (C14),
asignaciones (C19), catálogo (C11) y cuentas (C12).}
\end{frame}

%--------------------------- 19. ARQUITECTURA ------------------------
\begin{frame}{Un solo punto decide, en el servidor,\\ antes de tocar el dato}
\quien{Axel}
\vspace{0.1in}
\begin{tikzpicture}[x=1in,y=1in,
  comp/.style={draw=regla,line width=0.8pt,fill=white,align=center,inner sep=3pt,minimum height=0.8in},
  pdp/.style={draw=acento,line width=1.6pt,fill=acentoclaro,align=center,inner sep=3pt,minimum height=0.8in},
  zona/.style={draw=gris,densely dashed,line width=0.7pt},
  tz/.style={font=\tCuatro\bfseries,text=gris,anchor=north west},
  dos/.style={{Stealth[length=8pt,width=7pt]}-{Stealth[length=8pt,width=7pt]},draw=gris,line width=1.1pt}]
\draw[zona] (0,0) rectangle (2.0,3.3);
\draw[zona] (2.35,0) rectangle (3.95,3.3);
\draw[zona] (4.3,0) rectangle (8.9,3.3);
\draw[zona] (9.25,0) rectangle (12.1,3.3);
\node[tz] at (0.08,3.25) {NO CONFIABLE};
\node[tz] at (2.43,3.25) {BORDE};
\node[tz] at (4.38,3.25) {APLICACIÓN};
\node[tz] at (9.33,3.25) {DATOS Y SECRETOS};
\node[comp,minimum width=1.6in] (nav) at (1.0,1.25) {Navegador};
\node[comp,minimum width=1.2in] (px) at (3.15,1.25) {Proxy};
\node[comp,minimum width=4.0in] (app) at (6.6,1.25) {\textbf{Aplicación}\\{\tCuatro identifica, valida y hace cumplir}};
\node[pdp,minimum width=4.0in] (opa) at (6.6,2.45) {\textbf{Motor de decisión}\\{\tCuatro la matriz, como política}};
\node[comp,minimum width=2.4in] (pg) at (10.68,2.1) {\textbf{Base de datos}\\{\tCuatro segundo candado}};
\node[comp,minimum width=2.4in] (bao) at (10.68,0.75) {\textbf{Bóveda}\\{\tCuatro secretos y sello}};
\draw[dos] (nav) -- (px);
\draw[dos] (px) -- (app);
\draw[dos,draw=acento] (app) -- (opa);
\draw[dos] ($(app.east)+(0,0.18)$) -- ++(0.3,0) |- (pg.west);
\draw[dos] ($(app.east)+(0,-0.18)$) -- ++(0.3,0) |- (bao.west);
\end{tikzpicture}

\vspace{0.2in}
\dato{Esconder un botón no es autorización.} Si el servidor no decide, cualquiera llama directo.
\note{Flujo de una petición: el proxy termina la conexión cifrada; la aplicación identifica la sesión,
valida la forma de los datos y lee de la base los hechos; entrega al motor el sujeto, la acción, el
objeto y esos hechos; el motor devuelve la decisión con la regla que la determinó. Solo entonces se toca
el dato, en una transacción que también escribe la bitácora. La base aplica sus propias restricciones.

Fronteras de confianza: nada de lo que llega del navegador se da por bueno; los hechos se toman de la
base, nunca de la petición; si el motor no responde, se deniega; la base no confía en la aplicación.

Si preguntan qué pasa si el control está solo en la interfaz: el técnico envía a mano la petición de
aprobar su propia medición y el servidor la ejecuta. Es el error de referencia directa a objetos de la
clase 13.

Hospedaje, que el cliente dejó a nuestro criterio (minuta, P-23): un servidor virtual contratado, operado
por personal que la empresa designe. Hay además una ruta pública de solo lectura para comprobar
certificados.}
\end{frame}

%--------------------------- 20. TECNOLOGIAS -------------------------
\begin{frame}{Seis tecnologías libres,\\ elegidas por lo que protegen}
\quien{Axel}
\vspace{-0.05in}
{\renewcommand{\arraystretch}{1.36}
\begin{tabular}{@{}L{2.9in} L{3.7in} L{5.1in}@{}}
\toprule
\cabt{DECISIÓN} & \cabt{ELECCIÓN} & \cabt{LO QUE APORTA} \\
\midrule
Lenguaje y framework & Python 3.14, Django 5.2 LTS & Consultas parametrizadas; escape por omisión \\
Base de datos & PostgreSQL 18 & Seguridad por fila; emisión única \\
Sesión & De servidor, cookie opaca & \dato{Se revoca en el acto} \\
Autorización & Open Policy Agent 1.21 & Un solo lugar; la matriz como política \\
Contenedores & Podman sobre AlmaLinux 10 & Sin privilegios; confinados por SELinux \\
\bottomrule
\end{tabular}}
\note{Para cada una el documento indica versión, alternativas descartadas, aporte, limitaciones y
herramientas de análisis.

Descartes por licencia: HashiCorp Vault (BUSL desde 2023) y Docker Desktop. Se usa OpenBao, que es
software libre.

Descartes por seguridad: tokens JWT, porque no se pueden revocar en el acto al dar de baja o cambiar
de rol; MariaDB, porque no tiene seguridad por fila; Docker Engine, por su servicio con privilegios; una
cadena de bloques, porque resuelve la desconfianza entre organizaciones, que aquí no existe.

Limitaciones que hay que saber decir: los permisos de Django son por tipo de objeto, no por objeto; su
sitio de administración no se instala; el superusuario de PostgreSQL se salta todo; el motor no conoce el
historial y no filtra listados.

Contraseñas: Argon2id con sal por usuario. Segundo factor: llave de acceso o códigos temporales; nunca SMS.}
\end{frame}

%--------------------------- 21. BITACORA ----------------------------
\begin{frame}{La bitácora no la reescribe\\ ni quien administra el servidor}
\quien{Axel}
\vspace{0.2in}
\begin{tikzpicture}[x=1in,y=1in,
  reg/.style={draw=regla,line width=0.8pt,fill=plano,minimum width=2.25in,minimum height=1.0in,
              align=center,inner sep=3pt},
  ext/.style={draw=acento,line width=1.6pt,fill=acentoclaro,minimum width=2.6in,minimum height=1.0in,
              align=center,inner sep=3pt},
  fl/.style={-{Stealth[length=9pt,width=8pt]},draw=gris,line width=1.2pt}]
\draw[draw=gris,densely dashed,line width=0.7pt] (-0.1,-0.8) rectangle (8.55,1.15);
\node[font=\tCuatro\bfseries,text=gris,anchor=north west] at (0,1.1) {SERVIDOR: SOLO SE PUEDE INSERTAR};
\node[reg] (a) at (1.25,-0.05) {Registro 1041\\{\tCuatro huella del 1040}};
\node[reg] (b) at (4.2,-0.05) {Registro 1042\\{\tCuatro huella del 1041}};
\node[reg] (c) at (7.15,-0.05) {Registro 1043\\{\tCuatro huella del 1042}};
\draw[fl] (b) -- (a); \draw[fl] (c) -- (b);
\node[ext,minimum width=2.25in] (t) at (10.97,-0.05) {\textbf{Testigo externo}\\{\tCuatro sello de tiempo y auditor}};
\draw[fl,draw=acento] (c) -- node[above,pos=0.66,font=\tCuatro,text=gris]{cada día} (t);
\end{tikzpicture}

\vspace{0.3in}
Cambiar un registro rompe la cadena. Reescribirla entera no coincide con lo que ya salió.
\note{Cuatro mecanismos que se apoyan: (1) la aplicación solo puede insertar, no modificar ni borrar;
(2) cada registro incluye la huella SHA-256 del anterior; (3) una vez al día la última huella sale del
servidor, recibe un sello de tiempo RFC 3161 y va al buzón de la Gerencia y del auditor; (4) la acción y
su registro se confirman en la misma transacción.

Qué se registra: autenticación, cambios, consultas de datos sensibles, administración y toda decisión
denegada. Con qué campos: instante, cuenta y rol, acción, objeto y versión, resultado y regla, origen,
huellas.

Conservación: diez años para cambios y decisiones, que es la trazabilidad que el cliente debe poder
reconstruir (minuta, P-18); dos para consultas (supuesto P-19).

Secretos: viven en una bóveda; la llave de sello no sale, se le pide a la bóveda que firme. La bóveda se
abre con dos de tres custodios.}
\end{frame}

%--------------------------- 22. RESIDUALES --------------------------
\begin{frame}{Lo que no se cubre está declarado}
\quien{Javier}
\vspace{0.0in}
\begin{tabular}{@{}L{5.0in} L{6.7in}@{}}
\toprule
\cabt{RIESGO RESIDUAL} & \cabt{POR QUÉ SE ACEPTA} \\
\midrule
Colusión entre técnico y verificador & Evitarla exige una tercera persona por certificado \\
Una medición falsa desde el origen & El sistema no ve el árbol \\
Extracción lenta, bajo los umbrales & Bajarlos castiga a quien trabaja bien \\
Un vínculo con un dueño que nadie declaró & No existe un registro que lo descarte \\
El histórico de las hojas de cálculo & No se le puede dar la trazabilidad que no tuvo \\
\bottomrule
\end{tabular}
\note{Son trece riesgos residuales en el documento de arquitectura. Otros: equipo de usuario infectado; el
operador puede detener o destruir (se detecta, no se impide); un solo servidor; sin hardware de
seguridad dedicado; la fecha de medición es declarada, porque no hay aplicación sin conexión; el
proveedor de hospedaje tiene acceso físico; la aprobación no lleva firma digital certificada; y la
misma regla vive en dos lugares, el motor y la base.

Si preguntan qué requerimiento quedó sin cubrir: la trazabilidad de diez años hacia atrás. Las hojas de
cálculo se congelan, se les calcula la huella y lo migrado queda marcado como histórico.

Antes de entrar en operación hay diez criterios de aceptación, entre ellos una prueba de penetración
autorizada por escrito.}
\end{frame}

%--------------------------- 23. RESPUESTAS DEL CLIENTE -------------
\begin{frame}{El cliente respondió,\\ y el diseño cambió en cinco puntos}
\quien{Javier}
\vspace{0.0in}
\begin{tabular}{@{}L{0.9in} L{5.0in} L{5.6in}@{}}
\toprule
 & \cabt{LO QUE RESPONDIÓ} & \cabt{QUÉ CAMBIÓ} \\
\midrule
P-01 & Un verificador puede medir, por excepción & Habilitación temporal; no aprueba lo que midió \\
P-05 & Las fincas las asignan los supervisores & \dato{Rol nuevo: supervisor de campo} \\
P-06 & La revisión sigue después del cierre & El cierre detiene registrar, no aprobar \\
P-07 & Revisión técnica antes de autorizar & Solicitar, corregir, revisar, autorizar \\
P-10 & Avisar primero; no bloquear de inmediato & Aviso, verificación, restricción temporal \\
\bottomrule
\end{tabular}

\vspace{0.2in}
\pie{Catorce de veintiséis consultas, respondidas por escrito. Las otras doce son supuestos declarados. Minuta: informe de políticas, Anexo A.}
\note{Qué decir: el diseño se hizo completo con supuestos declarados, y la consulta sirvió como validación.
De catorce respuestas: dos confirmaron el supuesto (P-03, el verificador devuelve y no edita; P-11, el
técnico solo ve sus fincas), cuatro lo ajustaron, seis lo cambiaron y dos quedaron a nuestro criterio.

Lo que resistió: la técnica de separación. La respuesta a P-01 no la cambió; solo volvió dinámica la
exclusión entre técnico y verificador.

Lo que hubo que rehacer: el orden de la enmienda (P-07), lo que el cierre detiene (P-06) y quién asigna
las fincas (P-05). Y el diseño era más estricto de lo que el cliente quería: bloqueaba de inmediato
donde bastaba avisar (P-10).

Otros cambios para citar si preguntan: sociedades y representantes como dueños, y declaración de
vínculos con abstención (P-15); ninguna norma de auditoría se da por supuesta, y los auditores exportan
evidencia sellada (P-18); registro después de la visita, con fecha de medición y fecha de registro
(P-22); diseño sin cantidades fijas (P-02).

Cada cambio llegó hasta la matriz: diecisiete casos de prueba nuevos, T-21 a T-37, y cinco modificados.}
\end{frame}

%--------------------------- 24. DECISIONES DELEGADAS ----------------
\begin{frame}{Cuatro decisiones que el cliente\\ dejó en nuestras manos}
\quien{Javier}
\vspace{0.0in}
\begin{tabular}{@{}L{4.7in} L{7.0in}@{}}
\toprule
\cabt{LA PREGUNTA} & \cabt{LO QUE DECIDIMOS} \\
\midrule
¿Revisa un verificador distinto del original? & \dato{Sí, cuando la enmienda es material} \\
¿Hace falta firma digital certificada? & No: sello de la empresa y comprobación en línea \\
¿Llaves físicas para todo el personal? & No: llave de acceso del propio equipo \\
¿Dónde se hospeda y quién lo opera? & Servidor virtual contratado; personal designado \\
\bottomrule
\end{tabular}

\vspace{0.2in}
\pie{El cliente pidió que cada decisión tenga una justificación clara y responda a un riesgo real, no el sistema más complejo.}
\note{Verificador distinto (minuta, P-07): tres opciones. Siempre distinto: toda corrección, hasta una de
forma, espera a otra persona. El mismo: quien aprobó el dato erróneo dictamina sobre su propio error.
Elegimos según lo que cambia: es material si cambia un dato del cálculo de toneladas o si ya hay un
certificado; entonces revisa otro. Lo determina el sistema, no una persona.

Firma digital (P-04): la firma certificada de cada verificador da autenticidad y autoría con presunción
legal, pero no dice si el certificado fue revocado después, y cuesta por persona. Elegimos sello de la
empresa, comprobación en línea por código y aprobación personal con segundo factor, que queda en la
bitácora.

Llaves físicas (P-23): no hay presupuesto para todos. La llave de acceso del propio dispositivo
(passkey) resiste igual el correo falso y no cuesta nada. Se recomiendan llaves físicas solo para
Gerencia, administrador, operador y custodios de la bóveda.

Hospedaje (P-23): servidor propio, servidor virtual contratado o servicios administrados, con los
cuatro criterios del cliente. El propio exige inversión y cuidar energía y enlace en una oficina
rural; los administrados no son software libre. Lo cotidiano está automatizado.}
\end{frame}

%--------------------------- 25. REFERENCIAS -------------------------
\begin{frame}{Referencias principales}
\quien{}
\vspace{0.0in}
{\tCuatro
\begin{columns}[T,onlytextwidth]
\begin{column}{0.485\textwidth}
\raggedright
Abadi, M., \& Fournet, C. (2003). Access control based on execution history. \emph{NDSS}.\\[5pt]
Clark, D. D., \& Wilson, D. R. (1987). A comparison of commercial and military computer security
policies. \emph{IEEE S\&P}.\\[5pt]
Kuhn, D. R., Coyne, E. J., \& Weil, T. R. (2010). Adding attributes to role-based access control.
\emph{Computer, 43}(6).\\[5pt]
Li, N., Tripunitara, M. V., \& Bizri, Z. (2007). On mutually exclusive roles and separation-of-duty.
\emph{ACM TISSEC, 10}(2).\\[5pt]
Saltzer, J. H., \& Schroeder, M. D. (1975). The protection of information in computer systems.
\emph{Proc. IEEE, 63}(9).\\[5pt]
Sandhu, R. S., Coyne, E. J., Feinstein, H. L., \& Youman, C. E. (1996). Role-based access control
models. \emph{Computer, 29}(2).
\end{column}
\begin{column}{0.485\textwidth}
\raggedright
Schneier, B., \& Kelsey, J. (1999). Secure audit logs to support computer forensics. \emph{ACM
TISSEC, 2}(2).\\[5pt]
Simon, R. T., \& Zurko, M. E. (1997). Separation of duty in role-based environments. \emph{CSFW}.\\[5pt]
Temoshok, D., et al. (2025). \emph{Digital identity guidelines: Authentication and authenticator
management} (NIST SP 800-63B-4).\\[5pt]
Ley N.º 8968 (2011) y su reglamento, Decreto N.º 37554-JP. Ley N.º 7975 (2000). Ley N.º 8454
(2005). Ley N.º 9048 (2012).\\[5pt]
ISO/IEC 17029:2019. ISO 14065:2020. ISO/IEC 27001:2022.\\[5pt]
Vargas Rivera, A. (2026). \emph{Material de estudio del curso CE-1115 y enunciado del proyecto,
parte 1}. ITCR.
\end{column}
\end{columns}
}
\note{La lista completa, en formato APA, está en cada uno de los cinco documentos.}
\end{frame}

\end{document}
"
%=====================================================================
\documentclass[aspectratio=169,t]{beamer}

\usepackage[utf8]{inputenc}
\usepackage[T1]{fontenc}
\usepackage[spanish,es-noshorthands,es-tabla]{babel}
\usepackage[default]{lato}
\usepackage{booktabs,array,colortbl}
\usepackage{tikz}
\usetikzlibrary{positioning,arrows.meta,calc}

\ifdefined\connotas\setbeameroption{show notes}\fi

%--------------------------- lamina y margenes -----------------------
\geometry{papersize={13.333in,7.5in}}
\setbeamersize{text margin left=0.6in,text margin right=0.6in}

%--------------------------- color -----------------------------------
\definecolor{acento}{HTML}{0F6E6A}
\definecolor{acentoclaro}{HTML}{E7F2F1}
\definecolor{tinta}{HTML}{1A1A1A}
\definecolor{gris}{HTML}{5F5E5A}
\definecolor{regla}{HTML}{B9B8B0}
\definecolor{plano}{HTML}{FAFAFA}
\definecolor{negado}{HTML}{F2F1EC}

%--------------------------- escala: cuatro tamanos ------------------
\newcommand{\tUno}{\fontsize{40}{45}\selectfont}
\newcommand{\tDos}{\fontsize{28}{35}\selectfont}
\newcommand{\tTres}{\fontsize{18}{25}\selectfont}
\newcommand{\tCuatro}{\fontsize{13}{17.5}\selectfont}
\makeatletter
\renewcommand\normalsize{\@setfontsize\normalsize{18}{25}}
\renewcommand\small{\@setfontsize\small{13}{17.5}}
\renewcommand\footnotesize{\@setfontsize\footnotesize{13}{17.5}}
\renewcommand\scriptsize{\@setfontsize\scriptsize{13}{17.5}}
\renewcommand\tiny{\@setfontsize\tiny{13}{17.5}}
\renewcommand\large{\@setfontsize\large{28}{35}}
\renewcommand\Large{\@setfontsize\Large{40}{45}}
\makeatother

%--------------------------- cromo -----------------------------------
\usetheme{default}
\setbeamertemplate{navigation symbols}{}
\setbeamercolor{structure}{fg=acento}
\setbeamercolor{normal text}{fg=tinta}
\setbeamercolor{background canvas}{bg=white}
\setbeamerfont{note page}{size=\fontsize{20}{27}\selectfont}

\setbeamertemplate{frametitle}{%
  \vskip0.0in%
  \begin{minipage}[t]{\textwidth}
    \raggedright\color{tinta}\tUno\bfseries\insertframetitle\par
  \end{minipage}\par
  \vskip0.12in}

\gdef\presentador{}
\newcommand{\quien}[1]{\gdef\presentador{#1}}
\setbeamertemplate{footline}{%
  \hbox to\paperwidth{\hskip0.6in{\tCuatro\color{gris}\insertframenumber}\hfill
    {\tCuatro\color{gris}\presentador}\hskip0.6in}%
  \vskip0.45in}

%--------------------------- atajos ----------------------------------
\newcommand{\dato}[1]{\textcolor{acento}{\textbf{#1}}}
\newcommand{\pie}[1]{\leavevmode{\tCuatro\color{gris}#1}}
\newcommand{\cabt}[1]{\leavevmode{\tCuatro\bfseries\color{gris}#1}}
\newcolumntype{L}[1]{>{\raggedright\arraybackslash}p{#1}}
\arrayrulecolor{regla}
\setlength{\arrayrulewidth}{0.8pt}
\renewcommand{\arraystretch}{1.5}
% celdas de la matriz
\newcommand{\vb}[2]{\mbox{#1\,{\color{acento}\bfseries #2}}}
\newcommand{\sinacc}{\cellcolor{negado}\leavevmode{\color{gris}---}}
\newcommand{\suj}[1]{\leavevmode{\tTres\bfseries #1}}

%=====================================================================
\begin{document}

%--------------------------- 1. PORTADA ------------------------------
\begin{frame}[plain]
\vspace*{1.25in}
{\tCuatro\color{gris}PROYECTO \ $\cdot$ \ PARTE 1 \ $\cdot$ \ DEFENSA ORAL}\\[0.3in]
{\tUno\bfseries\color{tinta}Diseño del control de acceso\\[4pt] para Verdeazul Carbono}\\[0.25in]
{\color{acento}\rule{1.6in}{4pt}}\\[0.25in]
{\tDos\color{gris}Políticas, matriz y arquitectura de seguridad}\\[0.7in]
{\tTres\bfseries Randall Bolaños \quad Javier Tenorio \quad Julio Varela \quad Axel Flores}\\[0.12in]
{\tCuatro\color{gris}Grupo 6 \ $\cdot$ \ CE-1115 Seguridad de la Información \ $\cdot$ \ Prof. MSc. Andrés
Vargas Rivera \ $\cdot$ \ ITCR \ $\cdot$ \ 15 de octubre de 2026}
\note{Reparto propuesto. Randall: cliente, activos y separación de responsabilidades. Javier:
amenazas, marco legal, fricción, riesgos residuales y respuestas del cliente. Julio: modelo de control
de acceso y matriz. Axel: controles de historial y arquitectura. Todos deben poder responder por cualquier lámina:
las preguntas del profesor son individuales.}
\end{frame}

%--------------------------- 2. EL NEGOCIO ---------------------------
\begin{frame}{Verdeazul vende confianza:\\ cada certificado respalda una tonelada real}
\quien{Randall}
\vspace{0.15in}
\begin{tikzpicture}[x=1in,y=1in,
  caja/.style={draw=regla,line width=0.8pt,fill=plano,minimum width=2.0in,minimum height=0.95in,
               align=center,font=\tTres,inner sep=2pt},
  fl/.style={-{Stealth[length=9pt,width=8pt]},draw=gris,line width=1.2pt}]
\node[caja] (a) at (1.00,0) {Finca\\forestal};
\node[caja] (b) at (3.53,0) {Técnico\\mide};
\node[caja] (c) at (6.06,0) {Verificador\\aprueba};
\node[caja,draw=acento,line width=1.6pt] (d) at (8.59,0) {\textbf{Certificado}\\de carbono};
\node[caja] (e) at (11.12,0) {Mercado\\internacional};
\foreach \p/\q in {a/b,b/c,c/d,d/e} \draw[fl] (\p) -- (\q);
\end{tikzpicture}

\vspace{0.4in}
{\tDos \dato{Un dato inflado es fraude.}}

\vspace{0.12in}
Auditores europeos pueden pedir la trazabilidad de cualquier certificado de los últimos diez años.
\note{Qué decir: Verdeazul certifica captura de carbono en fincas forestales. El técnico mide, el
verificador aprueba y se emite un certificado que el dueño vende afuera. El valor del certificado es la
confianza en el dato. Hoy todo vive en hojas de cálculo compartidas.

Si preguntan por qué empezamos por el negocio: porque el enunciado pide políticas que describan a este
cliente, y casi todas las decisiones salen de esta lámina.}
\end{frame}

%--------------------------- 3. LOS DOS SUSTOS -----------------------
\begin{frame}{Los dos sustos de la carta son fallas de integridad}
\quien{Randall}
\vspace{0.1in}
\begin{tabular}{@{}L{2.2in} L{9.5in}@{}}
\toprule
\cabt{PRIMER SUSTO} & Una medición cambió después de aprobada, y nadie sabe quién fue. \\
\midrule
\cabt{SEGUNDO SUSTO} & Un certificado se emitió dos veces sobre la misma hectárea. \\
\bottomrule
\end{tabular}

\vspace{0.45in}
{\tDos Orden de prioridad: \dato{integridad}, confidencialidad, disponibilidad.}

\vspace{0.12in}
Cuando dos pilares chocan, gana la integridad: sin registro en bitácora no hay acción.
\note{Qué decir: ninguno de los dos incidentes fue un robo de información. En los dos, el dato dejó de
ser confiable. Por eso la integridad va primero. La confidencialidad va segunda por la competencia
entre dueños y por los datos personales. La disponibilidad va tercera, con un matiz: importa más que el
dato exista en diez años que un día sin sistema.

Si preguntan por un ejemplo de «gana la integridad»: si la bitácora no se puede escribir, la operación
se rechaza aunque eso detenga el trabajo (POL-BT-06).}
\end{frame}

%--------------------------- 4. ACTIVOS ------------------------------
\begin{frame}{Lo que más vale es lo que sostiene un certificado}
\quien{Randall}
\vspace{0.05in}
\begin{tabular}{@{}L{4.6in} L{2.6in} L{4.4in}@{}}
\toprule
\cabt{ACTIVO} & \cabt{CLASIFICACIÓN} & \cabt{LO QUE NO PUEDE FALLAR} \\
\midrule
Mediciones y certificados & Confidencial & \dato{Integridad} \\
Catálogo de fincas: la cartera & Confidencial & Integridad y confidencialidad \\
Datos personales de los dueños & Confidencial & Confidencialidad \\
Bitácora y llave de sello & Restringido & \dato{Integridad} \\
\bottomrule
\end{tabular}

\vspace{0.3in}
\pie{El inventario completo tiene once activos, cada uno valorado en los tres pilares (informe de políticas, sección 2).}
\note{Qué decir: clasificamos en cuatro niveles (público, interno, confidencial, restringido) y valoramos
cada activo con la escala de FIPS 199. La Gerencia de Certificación es la responsable del dato; el
administrador y el operador son custodios.

Si preguntan por un activo que no es obvio: los parámetros de cálculo (A-10). Un factor alterado infla
todos los certificados a la vez. Y la llave de sello (A-09): quien la tenga fabrica certificados
creíbles.}
\end{frame}

%--------------------------- 5. AMENAZAS -----------------------------
\begin{frame}{\textcolor{acento}{7 de 11} amenazas son internas:\\ el atacante más probable ya tiene una cuenta}
\quien{Javier}
\vspace{0.05in}
\begin{tabular}{@{}L{1.4in} L{6.6in} L{3.6in}@{}}
\toprule
\cabt{ID} & \cabt{YA OCURRIÓ EN VERDEAZUL} & \cabt{PILAR} \\
\midrule
AM-03 & Una medición editada después de aprobarse & Integridad \\
AM-06 & Un certificado emitido dos veces & Integridad \\
AM-04 & Cientos de fincas consultadas en un fin de semana & Confidencialidad \\
\bottomrule
\end{tabular}

\vspace{0.3in}
El diseño no mira el perímetro: mira qué puede hacer cada persona autorizada.
\note{Qué decir: el modelo tiene once amenazas, cada una con agente, motivo, activo, la debilidad que
hoy lo permite y el riesgo. Siete son internas: técnico que infla, misma persona que mide y aprueba,
edición silenciosa, extracción de la cartera, administrador u operador que abusa, doble emisión por
error y, desde la consulta al cliente, el empleado con un vínculo con un dueño (AM-11): no existe un
registro de esos vínculos. Tres ya ocurrieron.

Externas: dueño que quiere ver a un competidor, robo de credenciales, ransomware y certificado
falsificado o cuenta de auditor tomada.

Si preguntan cómo calculamos el riesgo: probabilidad por impacto. Probabilidad alta es que ya ocurrió
aquí o que es la vía habitual de ataque en el país.}
\end{frame}

%--------------------------- 6. MARCO LEGAL --------------------------
\begin{frame}{Cada obligación legal termina\\ en una regla que se puede comprobar}
\quien{Javier}
\vspace{0.05in}
\begin{tabular}{@{}L{2.9in} L{4.2in} L{4.5in}@{}}
\toprule
\cabt{NORMA} & \cabt{EXIGE} & \cabt{REGLA EN VERDEAZUL} \\
\midrule
Ley 8968, art. 14 & Autorización para transferir datos & El auditor consulta en la plataforma \\
Reglamento, art. 38 & Avisar una fuga en 5 días hábiles & Procedimiento con ese plazo \\
Ley 7975, art. 2 & Medidas razonables de secreto & Acceso por necesidad y registro \\
ISO/IEC 17029 & Quien revisa no participó & \dato{Medir y aprobar, separados} \\
\bottomrule
\end{tabular}

\vspace{0.2in}
\pie{El cliente aclaró que no audita bajo una norma determinada: la ISO se cita como referencia. Su exigencia firme es la trazabilidad de diez años.}
\note{Qué decir: la Ley 8968 aplica porque los datos de los dueños se comparten con auditores externos;
no es una base de uso exclusivamente interno. Derechos del titular en cinco días hábiles; consentimiento
informado al dar de alta a un dueño; conservación hasta diez años y después se desasocia.

La Ley 7975 es la que permite reclamar si la cartera llega a la competencia, pero solo si la empresa
puede demostrar que tomó medidas razonables: por eso el acceso por asignación y la bitácora tienen
además valor legal.

También: Código Penal (arts. 196 bis, 217 bis, 229 bis y 231), Código de Trabajo art. 71 g) y Ley 8454
de firma digital.

Auditores, según la minuta (P-18): una vez al año, varios días, en sitio o a distancia. Pueden exportar
evidencia, pero como paquetes que genera el sistema, sellados, registrados y sin datos personales. El
acceso vence con la ventana autorizada; prorrogarlo es una autorización nueva.}
\end{frame}

%--------------------------- 7. MODELO: TRES CAPAS -------------------
\begin{frame}{Roles con restricciones,\\ y condiciones que solo pueden quitar}
\quien{Julio}
\vspace{0.2in}
\begin{tikzpicture}[x=1in,y=1in,
  capa/.style={draw=acento,line width=1.4pt,fill=acentoclaro,minimum width=2.4in,minimum height=1.1in,
               align=center,inner sep=3pt},
  ext/.style={draw=regla,line width=0.8pt,fill=plano,minimum height=1.1in,align=center,inner sep=3pt},
  fl/.style={-{Stealth[length=9pt,width=8pt]},draw=gris,line width=1.2pt},
  no/.style={-{Stealth[length=9pt,width=8pt]},draw=gris,line width=1.0pt,densely dashed}]
\node[ext,minimum width=1.6in] (p) at (0.8,0) {Petición\\autenticada};
\node[capa] (r) at (3.25,0) {\textbf{1 $\cdot$ Roles}\\{\tCuatro ¿Este rol puede?}};
\node[capa] (c) at (6.10,0) {\textbf{2 $\cdot$ Condiciones}\\{\tCuatro ¿Aquí y ahora?}};
\node[capa] (h) at (8.95,0) {\textbf{3 $\cdot$ Historial}\\{\tCuatro ¿Encaja con lo previo?}};
\node[ext,minimum width=1.5in] (d) at (11.35,0) {Conceder};
\draw[fl] (p) -- (r);
\draw[fl] (r) -- node[above,font=\tCuatro,text=gris]{sí} (c);
\draw[fl] (c) -- node[above,font=\tCuatro,text=gris]{sí} (h);
\draw[fl] (h) -- node[above,font=\tCuatro,text=gris]{sí} (d);
\node[draw=regla,line width=0.8pt,fill=negado,minimum width=8.1in,minimum height=0.6in] (n) at (6.10,-1.55)
  {Denegar y registrar en bitácora};
\foreach \k in {r,c,h} \draw[no] (\k.south) -- node[right,font=\tCuatro,text=gris]{no} (\k.south |- n.north);
\end{tikzpicture}

\vspace{0.3in}
\dato{Solo la primera capa concede.} La matriz muestra lo máximo que cada rol puede llegar a hacer.
\note{Qué decir: elegimos RBAC con restricciones, nivel RBAC2 de Sandhu. Sobre esa base, un conjunto
cerrado de atributos decide si el permiso aplica a esta finca, en este estado y en este momento.
Es el enfoque centrado en roles de Kuhn, Coyne y Weil.

La frontera del híbrido: la capa 1 dice si el rol tiene el verbo; la capa 2 si aplica a este objeto ahora;
la capa 3 si encaja con el historial. Se combinan sin contradecirse porque solo la tabla de roles concede
y las otras dos son filtros: una negativa en cualquier capa gana.

Si preguntan por qué RBAC: la carta describe funciones, no personas; la gente rota; y hay que poder
explicarlo a un auditor leyendo una fila.}
\end{frame}

%--------------------------- 8. DESCARTES ----------------------------
\begin{frame}{Los otros tres modelos fallan\\ por razones de este caso}
\quien{Julio}
\vspace{0.05in}
\begin{tabular}{@{}L{2.2in} L{9.5in}@{}}
\toprule
\cabt{MODELO} & \cabt{POR QUÉ NO EN VERDEAZUL} \\
\midrule
DAC & Es lo que hay hoy: las hojas compartidas donde ocurrieron los dos sustos. \\
MAC & No hay niveles de secreto. Haría falta un compartimento por finca. \\
ABAC puro & Desaparece la fila que responde «¿qué puede hacer un verificador?». \\
\bottomrule
\end{tabular}

\vspace{0.3in}
\pie{De cada uno se conserva algo: la titularidad como atributo (DAC), la política que el usuario no puede
alterar (MAC) y las condiciones acotadas (ABAC).}
\note{DAC: el técnico que crea la medición sería su dueño y podría seguir modificándola o compartirla.
Es el primer susto.

MAC: Bell-LaPadula ordena por niveles de secreto y protege la confidencialidad. Aquí el riesgo principal
es de integridad. La variante de integridad, Biba, pone una etiqueta fija, pero la confianza en una
medición cambia con su estado y depende de quién hizo el paso anterior. Es el argumento de Clark y
Wilson. MAC sí lo usamos en la plataforma: SELinux confina los contenedores.

ABAC puro: podría expresarlo todo, pero no se puede revisar leyendo una fila, y la empresa no tiene
quién gobierne un catálogo de atributos.}
\end{frame}

%--------------------------- 9. NIVEL Y EXPLOSION --------------------
\begin{frame}{\textcolor{acento}{Siete} roles hoy, y siete con el doble de fincas}
\quien{Julio}
\vspace{0.25in}
{\tDos El alcance nunca es un rol: es un atributo.}

\vspace{0.35in}
\begin{tabular}{@{}L{3.3in} L{8.4in}@{}}
\toprule
RBAC$_2$, sin jerarquía & Si el verificador heredara del técnico, podría registrar: se rompe la separación. \\
Explosión de roles & Aparece al usar un rol para decir dónde o cuándo. Aquí eso vive en las condiciones. \\
La consulta lo probó & Cinco necesidades nuevas: \dato{un rol y cuatro atributos}. \\
\bottomrule
\end{tabular}
\note{Los siete roles: técnico, verificador, supervisor de campo, gerencia de certificación,
administrador de catálogo y cuentas, auditor externo y dueño de finca. Más una identidad de servicio
(SIS), el operador de plataforma, que no tiene rol en la aplicación, y el público sin autenticar.

El séptimo rol salió de la consulta: el cliente dijo que las fincas las asignan los supervisores del
personal de campo (minuta, P-05). Las otras cuatro necesidades nuevas (el verificador que mide por
excepción, el técnico que apoya a un compañero, la sociedad con un representante y el empleado con un
vínculo con un dueño) no son funciones de trabajo: se resolvieron con atributos.

Por qué RBAC2 y no RBAC3: no hay cadena de mando entre las siete funciones; la herencia vuelve
implícitos los permisos y la matriz no debe tener nada implícito; y con siete roles no ahorra nada.

Restricciones: un rol por persona, una identidad por persona, al menos dos verificadores, cuenta de
auditor con ventana. La pareja técnico--verificador es la única excepción: es dinámica. Un verificador
puede recibir una habilitación temporal de técnico, con fecha de fin, y nunca aprueba lo que midió
(minuta, P-01).

Si preguntan dónde está el riesgo de explosión: un rol por finca para aislar a los dueños, o roles como
«Técnico Zona Sur 2027».}
\end{frame}

%--------------------------- 10. HBAC --------------------------------
\begin{frame}{Cuatro controles de historial, uno de cada tipo}
\quien{Axel}
\vspace{0.0in}
\begin{tabular}{@{}L{0.7in} L{1.7in} L{4.6in} L{4.4in}@{}}
\toprule
 & \cabt{TIPO} & \cabt{QUÉ LO DISPARA} & \cabt{QUÉ DECIDE} \\
\midrule
H1 & Frecuencia & Muchas fincas distintas en un día & \dato{Tres niveles}: avisa, confirma, restringe \\
H2 & Tiempo & Una decisión fuera de jornada & Pide segundo factor y motivo \\
H3 & Combinada & Fuera de jornada, muchas fincas, casi todas nuevas & Confirma y avisa; si insiste, restringe \\
H4 & Secuencia & Quien participó antes quiere decidir & \dato{Deniega} \\
\bottomrule
\end{tabular}

\vspace{0.2in}
\pie{Ninguna señal aislada restringe: el cliente pidió avisar primero y no bloquear de inmediato (minuta, P-10).}
\note{Son reglas explícitas, no aprendizaje automático: se le pueden leer a un auditor y se pueden probar.

Cómo se distingue lo legítimo: se cuentan señales: volumen, novedad, momento e insistencia. El trabajo
extraordinario legítimo muestra una; la extracción de la cartera las muestra juntas.

Tres niveles: aviso (el usuario no nota nada), verificación adicional (segundo factor y motivo) y
restricción temporal (no abre fincas nuevas hasta el siguiente día hábil; la cuenta sigue activa).

H1 mitiga AM-04. Umbrales por rol, calibrados con datos reales; iniciales 40, 60 y 80 fincas por día
(supuesto P-13). Los intentos denegados cuentan doble.

H2 mitiga AM-08, la cuenta robada usada de madrugada. H4 es la separación de responsabilidades como
control de historial: reglas C3, C8 y C10.

H3 es la firma del incidente: fin de semana, cientos de fincas, ninguna suya. Pide confirmación en la
finca once y avisa; restringe en la veintiuna.

Falsos positivos: verificador con la cola acumulada al cierre, o nuevo y sin historial. El primer
período corre en observación para calibrar.

Descartada: Muralla China. Todos los dueños compiten: cada empleado quedaría con una sola finca.}
\end{frame}

%--------------------------- 11. EL TECNICO --------------------------
\begin{frame}{El técnico del fin de semana\\ hoy no vería ni una finca ajena}
\quien{Axel}
\vspace{0.05in}
\begin{tabular}{@{}L{4.6in} L{7.1in}@{}}
\toprule
\cabt{LO QUE HACE} & \cabt{LO QUE HACE EL SISTEMA} \\
\midrule
Abre una finca que no tiene asignada & \dato{Deniega}: su rol no ve fincas ajenas \\
Insiste tres, cinco, diez veces & Avisa a su supervisor; pide segundo factor; lo detiene media hora \\
La necesitaba de verdad & El supervisor le da una asignación temporal, que vence sola \\
Presenta su renuncia & Reporte de sus consultas de 90 días \\
\bottomrule
\end{tabular}

\vspace{0.25in}
\pie{El cliente lo confirmó: fincas asignadas con flexibilidad, y respuesta según la gravedad (minuta, P-10 y P-11).}
\note{Qué decir: la carta dice «no sabemos qué hacer con eso, pero nos gustaría que el sistema al menos se
diera cuenta». Nuestra primera respuesta no es vigilar el permiso: es quitarlo. Un técnico solo ve las
fincas que tiene asignadas (regla C1).

El cliente lo confirmó, con una condición: flexibilidad para apoyar a un compañero, sustituir a un
ausente o atender una visita extraordinaria, sin acceso permanente a todas las fincas. Se resuelve con
asignaciones temporales que crea el supervisor de campo, con motivo y fecha de fin (regla C19).

El diseño original suspendía la sesión a la quinta denegación. El cliente pidió no bloquear de
inmediato, y se cambió: aviso al supervisor a la tercera, segundo factor a la quinta, media hora de
restricción a la décima.

Los controles de historial quedan para quienes sí necesitan leer muchas fincas: verificador, supervisor,
gerencia, administrador y auditor. Si el técnico hubiera sido verificador, H3 le pide confirmación en la
finca once y lo restringe en la veintiuna.

Además hay una vía legal: confidencialidad firmada, que persiste después de la salida (Ley 8968 art. 11
y Ley 7975).}
\end{frame}

%--------------------------- 12. NADIE EMITE -------------------------
\begin{frame}{Nadie emite un certificado: lo hace el sistema}
\quien{Randall}
\vspace{0.25in}
\begin{tikzpicture}[x=1in,y=1in,
  paso/.style={draw=acento,line width=1.4pt,fill=acentoclaro,minimum width=3.3in,minimum height=1.1in,
               align=center,inner sep=3pt},
  sis/.style={draw=gris,line width=1.2pt,fill=negado,minimum width=3.3in,minimum height=1.1in,
              align=center,inner sep=3pt},
  fl/.style={-{Stealth[length=9pt,width=8pt]},draw=gris,line width=1.2pt},
  nota/.style={font=\tCuatro,text=gris,align=center}]
\node[paso] (a) at (1.65,0) {\textbf{Técnico}\\registra y envía};
\node[paso] (b) at (6.05,0) {\textbf{Verificador}\\aprueba o devuelve};
\node[sis]  (c) at (10.45,0) {\textbf{Sistema}\\emite};
\draw[fl] (a) -- (b); \draw[fl] (b) -- (c);
\node[nota,below=0.08in of a] {solo en fincas asignadas\\y con el período abierto};
\node[nota,below=0.08in of b] {no registró ninguna versión\\de esa medición};
\node[nota,below=0.08in of c] {no hay otro certificado vigente\\para la parcela y el período};
\end{tikzpicture}

\vspace{0.35in}
Ninguna persona tiene el permiso de emitir. \dato{Hacen falta dos, y distintas.}
\note{La acción crítica: emitir un certificado. Si una sola persona pudiera medir y aprobar, podría
inflar la biomasa, aprobar su propio dato y obtener un certificado por toneladas que no existen.

La segunda acción crítica: corregir una medición ya congelada. Es el primer susto, y se trata igual,
con el orden que pidió el cliente (minuta, P-07): solicita el técnico responsable, un verificador o la
Gerencia; corrige el técnico; revisa un verificador que no es autor ni solicitante; y la Gerencia da la
autorización final. Son al menos tres personas.

Si preguntan qué pasa con el certificado tras una enmienda (minuta, P-08): si las toneladas no cambian,
sigue vigente con la enmienda anotada. Si cambian, el sistema lo revoca y emite el reemplazo en un solo
acto, enlazados: nunca hay dos vigentes y ninguno se borra. Si ya se vendió, la Gerencia lo registra y
coordina con las partes.}
\end{frame}

%--------------------------- 13. POR MEDICION ------------------------
\begin{frame}{La separación se comprueba\\ por medición, no por puesto}
\quien{Randall}
\vspace{0.05in}
\begin{tabular}{@{}L{3.4in} L{8.3in}@{}}
\toprule
\cabt{SI SE COMPRUEBA\ldots} & \cabt{UN TÉCNICO ASCENDIDO A VERIFICADOR\ldots} \\
\midrule
por el puesto de hoy & aprueba lo que él mismo midió el año anterior. \\
\dato{por la medición} & \dato{no puede: figura como autor.} \\
\bottomrule
\end{tabular}

\vspace{0.35in}
Técnica: \textbf{separación por etapas del flujo}, reforzada con roles mutuamente excluyentes.
\note{Qué decir: quien inicia el proceso no puede cerrarlo. El sistema no pregunta qué rol tiene la
persona, pregunta si participó en esta medición. En la literatura es la separación basada en el objeto
o en el historial (Nash y Poland 1990; Simon y Zurko 1997).

Por qué es la adecuada: apunta al riesgo tal como lo dijo el cliente; sobrevive a los cambios de puesto
(la lección de Kerviel); no depende de que los roles estén bien administrados; y deja dos nombres
distintos en cada certificado.

El cliente confirmó que algunos verificadores fueron técnicos y que por excepción pueden medir, y pidió
«garantizar una revisión independiente, sin impedir que el personal pueda asumir otras
responsabilidades» (minuta, P-01). Una separación por puesto obligaría a escoger; esta da las dos
cosas: el verificador mide con una habilitación temporal y no podrá aprobar nunca lo que midió.

Coincide con el principio de imparcialidad de ISO/IEC 17029, que citamos solo como referencia.}
\end{frame}

%--------------------------- 14. ALTERNATIVAS ------------------------
\begin{frame}{Las alternativas se descartaron\\ por lo que cuestan aquí}
\quien{Randall}
\vspace{0.0in}
\begin{tabular}{@{}L{4.3in} L{7.4in}@{}}
\toprule
\cabt{TÉCNICA} & \cabt{POR QUÉ NO ES LA PRINCIPAL} \\
\midrule
Roles excluyentes, solos & No cubren a quien cambia de puesto \\
Doble control (cuatro ojos) & \dato{Duplica la verificación de cada certificado} \\
Segregación temporal & Limita cuándo, no quién \\
Rotación & No impide un acto; se recomienda contra la colusión \\
\bottomrule
\end{tabular}

\vspace{0.2in}
\pie{La segregación temporal sí se usa donde el problema es la duración: la asignación temporal, la habilitación del verificador, el plazo de la enmienda y la cuenta del auditor.}
\note{Roles mutuamente excluyentes estáticos (RBAC2): se conservan como refuerzo. Li, Tripunitara y Bizri
muestran que excluir roles es un mecanismo para aproximar la separación, no la separación misma.

Exclusión dinámica: basta cerrar sesión y abrir otra. Sola, demora el fraude y no lo impide. Se usa
como refuerzo en la pareja técnico--verificador, detrás de la regla C3.

Doble control: era la opción más fuerte contra el fraude. Se descartó por costo operativo: duplica la
verificación de cada certificado, con cualquier tamaño de equipo.

Rotación: cuesta conocimiento de la finca y depende de cuántas personas haya, un dato que el cliente
pidió no fijar (minuta, P-02).}
\end{frame}

%--------------------------- 15. FRICCION ----------------------------
\begin{frame}{La fricción existe y está acotada}
\quien{Javier}
\vspace{0.0in}
\begin{tabular}{@{}L{5.2in} L{6.5in}@{}}
\toprule
\cabt{LO QUE CUESTA} & \cabt{CÓMO SE MANEJA} \\
\midrule
Siempre hacen falta dos personas & Mínimo dos verificadores; \dato{sin atajo de emergencia} \\
Las devoluciones quitan tiempo de campo & Validación de rangos al registrar \\
Una enmienda pide tres personas & Plazo de diez días; el sistema la clasifica \\
La Gerencia concentra las segundas firmas & Nunca inicia lo que aprueba \\
\bottomrule
\end{tabular}
\note{Qué decir: no existe un atajo para aprobar lo propio. Un atajo así sería el hueco que el auditor
buscaría. Si el único verificador presente participó, la medición espera; la cola avisa a la Gerencia.

Lo que la técnica no resuelve: la colusión. Si un técnico y un verificador se ponen de acuerdo, el flujo
se cumple. La separación sube el costo del fraude de una persona a dos. Quedan la rotación, un indicador
de parejas repetidas y la auditoría externa.

La Gerencia es un cuello de botella y un punto de confianza: firma en varios flujos, pero nunca inicia
lo que aprueba, no registra ni aprueba mediciones, ya no asigna fincas y no resuelve sus propias alertas.

Enmienda: revisa un verificador distinto del original solo cuando es material (cambia toneladas o ya hay
certificado); lo determina el sistema. Con solo dos verificadores, si uno solicita sobre lo que aprobó
el otro no queda revisor: la presenta el técnico o la Gerencia.

La flexibilidad también cuesta un trámite: la asignación temporal la crea el supervisor y la
habilitación de un verificador para medir pasa por el administrador y la Gerencia.}
\end{frame}

%--------------------------- 16. MATRIZ (1 de 2) ---------------------
\begin{frame}{Lo que no está escrito está denegado}
\quien{Julio}
\vspace{-0.02in}
{\tCuatro
\renewcommand{\arraystretch}{1.3}
\setlength{\tabcolsep}{5pt}
\begin{tabular}{@{}L{0.75in}|L{1.95in}|L{2.25in}|L{2.0in}|L{1.85in}|L{2.3in}@{}}
\rowcolor{acentoclaro}
 & \textbf{Fincas y parcelas} & \textbf{Datos personales del dueño} & \textbf{Períodos de medición} & \textbf{Asignaciones} & \textbf{Mediciones} \\
\hline
\suj{TEC} & \vb{leer}{C1} & \vb{leer}{C1} (solo contacto) & leer & \vb{leer}{C1} (las propias) & \vb{leer}{C1}, \vb{crear}{C2}, \vb{modificar}{C2}, \vb{enviar}{C2} \\
\hline
\suj{VER} & \vb{leer}{H} & \sinacc & leer & \sinacc & \vb{leer}{H}, \vb{aprobar}{C3}, \vb{devolver}{C3} \\
\hline
\suj{SUP} & \vb{leer}{H} & \sinacc & leer & leer, \vb{crear}{C19}, \vb{modificar}{C19} & \vb{leer}{H} (solo estado y fechas) \\
\hline
\suj{GER} & \vb{leer}{H}, \vb{aprobar}{C11} & \vb{leer}{H} & leer, \vb{abrir}{C14}, \vb{cerrar}{C14} & leer & \vb{leer}{H} \\
\hline
\suj{ADM} & \vb{leer}{H}, \vb{proponer}{C11} & \vb{leer}{H}, crear, modificar, \vb{eliminar}{C16} & leer, crear, \vb{proponer}{C14} & \sinacc & \sinacc \\
\hline
\suj{AUD} & \vb{leer}{C6}, \vb{exportar}{C18} & \vb{leer}{C6} & \vb{leer}{C6}, \vb{exportar}{C18} & \vb{leer}{C6}, \vb{exportar}{C18} & \vb{leer}{C6}, \vb{exportar}{C18} \\
\hline
\suj{DUE} & \vb{leer}{C5} & \vb{leer}{C5} & leer & \vb{leer}{C5} & \vb{leer}{C5} (solo aprobadas) \\
\hline
\suj{SIS} & leer & leer (solo el nombre del titular) & leer & leer & leer \\
\hline
\suj{OPE} & \sinacc & \sinacc & \sinacc & \sinacc & \sinacc \\
\hline
\suj{PUB} & \sinacc & \sinacc & \sinacc & \sinacc & \sinacc \\
\end{tabular}

\vspace{0.06in}
{\color{gris}Matriz real, objetos 1 a 5 de 10. \quad --- sin acceso \quad {\color{acento}\textbf{C}$n$} permiso sujeto a la
regla $n$ \quad {\color{acento}\textbf{H}} lectura vigilada por historial}
}
\note{Cómo leerla: sujetos en las filas, objetos en las columnas. Una celda gris es sin acceso. Una marca C
significa que el permiso solo vale si se cumple esa regla; las condiciones solo pueden quitar.

Recorrido sugerido: fila TEC, columna Mediciones: crea, modifica y envía, pero no aprueba. Fila VER:
aprueba y devuelve con C3, no crea ni modifica, y no ve los datos personales del dueño. Fila SUP, que
es nueva: asigna fincas (C19) y de las mediciones solo ve estado y fechas. Fila GER: abre y cierra
períodos que registró el administrador (C14); ya no crea asignaciones. Fila AUD: leer y exportar
evidencia, con ventana (C6 y C18), y ningún verbo de escritura. Fila DUE: solo lo suyo (C5). Filas OPE
y PUB: nada en esta mitad.

Diez sujetos: siete roles, el sistema, el operador y el público sin autenticar.

Mínimo privilegio: cada permiso amplio tiene su justificación en la sección 6 del documento de la matriz.}
\end{frame}

%--------------------------- 17. MATRIZ (2 de 2) ---------------------
\begin{frame}{Nadie tiene «eliminar», y nadie emite}
\quien{Julio}
\vspace{-0.02in}
{\tCuatro
\renewcommand{\arraystretch}{1.3}
\setlength{\tabcolsep}{5pt}
\begin{tabular}{@{}L{0.75in}|L{2.3in}|L{1.9in}|L{2.75in}|L{1.4in}|L{2.0in}@{}}
\rowcolor{acentoclaro}
 & \textbf{Enmiendas} & \textbf{Certificados} & \textbf{Cuentas y roles} & \textbf{Bitácora} & \textbf{Alertas de historial} \\
\hline
\suj{TEC} & \vb{leer}{C1}, \vb{solicitar}{C7}, \vb{corregir}{C9} & \sinacc & \vb{modificar}{C13} & \sinacc & \sinacc \\
\hline
\suj{VER} & leer, \vb{solicitar}{C7}, \vb{revisar}{C10} & \vb{leer}{H} & \vb{modificar}{C13} & \sinacc & \sinacc \\
\hline
\suj{SUP} & leer (solo estado) & \sinacc & \vb{modificar}{C13} & \sinacc & \vb{leer}{C15}, \vb{resolver}{C15} \\
\hline
\suj{GER} & leer, \vb{solicitar}{C7}, \vb{autorizar}{C8} & \vb{leer}{H} & leer, \vb{aprobar}{C12}, \vb{modificar}{C13} & leer & leer, \vb{resolver}{C15} \\
\hline
\suj{ADM} & \sinacc & \sinacc & leer, \vb{crear}{C12}, \vb{desactivar}{C12}, \vb{proponer}{C12}, \vb{modificar}{C13} & \sinacc & \sinacc \\
\hline
\suj{AUD} & \vb{leer}{C6}, \vb{exportar}{C18} & \vb{leer}{C6}, \vb{exportar}{C18} & \vb{leer}{C6}, \vb{modificar}{C13} & \vb{leer}{C6}, \vb{exportar}{C18} & \vb{leer}{C6}, \vb{exportar}{C18} \\
\hline
\suj{DUE} & \vb{leer}{C5} & \vb{leer}{C5} & \vb{modificar}{C13} & \sinacc & \sinacc \\
\hline
\suj{SIS} & leer & \dato{emitir C4}, \vb{revocar}{C4} & \sinacc & anexar & crear \\
\hline
\suj{OPE} & \sinacc & \sinacc & \sinacc & \sinacc & \sinacc \\
\hline
\suj{PUB} & \sinacc & \vb{comprobar}{C20} & \sinacc & \sinacc & \sinacc \\
\end{tabular}

\vspace{0.06in}
{\color{gris}Matriz real, objetos 6 a 10 de 10. \quad Toda decisión exige segundo factor reciente. \quad {\color{acento}\textbf{C17}}: nadie decide sobre una finca con cuyo dueño declaró un vínculo.}
}
\note{Dos cosas que se ven en esta mitad. Una: el verbo «eliminar» no aparece en enmiendas, certificados,
cuentas ni bitácora. Nada se borra; todo cambio es una versión nueva o un cambio de estado. Dos: «emitir»
solo está en la fila SIS. Ninguna fila de persona lo tiene.

Columna Enmiendas: solicitar (técnico, verificador o Gerencia), corregir (técnico), revisar
(verificador) y autorizar (Gerencia), con las reglas C7 a C10. Cuatro pasos, al menos tres personas.

Columna Cuentas y roles: C13 es «credenciales y declaración propias»; C12 es «nadie se administra a sí
mismo»: el administrador propone el rol y la Gerencia lo aprueba.

Columna Bitácora: solo se agrega. La leen la Gerencia y el auditor.

Fila PUB: quien tenga el código de un certificado comprueba si es auténtico y si sigue vigente, y nada
más (C20). Sale de la respuesta del cliente sobre autenticidad y estado del certificado (minuta, P-04).

C17, abstención por vínculo: quien declaró un vínculo familiar o comercial con un dueño no mide, no
aprueba, no asigna ni autoriza nada sobre esa finca (minuta, P-15).

Al lado de la matriz hay una tabla de exclusiones: cada persona tiene un solo rol, salvo la habilitación
temporal de un verificador como técnico.}
\end{frame}

%--------------------------- 18. TRES CELDAS -------------------------
\begin{frame}{Tres celdas bastan para leer la separación}
\quien{Julio}
\vspace{0.0in}
\begin{tabular}{@{}L{1.0in} L{5.9in} L{4.8in}@{}}
\toprule
 & \cabt{MEDICIONES} & \cabt{CERTIFICADOS} \\
\midrule
\textbf{TEC} & Crea, modifica y envía. Solo sus borradores, con el período abierto. & Sin acceso \\
\textbf{VER} & \dato{Aprueba solo si no registró ninguna versión.} & Lee \\
\textbf{SIS} & Lee & \dato{Emite solo si no hay otro vigente.} \\
\bottomrule
\end{tabular}

\vspace{0.25in}
\pie{Caso de prueba T-06: la verificadora que antes fue técnica intenta aprobar una medición que registró. Resultado esperado: deniega.}
\note{Son las reglas C2, C3 y C4. Cada regla de la matriz tiene un caso de prueba con resultado esperado:
treinta y siete casos, T-01 a T-37. El equipo que construya debe convertirlos en pruebas automáticas, más una por
cada celda sin acceso.

Otros casos para citar: T-07, segunda emisión sobre la misma parcela y período: no emite. T-08, el dueño
cambia el identificador en la dirección: deniega y responde «no encontrado». T-10, el auditor envía una
petición de modificación: deniega.

Caso nuevo para citar: T-17, después del cierre el verificador aprueba una medición que se envió a
tiempo: concede. El cliente pidió distinguir registrar, modificar, revisar y aprobar (minuta, P-06).

El mismo patrón, quien inicia no cierra, se repite en enmiendas (C7 a C10), períodos (C14),
asignaciones (C19), catálogo (C11) y cuentas (C12).}
\end{frame}

%--------------------------- 19. ARQUITECTURA ------------------------
\begin{frame}{Un solo punto decide, en el servidor,\\ antes de tocar el dato}
\quien{Axel}
\vspace{0.1in}
\begin{tikzpicture}[x=1in,y=1in,
  comp/.style={draw=regla,line width=0.8pt,fill=white,align=center,inner sep=3pt,minimum height=0.8in},
  pdp/.style={draw=acento,line width=1.6pt,fill=acentoclaro,align=center,inner sep=3pt,minimum height=0.8in},
  zona/.style={draw=gris,densely dashed,line width=0.7pt},
  tz/.style={font=\tCuatro\bfseries,text=gris,anchor=north west},
  dos/.style={{Stealth[length=8pt,width=7pt]}-{Stealth[length=8pt,width=7pt]},draw=gris,line width=1.1pt}]
\draw[zona] (0,0) rectangle (2.0,3.3);
\draw[zona] (2.35,0) rectangle (3.95,3.3);
\draw[zona] (4.3,0) rectangle (8.9,3.3);
\draw[zona] (9.25,0) rectangle (12.1,3.3);
\node[tz] at (0.08,3.25) {NO CONFIABLE};
\node[tz] at (2.43,3.25) {BORDE};
\node[tz] at (4.38,3.25) {APLICACIÓN};
\node[tz] at (9.33,3.25) {DATOS Y SECRETOS};
\node[comp,minimum width=1.6in] (nav) at (1.0,1.25) {Navegador};
\node[comp,minimum width=1.2in] (px) at (3.15,1.25) {Proxy};
\node[comp,minimum width=4.0in] (app) at (6.6,1.25) {\textbf{Aplicación}\\{\tCuatro identifica, valida y hace cumplir}};
\node[pdp,minimum width=4.0in] (opa) at (6.6,2.45) {\textbf{Motor de decisión}\\{\tCuatro la matriz, como política}};
\node[comp,minimum width=2.4in] (pg) at (10.68,2.1) {\textbf{Base de datos}\\{\tCuatro segundo candado}};
\node[comp,minimum width=2.4in] (bao) at (10.68,0.75) {\textbf{Bóveda}\\{\tCuatro secretos y sello}};
\draw[dos] (nav) -- (px);
\draw[dos] (px) -- (app);
\draw[dos,draw=acento] (app) -- (opa);
\draw[dos] ($(app.east)+(0,0.18)$) -- ++(0.3,0) |- (pg.west);
\draw[dos] ($(app.east)+(0,-0.18)$) -- ++(0.3,0) |- (bao.west);
\end{tikzpicture}

\vspace{0.2in}
\dato{Esconder un botón no es autorización.} Si el servidor no decide, cualquiera llama directo.
\note{Flujo de una petición: el proxy termina la conexión cifrada; la aplicación identifica la sesión,
valida la forma de los datos y lee de la base los hechos; entrega al motor el sujeto, la acción, el
objeto y esos hechos; el motor devuelve la decisión con la regla que la determinó. Solo entonces se toca
el dato, en una transacción que también escribe la bitácora. La base aplica sus propias restricciones.

Fronteras de confianza: nada de lo que llega del navegador se da por bueno; los hechos se toman de la
base, nunca de la petición; si el motor no responde, se deniega; la base no confía en la aplicación.

Si preguntan qué pasa si el control está solo en la interfaz: el técnico envía a mano la petición de
aprobar su propia medición y el servidor la ejecuta. Es el error de referencia directa a objetos de la
clase 13.

Hospedaje, que el cliente dejó a nuestro criterio (minuta, P-23): un servidor virtual contratado, operado
por personal que la empresa designe. Hay además una ruta pública de solo lectura para comprobar
certificados.}
\end{frame}

%--------------------------- 20. TECNOLOGIAS -------------------------
\begin{frame}{Seis tecnologías libres,\\ elegidas por lo que protegen}
\quien{Axel}
\vspace{-0.05in}
{\renewcommand{\arraystretch}{1.36}
\begin{tabular}{@{}L{2.9in} L{3.7in} L{5.1in}@{}}
\toprule
\cabt{DECISIÓN} & \cabt{ELECCIÓN} & \cabt{LO QUE APORTA} \\
\midrule
Lenguaje y framework & Python 3.14, Django 5.2 LTS & Consultas parametrizadas; escape por omisión \\
Base de datos & PostgreSQL 18 & Seguridad por fila; emisión única \\
Sesión & De servidor, cookie opaca & \dato{Se revoca en el acto} \\
Autorización & Open Policy Agent 1.21 & Un solo lugar; la matriz como política \\
Contenedores & Podman sobre AlmaLinux 10 & Sin privilegios; confinados por SELinux \\
\bottomrule
\end{tabular}}
\note{Para cada una el documento indica versión, alternativas descartadas, aporte, limitaciones y
herramientas de análisis.

Descartes por licencia: HashiCorp Vault (BUSL desde 2023) y Docker Desktop. Se usa OpenBao, que es
software libre.

Descartes por seguridad: tokens JWT, porque no se pueden revocar en el acto al dar de baja o cambiar
de rol; MariaDB, porque no tiene seguridad por fila; Docker Engine, por su servicio con privilegios; una
cadena de bloques, porque resuelve la desconfianza entre organizaciones, que aquí no existe.

Limitaciones que hay que saber decir: los permisos de Django son por tipo de objeto, no por objeto; su
sitio de administración no se instala; el superusuario de PostgreSQL se salta todo; el motor no conoce el
historial y no filtra listados.

Contraseñas: Argon2id con sal por usuario. Segundo factor: llave de acceso o códigos temporales; nunca SMS.}
\end{frame}

%--------------------------- 21. BITACORA ----------------------------
\begin{frame}{La bitácora no la reescribe\\ ni quien administra el servidor}
\quien{Axel}
\vspace{0.2in}
\begin{tikzpicture}[x=1in,y=1in,
  reg/.style={draw=regla,line width=0.8pt,fill=plano,minimum width=2.25in,minimum height=1.0in,
              align=center,inner sep=3pt},
  ext/.style={draw=acento,line width=1.6pt,fill=acentoclaro,minimum width=2.6in,minimum height=1.0in,
              align=center,inner sep=3pt},
  fl/.style={-{Stealth[length=9pt,width=8pt]},draw=gris,line width=1.2pt}]
\draw[draw=gris,densely dashed,line width=0.7pt] (-0.1,-0.8) rectangle (8.55,1.15);
\node[font=\tCuatro\bfseries,text=gris,anchor=north west] at (0,1.1) {SERVIDOR: SOLO SE PUEDE INSERTAR};
\node[reg] (a) at (1.25,-0.05) {Registro 1041\\{\tCuatro huella del 1040}};
\node[reg] (b) at (4.2,-0.05) {Registro 1042\\{\tCuatro huella del 1041}};
\node[reg] (c) at (7.15,-0.05) {Registro 1043\\{\tCuatro huella del 1042}};
\draw[fl] (b) -- (a); \draw[fl] (c) -- (b);
\node[ext,minimum width=2.25in] (t) at (10.97,-0.05) {\textbf{Testigo externo}\\{\tCuatro sello de tiempo y auditor}};
\draw[fl,draw=acento] (c) -- node[above,pos=0.66,font=\tCuatro,text=gris]{cada día} (t);
\end{tikzpicture}

\vspace{0.3in}
Cambiar un registro rompe la cadena. Reescribirla entera no coincide con lo que ya salió.
\note{Cuatro mecanismos que se apoyan: (1) la aplicación solo puede insertar, no modificar ni borrar;
(2) cada registro incluye la huella SHA-256 del anterior; (3) una vez al día la última huella sale del
servidor, recibe un sello de tiempo RFC 3161 y va al buzón de la Gerencia y del auditor; (4) la acción y
su registro se confirman en la misma transacción.

Qué se registra: autenticación, cambios, consultas de datos sensibles, administración y toda decisión
denegada. Con qué campos: instante, cuenta y rol, acción, objeto y versión, resultado y regla, origen,
huellas.

Conservación: diez años para cambios y decisiones, que es la trazabilidad que el cliente debe poder
reconstruir (minuta, P-18); dos para consultas (supuesto P-19).

Secretos: viven en una bóveda; la llave de sello no sale, se le pide a la bóveda que firme. La bóveda se
abre con dos de tres custodios.}
\end{frame}

%--------------------------- 22. RESIDUALES --------------------------
\begin{frame}{Lo que no se cubre está declarado}
\quien{Javier}
\vspace{0.0in}
\begin{tabular}{@{}L{5.0in} L{6.7in}@{}}
\toprule
\cabt{RIESGO RESIDUAL} & \cabt{POR QUÉ SE ACEPTA} \\
\midrule
Colusión entre técnico y verificador & Evitarla exige una tercera persona por certificado \\
Una medición falsa desde el origen & El sistema no ve el árbol \\
Extracción lenta, bajo los umbrales & Bajarlos castiga a quien trabaja bien \\
Un vínculo con un dueño que nadie declaró & No existe un registro que lo descarte \\
El histórico de las hojas de cálculo & No se le puede dar la trazabilidad que no tuvo \\
\bottomrule
\end{tabular}
\note{Son trece riesgos residuales en el documento de arquitectura. Otros: equipo de usuario infectado; el
operador puede detener o destruir (se detecta, no se impide); un solo servidor; sin hardware de
seguridad dedicado; la fecha de medición es declarada, porque no hay aplicación sin conexión; el
proveedor de hospedaje tiene acceso físico; la aprobación no lleva firma digital certificada; y la
misma regla vive en dos lugares, el motor y la base.

Si preguntan qué requerimiento quedó sin cubrir: la trazabilidad de diez años hacia atrás. Las hojas de
cálculo se congelan, se les calcula la huella y lo migrado queda marcado como histórico.

Antes de entrar en operación hay diez criterios de aceptación, entre ellos una prueba de penetración
autorizada por escrito.}
\end{frame}

%--------------------------- 23. RESPUESTAS DEL CLIENTE -------------
\begin{frame}{El cliente respondió,\\ y el diseño cambió en cinco puntos}
\quien{Javier}
\vspace{0.0in}
\begin{tabular}{@{}L{0.9in} L{5.0in} L{5.6in}@{}}
\toprule
 & \cabt{LO QUE RESPONDIÓ} & \cabt{QUÉ CAMBIÓ} \\
\midrule
P-01 & Un verificador puede medir, por excepción & Habilitación temporal; no aprueba lo que midió \\
P-05 & Las fincas las asignan los supervisores & \dato{Rol nuevo: supervisor de campo} \\
P-06 & La revisión sigue después del cierre & El cierre detiene registrar, no aprobar \\
P-07 & Revisión técnica antes de autorizar & Solicitar, corregir, revisar, autorizar \\
P-10 & Avisar primero; no bloquear de inmediato & Aviso, verificación, restricción temporal \\
\bottomrule
\end{tabular}

\vspace{0.2in}
\pie{Catorce de veintiséis consultas, respondidas por escrito. Las otras doce son supuestos declarados. Minuta: informe de políticas, Anexo A.}
\note{Qué decir: el diseño se hizo completo con supuestos declarados, y la consulta sirvió como validación.
De catorce respuestas: dos confirmaron el supuesto (P-03, el verificador devuelve y no edita; P-11, el
técnico solo ve sus fincas), cuatro lo ajustaron, seis lo cambiaron y dos quedaron a nuestro criterio.

Lo que resistió: la técnica de separación. La respuesta a P-01 no la cambió; solo volvió dinámica la
exclusión entre técnico y verificador.

Lo que hubo que rehacer: el orden de la enmienda (P-07), lo que el cierre detiene (P-06) y quién asigna
las fincas (P-05). Y el diseño era más estricto de lo que el cliente quería: bloqueaba de inmediato
donde bastaba avisar (P-10).

Otros cambios para citar si preguntan: sociedades y representantes como dueños, y declaración de
vínculos con abstención (P-15); ninguna norma de auditoría se da por supuesta, y los auditores exportan
evidencia sellada (P-18); registro después de la visita, con fecha de medición y fecha de registro
(P-22); diseño sin cantidades fijas (P-02).

Cada cambio llegó hasta la matriz: diecisiete casos de prueba nuevos, T-21 a T-37, y cinco modificados.}
\end{frame}

%--------------------------- 24. DECISIONES DELEGADAS ----------------
\begin{frame}{Cuatro decisiones que el cliente\\ dejó en nuestras manos}
\quien{Javier}
\vspace{0.0in}
\begin{tabular}{@{}L{4.7in} L{7.0in}@{}}
\toprule
\cabt{LA PREGUNTA} & \cabt{LO QUE DECIDIMOS} \\
\midrule
¿Revisa un verificador distinto del original? & \dato{Sí, cuando la enmienda es material} \\
¿Hace falta firma digital certificada? & No: sello de la empresa y comprobación en línea \\
¿Llaves físicas para todo el personal? & No: llave de acceso del propio equipo \\
¿Dónde se hospeda y quién lo opera? & Servidor virtual contratado; personal designado \\
\bottomrule
\end{tabular}

\vspace{0.2in}
\pie{El cliente pidió que cada decisión tenga una justificación clara y responda a un riesgo real, no el sistema más complejo.}
\note{Verificador distinto (minuta, P-07): tres opciones. Siempre distinto: toda corrección, hasta una de
forma, espera a otra persona. El mismo: quien aprobó el dato erróneo dictamina sobre su propio error.
Elegimos según lo que cambia: es material si cambia un dato del cálculo de toneladas o si ya hay un
certificado; entonces revisa otro. Lo determina el sistema, no una persona.

Firma digital (P-04): la firma certificada de cada verificador da autenticidad y autoría con presunción
legal, pero no dice si el certificado fue revocado después, y cuesta por persona. Elegimos sello de la
empresa, comprobación en línea por código y aprobación personal con segundo factor, que queda en la
bitácora.

Llaves físicas (P-23): no hay presupuesto para todos. La llave de acceso del propio dispositivo
(passkey) resiste igual el correo falso y no cuesta nada. Se recomiendan llaves físicas solo para
Gerencia, administrador, operador y custodios de la bóveda.

Hospedaje (P-23): servidor propio, servidor virtual contratado o servicios administrados, con los
cuatro criterios del cliente. El propio exige inversión y cuidar energía y enlace en una oficina
rural; los administrados no son software libre. Lo cotidiano está automatizado.}
\end{frame}

%--------------------------- 25. REFERENCIAS -------------------------
\begin{frame}{Referencias principales}
\quien{}
\vspace{0.0in}
{\tCuatro
\begin{columns}[T,onlytextwidth]
\begin{column}{0.485\textwidth}
\raggedright
Abadi, M., \& Fournet, C. (2003). Access control based on execution history. \emph{NDSS}.\\[5pt]
Clark, D. D., \& Wilson, D. R. (1987). A comparison of commercial and military computer security
policies. \emph{IEEE S\&P}.\\[5pt]
Kuhn, D. R., Coyne, E. J., \& Weil, T. R. (2010). Adding attributes to role-based access control.
\emph{Computer, 43}(6).\\[5pt]
Li, N., Tripunitara, M. V., \& Bizri, Z. (2007). On mutually exclusive roles and separation-of-duty.
\emph{ACM TISSEC, 10}(2).\\[5pt]
Saltzer, J. H., \& Schroeder, M. D. (1975). The protection of information in computer systems.
\emph{Proc. IEEE, 63}(9).\\[5pt]
Sandhu, R. S., Coyne, E. J., Feinstein, H. L., \& Youman, C. E. (1996). Role-based access control
models. \emph{Computer, 29}(2).
\end{column}
\begin{column}{0.485\textwidth}
\raggedright
Schneier, B., \& Kelsey, J. (1999). Secure audit logs to support computer forensics. \emph{ACM
TISSEC, 2}(2).\\[5pt]
Simon, R. T., \& Zurko, M. E. (1997). Separation of duty in role-based environments. \emph{CSFW}.\\[5pt]
Temoshok, D., et al. (2025). \emph{Digital identity guidelines: Authentication and authenticator
management} (NIST SP 800-63B-4).\\[5pt]
Ley N.º 8968 (2011) y su reglamento, Decreto N.º 37554-JP. Ley N.º 7975 (2000). Ley N.º 8454
(2005). Ley N.º 9048 (2012).\\[5pt]
ISO/IEC 17029:2019. ISO 14065:2020. ISO/IEC 27001:2022.\\[5pt]
Vargas Rivera, A. (2026). \emph{Material de estudio del curso CE-1115 y enunciado del proyecto,
parte 1}. ITCR.
\end{column}
\end{columns}
}
\note{La lista completa, en formato APA, está en cada uno de los cinco documentos.}
\end{frame}

\end{document}
